% ======================================================================
\chapter{Непараметрические (свободные от распределения) контрольные карты по количественным признакам для одной переменной}
% ======================================================================

\section*{Обзор главы}
В главе 3 мы рассматривали параметрические контрольные карты по количественным признакам для одной переменной. В этой главе мы рассматриваем непараметрические контрольные карты. Как и в статистической литературе, эти карты также называют картами, свободными от распределения, и мы используем термины «свободные от распределения» (distribution-free) и «непараметрические» как взаимозаменяемые. Подавляющее большинство доступных непараметрических карт предназначены для параметра положения, который представляет собой центр непрерывного распределения. Параметром положения, например, может быть среднее или медиана распределения процесса.

\section{Введение}
В рамках статистического контроля процессов (SPC) часто предполагается, что картина случайных причин, или сам процесс, следует некоторому известному параметрическому распределению. Контрольные карты, спроектированные и применяемые в предположении параметрического распределения, называются параметрическими картами. Например, наиболее распространённым является предположение о нормальном распределении. Статистические свойства этих контрольных карт являются точными только в том случае, если предположение о нормальности действительно выполняется. Однако во многих приложениях истинное лежащее в основе распределение процесса неизвестно или не является нормальным, и, следовательно, свойства или характеристики этих стандартных (параметрических) контрольных карт в таких ситуациях могут сильно пострадать. Поэтому необходима разработка и применение контрольных карт, не зависящих от нормальности или других конкретных параметрических предположений о распределении. Этой цели могут служить карты, свободные от распределения. И всё же, несмотря на доводы в пользу использования непараметрических, или свободных от распределения, контрольных карт, их внедрение и разработка шли несколько медленно. Причин этому несколько.

Во-первых, можно возразить, что центральная предельная теорема обеспечит работу контрольной карты так, как ожидается. Хотя это может быть верно для контрольных карт, основанных на средних некоторых выборочных статистик, в общем случае это не так. Там, где контрольные карты нужно применять, например, к отдельным наблюдениям, центральная предельная теорема неприменима и на неё нельзя ссылаться. Во-вторых, свою роль сыграли как прошлая недоступность достаточных вычислительных мощностей, так и представление о том, что при использовании этих (простых) методик, часто основанных на подсчёте и/или ранжировании, необходимо жертвовать «эффективностью». Однако первое уже не является проблемой в нынешний компьютерный век, а второе не обязательно верно. Например, в статистической литературе хорошо известно, что для некоторых распределений с тяжёлыми хвостами и/или асимметричных некоторые непараметрические процедуры превосходят свои параметрические аналоги. Примечательно, что даже когда лежащее в основе распределение нормально, эффективность некоторых непараметрических критериев по отношению к соответствующим (оптимальным) методам нормальной теории может достигать 0,955 (см., например, Gibbons and Chakraborti, 2010, с. 218). Для некоторых других распределений с тяжёлыми хвостами и скошенных распределений эффективность может быть даже выше (см. Conover, 1999); см. также Hodges and Lehmann (1956).

Следовательно, контрольные карты, свободные от распределения, могут быть полезны, а иногда даже предпочтительнее, когда лежащее в основе распределение процесса неизвестно или не может быть разумно смоделировано параметрическим распределением. Непараметрические карты устойчивы (робастны) в смысле контроля частоты ложных тревог, или вероятности ложной тревоги, в состоянии IC и, следовательно, перспективы их обнаружения сдвига являются надёжными. В такой постановке задача непараметрической контрольной карты для параметра положения сводится к задаче проверки гипотез о медиане (или другом квантиле) распределения. Наиболее известны два непараметрических критерия: знаковый критерий и критерий знаковых рангов Уилкоксона. Мы подробно описываем карты, основанные на этих критериях.

Цель этой главы --- рассмотреть как хорошо известные, так и некоторые новые непараметрические (свободные от распределения) контрольные карты по количественным признакам для параметра положения процесса, которые применяются, когда лежащее в основе распределение процесса непрерывно, но не является нормальным, или неизвестно, или не может быть разумно смоделировано известным распределением. Как и в главе 3, мы рассматриваем три основных класса карт: карты Шухарта, карты CUSUM и карты EWMA. В каждом случае мы выводим контрольные карты, обсуждаем их свойства и иллюстрируем применение на числовых примерах. Как и ранее, рассматриваются случаи K (параметры известны) и U (параметры неизвестны), и аналогично главе 3, в разделе о случае U вначале предполагается, что есть эталонная выборка из фазы I, на основании которой строятся границы контрольной карты. Мы также вводим понятие длины серии и связанных с ней характеристик для сравнения эффективности различных карт и обсуждаем, как они могут быть вычислены либо точно с использованием цепей Маркова (в случае K), либо с помощью симуляционного моделирования. По ходу дела мы вводим и иллюстрируем условный и безусловный подходы к распределению длины серии в случае U. Затем мы обсуждаем некоторые вопросы, связанные с оцениванием параметров, объёмом данных фазы I и повышением эффективности контрольных карт с помощью правил серий и правил чувствительности.

\section{Свободные от распределения контрольные карты по количественным признакам в случае K}
В этом разделе мы описываем непараметрические карты Шухарта, CUSUM и EWMA в случае K, то есть когда значение параметра в состоянии IC (например, медиана) известно или задано.

\subsection{Контрольные карты Шухарта}
\subsubsection{Контрольные карты Шухарта на основе знаков}
Рассмотрим ситуацию, когда выборки или подгруппы из \(n\) (\(\ge 1\)) наблюдений берутся случайно через регулярные интервалы времени из неизвестного непрерывного распределения процесса с представляющим интерес параметром \(\mu\). В общем случае \(\mu\) может быть любым параметром, но обычно в качестве него берут медиану процесса. Предположим, что известное или заданное значение \(\mu\) в состоянии IC равно \(\mu_0\). Тогда, если распределение симметрично, целевое значение совпадает со средним процесса. Пусть \(X_{i1}, X_{i2}, \ldots, X_{in}\), \(i = 1, 2, \ldots\), --- выборка или подгруппа в \(i\)-й момент времени, и предположим, что каждое наблюдение \(X_{ij}\), \(j = 1, 2, \ldots, n\), сравнивается с \(\mu_0\) и для каждой выборки фиксируется число наблюдений выше и ниже \(\mu_0\). Пусть \(n_i^+\) (\(n_i^-\)) обозначает число наблюдений \(x_{ij}\), больших (меньших) \(\mu_0\), в \(i\)-й выборке. Тогда пусть
\begin{equation}
	SN_i = \sum_{j=1}^{n} \phi(X_{ij} - \mu_0)
	\label{eq:4.1}
\end{equation}
где
\begin{equation}
	\phi(t) =
	\begin{cases}
		\phantom{-}1 & \text{если } t > 0, \\
		\phantom{-}0 & \text{если } t = 0, \\
		-1 & \text{если } t < 0.
	\end{cases}
	\label{eq:4.2}
\end{equation}

Статистику \(SN_i\) можно использовать как статистику для нанесения на карту для знаковой карты. Можно видеть, что \(SN_i = n_i^+ - n_i^-\), то есть \(SN_i\) --- разность между числом наблюдений, больших \(\mu_0\), и числом наблюдений, меньших \(\mu_0\), в \(i\)-й выборке. Для случая отдельных данных, то есть \(n = 1\), статистика \(SN_i\) равна \(+1\), когда наблюдение больше \(\mu_0\), и \(-1\), когда наблюдение меньше \(\mu_0\). Из \(SN_i = n_i^+ - n_i^-\) следует, что \(SN_i + n = 2n_i^+\) в предположении отсутствия связок, так что \(n_i^+ = \frac{SN_i + n}{2}\). Отметим, что можно записать \(n_i^+ \le n\), число наблюдений, больших заданного \(\mu_0\), как \(n_i^+ = T_i = \sum_{j=1}^{n} \Psi(X_{ij} - \mu_0)\), где \(\Psi(A)\) обозначает индикаторную функцию
\[
\Psi(A) =
\begin{cases}
	1 & \text{если } A > 0, \\
	0 & \text{если } A \le 0.
\end{cases}
\]
Из соотношения \(n_i^+ = \frac{SN_i + n}{2}\) получаем \(SN_i = 2T_i - n\), поскольку \(n_i^+ = T_i\), так что знаковую карту можно эквивалентно выразить через статистику \(T_i\).

Одно из преимуществ использования знакового критерия, основанного на \(SN_i\), состоит в том, что его математическое ожидание в состоянии IC равно 0 (это будет показано далее). Однако в (непараметрической) литературе более известна версия статистики \(T_i\) (поскольку распределение \(T_i\) биномиальное), на которой основан знаковый критерий, и она называется статистикой знакового критерия (см., например, Gibbons and Chakraborti, 2010, с. 168). Для знаковой карты Шухарта (обозначаемой Shewhart-SN) статистика \(SN_i\) используется как статистика для нанесения на карту.

Поскольку распределение процесса предполагается непрерывным, \(P(X_{ij} = \mu_0) = 0\), и сценарий, при котором \(\phi(X_{ij} - \mu_0) = 0\), то есть «нули», или нулевые разности с \(\mu_0\), по крайней мере теоретически не должен встречаться. Однако там, где наблюдения округлены, изредка может наблюдаться «ноль»; но, пока число «нулей» не слишком велико, \(SN_i\) можно безопасно вычислять по определению из уравнения \eqref{eq:4.1}.

Контрольные границы и центральная линия двусторонней карты Shewhart-SN имеют вид
\begin{equation}
	\begin{aligned}
		UCL &= a \\
		CL &= 0 \\
		LCL &= -a
	\end{aligned}
	\label{eq:4.3}
\end{equation}
где \(a\) --- положительное целое число такое, что \(0 < a \le n\). Следовательно, карта Shewhart-SN подаёт сигнал о том, что произошёл сдвиг, если \(|SN_i| \ge a\), то есть если \(SN_i \ge a\) или если \(SN_i \le -a\). Напротив, пока \(-a < SN_i < a\), процесс считается находящимся в состоянии IC относительно целевой медианы \(\mu_0\). Односторонние и двусторонние формы карт Шухарта введены в разделе 3.2.1; здесь нижняя односторонняя карта Shewhart-SN подаёт сигнал о том, что произошёл нисходящий сдвиг медианы, если \(SN_i \le -a\) (карта имеет только нижнюю контрольную границу, или LCL), тогда как верхняя односторонняя карта Shewhart-SN подаёт сигнал о том, что произошёл восходящий сдвиг, если \(SN_i \ge a\) (карта имеет только верхнюю контрольную границу, или UCL). Здесь мы рассматриваем двусторонние карты, поскольку они более общие (так как практик может не знать направления сдвига медианы).

\paragraph{Вычисление константы для нанесения на карту \(a\).} Константа для нанесения на карту \(a\) --- это наименьшее целое число в интервале \(0 < a \le n\) такое, что \(P_{\mu_0 | IC}(SN_i \ge a) \le \frac{FAR}{2}\). Подстрочный индекс \(\mu_0 | IC\) указывает, что константа для нанесения на карту находится в случае K и в состоянии IC. Здесь важно отметить, что, поскольку распределение \(SN_i\) дискретно, есть вероятность, что равенство выше не будет достигнуто точно. Поэтому, используя консервативный подход, \(a\) находится так, чтобы достигнутая FAR была меньше или равна заданной или номинальной FAR. Например, полагая номинальную FAR, обозначаемую \(FAR_0\), равной 0.0027, получаем \(\frac{FAR_0}{2} = \frac{0.0027}{2} = 0.00135\), так что \(a\) --- наименьшее целое число такое, что \(P_{\mu_0 | IC}(SN_i \ge a) \le 0.00135\).

\noindent\textbf{Замечание 4.1. Использование биномиальных вероятностей для вычисления константы для нанесения на карту \(a\) карты Shewhart-SN.}

Мы обращаем внимание на запись «правый хвост» в биномиальной таблице (см. таблицу A в Приложении A), поскольку мы ищем UCL, то есть сосредоточены на вероятности верхнего хвоста \(P_{\mu_0 | IC}(T_i \ge t)\), чтобы вычислить \(P_{\mu_0 | IC}(SN_i \ge a)\).

Далее рассмотрим два примера. Первый, для \(n = 5\), иллюстрирует случай, когда нельзя найти достигнутое значение FAR, которое меньше 0.0027. Обычно так бывает при малых значениях \(n\) и номинальной FAR. Второй пример, для \(n = 20\), иллюстрирует случай, когда достигнутое значение FAR меньше 0.0027 получить можно. Таким образом, объём выборки (подгруппы), которая контролируется, важен, когда речь идёт о непараметрических контрольных картах и о выборе номинальных значений FAR. Отметим, что так обычно обстоит дело с картами, основанными на статистиках для нанесения на карту с дискретным распределением, что имеет место и для карт по альтернативным признакам (в этой книге не рассматриваемых).

\emph{Для \(n = 5\).} По записи «правый хвост» для 5 имеем \(\frac{FAR}{2} = 0.03125\), что является наименьшей вероятностью, связанной с \(n = 5\). Хотя эта вероятность не меньше и не равна 0.00135, это наименьшая вероятность, которую мы можем получить при выборке объёма 5. Теперь, чтобы получить UCL, то есть \(a\), используем линейное соотношение \(SN_i = 2T_i - n\), так что \(a = 2(5) - 5 = 5\) при \(FAR = (2)(0.03125) = 0.0625\). Таким образом, для \(n = 5\) имеем \(LCL / UCL = \pm 5\) и \(CL = 0\) при соответствующей \(FAR = 0.0625\), которая является наименьшей достижимой FAR. Поскольку значение медианы в состоянии IC известно или задано (случай K), \(ARL_{IC} = \frac{1}{FAR} = \frac{1}{0.0625} = 16\). Опять же, это плохо, поскольку далеко от типичного желаемого значения 370. Однако это наибольшая достижимая \(ARL_{IC}\) при работе с выборкой объёма 5.

Фактически можно видеть, что, если объём выборки не равен 10 или больше, знаковая карта может быть несколько непривлекательной с точки зрения малой FAR (или большой \(ARL_{IC}\)), обычно используемой в SPC. Читателю предлагается вычислить значения FAR для \(n = 6, 7, 8, 9\) и \(10\), чтобы в этом убедиться.

\emph{Для \(n = 20\).} По записи «правый хвост» для 17 имеем \(\frac{FAR}{2} = 0.001288\), что меньше 0.00135. Теперь, чтобы получить UCL, то есть \(a\), используем линейное соотношение \(SN_i = 2T_i - n\), так что \(a = 2(17) - 20 = 14\) при \(FAR = (2)(0.001288) = 0.002576\). Таким образом, для \(n = 20\) имеем \(LCL / UCL = \pm 14\) и \(CL = 0\) при соответствующей \(FAR = 0.002576\) и \(ARL_{IC} = \frac{1}{FAR} = \frac{1}{0.002576} \approx 388\), что хорошо, поскольку близко к типичному желаемому значению 370.

\subsubsection{Знаковая карта Шухарта с правилами сигналов на основе серий}
Хотя карта Shewhart-SN мощна при обнаружении больших сдвигов параметра положения процесса \(\mu\), как и любая карта Шухарта, она, как ожидается, будет недостаточно чувствительной при обнаружении меньших сдвигов. Следовательно, как и в случае параметрической карты Шухарта, мы можем повысить эффективность карты Shewhart-SN при обнаружении малых сдвигов, добавив правила сигналов на основе серий. Стандартное правило сигналов на основе серий «\(k\) из \(k\)» подаёт сигнал, когда либо \(k\) последовательных точек нанесены на UCL или выше неё, либо \(k\) последовательных точек нанесены на LCL или ниже неё. Иллюстрация приведена в разделе 4.3. Читателю также рекомендуется обратиться к Human, Chakraborti and Smit (2010b) за дополнительными примерами и подробностями о картах Shewhart-SN, основанных на стандартном правиле «2 из 2».

\subsubsection{Методы вычисления распределения длины серии}
\paragraph{Точный подход.} Для карт Шухарта и знаковых контрольных карт типа Шухарта можно точно вычислить распределение длины серии и различные характеристики длины серии. Здесь следует отметить, что исходные карты Шухарта предназначены для параметрического случая (нормального распределения), и некоторые авторы использовали термин «карты типа Шухарта» в непараметрическом случае, чтобы подчеркнуть, что они аналогичны. Однако для простоты мы используем термин «Шухарт» на протяжении всей главы как для параметрических, так и для непараметрических карт. Как и в случае параметрической карты Шухарта в случае K, распределение длины серии карты Shewhart-SN есть \(GEO(\theta)\), то есть геометрическое с вероятностью \(\theta = P(SN_i \ge a \text{ или } SN_i \le -a)\), которая является вероятностью сигнала. Поскольку \(SN_i = 2T_i - n\), вероятность \(\theta\) удобно выразить через распределение \(T_i\), которое, как известно, в общем случае является биномиальным с параметрами \(n\) и \(p = P(X_{ij} \ge \mu_0)\). В таблице \ref{tab:4.1} показаны распределительные свойства статистик \(T_i\) и \(SN_i\) соответственно.

\begin{table}[!ht]
	\centering
	\caption{Распределительные свойства статистик \(T_i\) и \(SN_i\) соответственно.}
	\label{tab:4.1}
	\begin{tabular}{@{}lll@{}}
		\toprule
		 & \(T_i\) & \(SN_i\) \\
		\midrule
		Математическое ожидание & \(E(T_i) = np\) & \(E(SN_i) = E(2T_i - n)\) \\
		 & & \(= 2E(T_i) - n\) \\
		 & & \(= 2np - n\) \\
		 & & \(= n(2p - 1)\) \\
		\midrule
		Стандартное отклонение & \(SD(T_i) = \sqrt{np(1 - p)}\) & \(SD(SN_i) = 2\sqrt{np(1 - p)}\) \\
		\bottomrule
	\end{tabular}
\end{table}

\begin{figure}[!ht]
	\centering
	\includegraphics[width=0.6\linewidth]{"./images/fig-4.1.png"}
	\caption[]{Распределение вероятностей в состоянии IC статистик \(T_i\) и \(SN_i\) для \(n = 10\).}
	\label{fig-4.1}
\end{figure}

На рисунке \ref{fig-4.1} показано распределение вероятностей в состоянии IC статистик \(T_i\) и \(SN_i\) для \(n = 10\). Заметим, что в состоянии IC при \(p = 0.5\) распределение вероятностей \(T_i\) симметрично относительно своего математического ожидания \(E(T_i) = np = (10)(0.5) = 5\). Напротив, распределение вероятностей \(SN_i\) в состоянии IC симметрично относительно своего математического ожидания \(E(SN_i) = n(2p - 1) = 10(2(0.5) - 1) = 0\).

Распределение длины серии в состоянии IC двусторонней карты Shewhart-SN, например, следует геометрическому распределению \(GEO(\theta)\) с вероятностью успеха (FAR) \(\theta\), где
\[
\theta = (0.5)^n \left[ \binom{n}{a} + \binom{n}{a+1} + \cdots + \binom{n}{\left[\frac{n-a}{2}\right] + n} \right],
\]
и \([d]\) обозначает наибольшее целое, не превосходящее \(d\). Это потому, что в состоянии IC \(T_i\) следует биномиальному распределению с параметрами \(n\) и \(p = P(X_{ij} \ge \mu_0) = 0.5\).

Как и в случае параметрических карт в случае K, \(ARL_{IC}\) карты Shewhart-SN есть величина, обратная FAR, а другие моменты распределения длины серии в состоянии IC можно найти, используя свойства геометрического распределения. Распределение длины серии карты Shewhart-SN в состоянии OOC также геометрическое с вероятностью успеха \(\theta\), которую можно вычислить аналогичным образом с помощью биномиального распределения.

Далее рассмотрим карту Shewhart-SN с правилами сигналов на основе серий. Human et al. (2010b) рассматривали использование правил серий «1 из 1», «\(k\) из \(k\)» и «\(k\) из \(w\)»; читателю рекомендуется обратиться к их работе за подробностями. Напомним, что единственные предположения, которые должны выполняться для карт SN, --- это независимость наблюдений и непрерывность их распределения. Отметим, что статистика для нанесения на карту, используемая Human et al. (2010b), --- не статистика \(SN_i\), введённая ранее (см. уравнение \eqref{eq:4.1}), а классическая знаковая статистика \(T_i\) (см., например, Gibbons and Chakraborti, 2010, с. 168). Контрольные границы и CL задаются как \(UCL = n - b\), \(LCL = a\) и \(CL = n\theta_0\) соответственно, где \(a\) и \(b\) --- значения между 0 и \(n\), которые обычно выбираются для получения нужной \(ARL_{IC}\). Поскольку это первая карта с правилами серий, рассматриваемая в этой главе, мы начинаем с простейшего случая, то есть правила серий «1 из 1», чтобы ввести обозначения. Для верхней односторонней карты SN «1 из 1» индикаторные переменные «1--0» \(Z_h^+\) определяются в связи со статистикой для нанесения на карту \(T_h\), \(h = 1, 2, 3, \ldots\), каждая из которых означает либо успех (обозначается 1), указывающий на сигнал, либо неудачу (обозначается 0), указывающую на отсутствие сигнала, с соответствующими вероятностями \(p^+\) и \(1 - p^+\). Эта карта «1 из 1» SN подаёт сигнал, когда \(T_i \ge UCL\). Аналогично для нижней односторонней карты «1 из 1» индикаторная переменная \(Z_h^-\) обозначает успех (1) или неудачу (0) с вероятностями \(p^-\) и \(1 - p^-\) соответственно; эта карта подаёт сигнал, когда \(T_i \le LCL\). Для двусторонней карты «1 из 1» индикаторные переменные «1--0» \(Z_h\), \(h = 1, 2, 3, \ldots\), снова обозначают успех (1) (то есть сигнал) или неудачу (0) (то есть отсутствие сигнала) с вероятностями \(p\) и \(1 - p\); эта карта подаёт сигнал, когда \(T_i \le LCL\) или \(T_i \ge UCL\).

Теперь, когда основы определены, рассмотрим верхнюю одностороннюю карту Shewhart-SN, дополненную правилом серий «2 из 2». Односторонняя карта объясняется первой, поскольку это первая карта с правилами серий, рассматриваемая в этой главе. Двусторонние карты будут обсуждаться в более позднем разделе. Пусть состояние \(S_1\) представляет область на UCL и выше неё, а состояние \(S_0\) --- область ниже UCL. Пусть \(p\) обозначает вероятность сигнала, то есть вероятность того, что статистика для нанесения на карту окажется на UCL или выше неё. Следовательно, вероятность сигнала \(p = P(T_i \ge UCL) = P(T_i \ge n - b) = I_p(n - b, b + 1)\), где \(I_p(n - b, b + 1)\) обозначает неполный бета-интеграл, или функцию распределения бета-распределения, в точке \(p\) с параметрами \(n - b\) и \(b + 1\). Соотношение между вероятностью верхнего хвоста биномиального распределения и неполным бета-интегралом хорошо известно; см., например, Gibbons and Chakraborti (2010, с. 66). Пусть также \(q = 1 - p\) обозначает вероятность отсутствия сигнала, а \(\phi\) --- фиктивное состояние. Верхняя односторонняя карта Shewhart-SN, дополненная правилами серий «2 из 2», подаёт сигнал, когда две последовательные статистики для нанесения на карту SN оказываются на UCL или выше неё. Читателю предлагается показать, что существенная матрица переходных вероятностей (TPM) в этом случае задаётся элементами таблицы \ref{tab:4.2}; заинтересованному читателю рекомендуется обратиться к работе Human et al. (2010b) за подробностями и дополнительными примерами.

\begin{table}[!ht]
	\centering
	\caption{Существенная TPM для верхней односторонней контрольной карты, дополненной правилом сигналов на основе серий «2 из 2».}
	\label{tab:4.2}
	\begin{tabular}{@{}lccc@{}}
		\toprule
		Состояния & \(\phi\) & \(S_0\) & \(S_1\) \\
		\midrule
		\(\phi\) & 0 & \(p_0\) & \(p_1\) \\
		\(S_0\) & 0 & \(p_0\) & \(p_1\) \\
		\(S_1\) & 0 & \(p_0\) & 0 \\
		\bottomrule
	\end{tabular}
\end{table}

Как только существенная TPM получена, распределение длины серии и его различные характеристики, такие как ARL и т. д., можно вычислить с помощью теории марковских цепей (MC), как показано в уравнениях 3.22--3.24. Опять же читателю рекомендуется обратиться к Human et al. (2010b) за подробностями. Можно видеть, например, что \(ARL = (1 + p)/p^2\) и \(FAR = (p_0)^2\), где \(p_0\) --- вероятность сигнала, когда процесс находится в состоянии IC. Таким образом, теория MC играет полезную роль при выводе распределения длины серии для карт Шухарта, дополненных правилами серий, где статистики для нанесения на карту зависимы.

\paragraph{Симуляционный подход.} Хотя для карты Shewhart-SN распределение длины серии и его различные характеристики можно вычислить аналитически, в общем случае симуляция может быть полезной и предпочтительной. Для получения характеристик длины серии карты Shewhart-SN можно использовать следующий алгоритм симуляции. Соответствующая SAS-программа приведена в Приложении B (SAS-программа № 1), а вывод показан ниже на рисунке \ref{fig-4.2}.

\begin{enumerate}
	\item[\textbf{Шаг 1:}] После задания объёма подгруппы (\(n \ge 1\)) и величины сдвига (\(\delta\)), который нужно обнаружить, генерируем случайные подгруппы из некоторого непрерывного распределения, например стандартного нормального, Стьюдента \(t\), Лапласа, логистического, гамма и т. д. Для генерации данных можно использовать любое непрерывное распределение, симметричное или асимметричное.
	\item[\textbf{Шаг 2:}] Вычисляем статистики \(SN_t\) для каждой подгруппы.
	\item[\textbf{Шаг 3:}] Вычисляем контрольные границы, то есть константу для нанесения на карту \(a\), и сравниваем \(SN_t\) с контрольными границами. Вычисление константы для нанесения на карту \(a\) подробно объяснено в разделе 4.2.1.1.
	\item[\textbf{Шаг 4:}] Число подгрупп, необходимое до первого попадания \(SN_t\) на контрольные границы или за них, фиксируется как наблюдение из распределения длины серии.
	\item[\textbf{Шаг 5:}] Повторяем шаги 1--4 в общей сложности 100\,000 раз.
	\item[\textbf{Шаг 6:}] Как только получен «набор данных» из 100\,000 наблюдений из распределения длины серии, для получения характеристик длины серии можно использовать \texttt{proc univariate} SAS. Для \(\delta = 0\) это будут характеристики длины серии в состоянии IC, тогда как для \(\delta \ne 0\) --- характеристики длины серии в состоянии OOC.
\end{enumerate}

\begin{figure}[!ht]
	\centering
	\includegraphics[width=0.7\linewidth]{"./images/fig-4.2.png"}
	\caption[]{Вывод SAS распределения RL при параметрах проектирования (\(n = 10\), \(a = 10\)) карты Shewhart-SN, выбранных для номинальной \(ARL_{IC} = 500\).}
	\label{fig-4.2}
\end{figure}

\subsubsection{Контрольные карты Шухарта на основе знаковых рангов}
С одной стороны, хотя знаковый критерий является универсальным непараметрическим критерием положения, он один из простейших и, поскольку он, очевидно, не использует большую часть информации из данных, может оказаться не очень мощным. С другой стороны, знаковый критерий можно применять при минимальных предположениях о распределении и даже с дихотомизированными данными. Последнее может быть ситуацией, когда единственная доступная информация для каждой испытуемой единицы --- это то, было ли измерение выше (или ниже) целевого значения представляющего интерес параметра. Это может быть плюсом в некоторых практических ситуациях.

Популярной альтернативой знаковому критерию является критерий знаковых рангов. Этот критерий использует как знаки, так и ранги данных и является популярным непараметрическим критерием для медианы симметричной непрерывной совокупности. Хорошо известно, что критерий знаковых рангов эффективнее знакового критерия, но дополнительная эффективность достигается ценой дополнительного предположения о симметрии. Та же дискуссия применима к знаковой и знаково-ранговой контрольным картам, которые основаны на соответствующих критериях. Кроме того, отметим, что универсальный знаковый критерий (контрольную карту) можно использовать для всех квантилей, тогда как критерий знаковых рангов (контрольную карту) можно использовать только для медианы.

\paragraph{Карта Шухарта по знаковым рангам.} Начнём с контрольной карты, основанной на знаково-ранговой статистике (обозначается Shewhart-SR). Пусть \(X_{i1}, X_{i2}, \ldots, X_{in}\), \(i = 1, 2, \ldots\) и \(n > 1\), --- выборки или подгруппы независимых наблюдений, взятых случайно (и последовательно) из неизвестного, но непрерывного распределения процесса, симметричного относительно известного значения \(\mu_0\) в состоянии IC --- обычно называемого целевым значением. Поскольку распределение процесса предполагается симметричным, среднее и медиана совпадают, и таким образом постановка задачи для знаково-ранговой карты медианы сопоставима с постановкой задачи для параметрической карты \(\overline{X}\) Шухарта, которая предполагает нормальность (являющуюся, конечно, симметричным распределением). Пусть \(R_{ij}\) обозначает ранг абсолютного отклонения (разности) \(|x_{ij} - \mu_0|\) среди \((|x_{i1} - \mu_0|, |x_{i2} - \mu_0|, \ldots, |x_{in} - \mu_0|)\) для \(i = 1, 2, 3, \ldots\) и \(j = 1, 2, 3, \ldots, n\). Тогда пусть
\begin{equation}
	T_n^+ = \sum_{j=1}^{n} \Psi(x_{ij} - \mu_0) R_{ij}
	\label{eq:4.4}
\end{equation}
где \(\Psi(t)\) --- индикаторная функция, определённая ранее. Статистика \(T_n^+\) --- известная статистика знаковых рангов Уилкоксона (см., например, Gibbons and Chakraborti, 2010, с. 195). Очевидно, \(T_n^+\) --- сумма рангов абсолютных значений отклонений, соответствующих положительным разностям. Теперь определим
\begin{equation}
	SR_i = \sum_{j=1}^{n} \phi(x_{ij} - \mu_0) R_{ij},
	\label{eq:4.5}
\end{equation}
которая является разностью между суммой рангов, связанных с положительными разностями (то есть с теми наблюдениями, для которых \(X_{ij} - \mu_0 > 0\)), и суммой рангов, связанных с отрицательными разностями (то есть с теми наблюдениями, для которых \(X_{ij} - \mu_0 < 0\)). Статистика \(SR_i\) линейно связана с более классической статистикой знаковых рангов Уилкоксона \(T_n^+\) формулой
\[
SR_i = 2T_n^+ - \frac{n(n + 1)}{2},
\]
поскольку \(T_n^+ + T_n^- = n(n + 1)/2\).

Теперь продолжим вводить карту Shewhart-SR. Контрольные границы и центральная линия двусторонней карты Shewhart-SR имеют вид
\begin{equation}
	\begin{aligned}
		UCL &= a \\
		CL &= 0 \\
		LCL &= -a
	\end{aligned}
	\label{eq:4.6}
\end{equation}
где \(a\) --- некоторое положительное целое число в интервале \(0 < a \le \frac{n(n + 1)}{2}\). Следовательно, карта Shewhart-SR подаёт сигнал о том, что произошёл сдвиг медианы, если \(|SR_i| \ge a\), то есть если \(SR_i \ge a\) или если \(SR_i \le -a\). Напротив, пока \(-a < SR_i < a\), процесс считается находящимся в состоянии IC относительно целевого значения \(\mu_0\).

Односторонние и двусторонние формы карт Шухарта введены в разделе 3.2.1; здесь нижняя односторонняя карта Shewhart-SR подаёт сигнал о том, что произошёл нисходящий сдвиг, если \(SR_i \le -a\) (карта имеет только LCL), тогда как верхняя односторонняя карта Shewhart-SR подаёт сигнал о том, что произошёл восходящий сдвиг, если \(SR_i \ge a\) (карта имеет только UCL). Здесь мы рассматриваем двусторонние карты, поскольку они более общие (так как практик может не знать направления сдвига) и поскольку односторонние карты в литературе получили мало внимания.

\paragraph{Вычисление константы для нанесения на карту \(a\).} Константа для нанесения на карту \(a\) --- это наименьшее целое число в интервале \(0 < a \le \frac{n(n + 1)}{2}\) такое, что \(P_{\mu_0 | IC}(SR_i \ge a) \le \frac{FAR}{2}\). Как и в случае знаковой карты, опять же, поскольку распределение статистики для нанесения на карту \(SR_i\) дискретно, есть вероятность, что равенство выше не будет достигнуто точно. Поэтому, используя, как и ранее, консервативный подход, \(a\) находится так, чтобы достигнутая FAR была меньше или равна заданной FAR. Полагая идеальную FAR равной 0.0027, где \(\frac{FAR}{2} = \frac{0.0027}{2} = 0.00135\), получаем, что \(a\) --- наименьшее целое число такое, что \(P_{\mu_0 | IC}(SR_i \ge a) \le 0.00135\).

\noindent\textbf{Замечание 4.2. Использование вероятностей знаковых рангов Уилкоксона для вычисления константы для нанесения на карту \(a\) карты Shewhart-SR.}

Мы обращаем внимание на запись «правый хвост» в таблице B в Приложении A, поскольку мы сосредоточены на UCL, то есть на вероятности верхнего хвоста \(P_{\mu_0 | IC}(T_n^+ \ge t)\), чтобы в итоге получить \(P_{\mu_0 | IC}(SR_i \ge a)\).

Как и в случае знаковой карты, далее рассмотрим два примера. Первый, для \(n = 5\), иллюстрирует случай, когда нельзя найти достигнутое значение FAR, которое меньше номинального значения FAR 0.0027. Как отмечалось ранее, обычно так бывает при малых значениях \(n\). После него следующий пример для \(n = 13\) иллюстрирует случай, когда достигнутое значение FAR меньше 0.0027 получить можно. Обычно так бывает при больших значениях \(n\), как было видно для знаковой карты.

\emph{Для \(n = 5\).} По записи «правый хвост» для 15 имеем \(\frac{FAR}{2} = 0.031\), что является наименьшей вероятностью, связанной с \(n = 5\). Хотя эта вероятность не меньше и не равна 0.00135, это наименьшая вероятность, которую мы можем получить при выборке объёма 5. Теперь, чтобы получить UCL, то есть \(a\), используем линейное соотношение \(SR_i = 2T_n^+ - n(n + 1)/2\), так что \(a = 2(15) - 15 = 15\) при \(FAR = (2)(0.031) = 0.062\). Таким образом, для \(n = 5\) имеем \(LCL / UCL = \pm 15\) и \(CL = 0\) при соответствующей \(FAR = 0.062\), которая является наименьшей достижимой FAR. Поскольку значение медианы в состоянии IC известно или задано (случай K), \(ARL_{IC} = \frac{1}{FAR} = \frac{1}{0.062} \approx 16\). Опять же, это плохо, поскольку далеко от типичного желаемого значения 370. Однако это наибольшая достижимая \(ARL_{IC}\) при работе с выборкой объёма 5.

\emph{Для \(n = 13\).} По записи «правый хвост» для 86 имеем \(\frac{FAR}{2} = 0.001\), что меньше 0.00135. Теперь, чтобы получить UCL, то есть \(a\), используем линейное соотношение \(SR_i = 2T_n^+ - n(n + 1)/2\), так что \(a = 2(86) - \frac{(13)(14)}{2} = 172 - 91 = 81\) при \(FAR = (2)(0.001) = 0.002\), что меньше 0.0027. Таким образом, для \(n = 13\) имеем \(LCL / UCL = \pm 81\) и \(CL = 0\) при соответствующей \(FAR = 0.002\), которая является наименьшей достижимой FAR. Поскольку значение медианы в состоянии IC известно или задано (случай K), \(ARL_{IC} = \frac{1}{FAR} = \frac{1}{0.002} = 500\), что хорошо, поскольку близко к типичному желаемому значению 500.

\subsubsection{Карта Шухарта по знаковым рангам с правилами сигналов на основе серий}
Как отмечалось ранее, известно, что карты Шухарта недостаточно чувствительны при обнаружении малых сдвигов, и поэтому, как и в случае знаковой карты, мы можем повысить эффективность карты Shewhart-SR при обнаружении меньших сдвигов, добавив правила сигналов на основе серий. Это изучено в Chakraborti and Eryilmaz (2007). Иллюстрация улучшенной карты Shewhart-SR «2 из 2» приведена в разделе 4.3.

\subsubsection{Методы вычисления распределения длины серии}
\paragraph{Точный подход.} Как отмечалось ранее, для контрольных карт Шухарта и некоторых контрольных карт типа Шухарта можно точно вычислить характеристики распределения длины серии. Опять же распределение длины серии контрольной карты Шухарта по знаковым рангам геометрическое с вероятностью успеха, равной вероятности сигнала. Например, читателю рекомендуется обратиться к Bakir (2004), где ARL непараметрической контрольной карты Шухарта по знаковым рангам вычислена точно.

Свойства статистики \(T_n^+\) хорошо известны (см., например, Gibbons and Chakraborti, 2010, с. 196) и приведены в таблице \ref{tab:4.3}.

\begin{table}[!ht]
	\centering
	\caption{Распределительные свойства статистик \(T_n^+\) и \(SR_i\) соответственно.}
	\label{tab:4.3}
	\begin{tabular}{@{}lll@{}}
		\toprule
		 & \(T_n^+\) & \(SR_i\) \\
		\midrule
		Математическое ожидание & \(E(T_n^+) = \dfrac{n(n + 1)}{4}\) & \(E(SR_i) = E\left(2T_n^+ - \dfrac{n(n+1)}{2}\right) = 0\) \\
		\midrule
		Стандартное отклонение & \(SD(T_n^+) = \sqrt{\dfrac{n(n + 1)(2n + 1)}{24}}\) & \(SD(SR_i) = \sqrt{\dfrac{n(n + 1)(2n + 1)}{6}}\) \\
		\bottomrule
	\end{tabular}
\end{table}

Распределение \(SR_i\) можно найти через линейное соотношение \(SR_i = 2T_n^+ - n(n + 1)/2\); оно приведено в последнем столбце таблицы \ref{tab:4.3}.

Распределения вероятностей \(T_n^+\) и \(SR_i\) оба симметричны\footnote{\(T_n^+\) и \(SR_i\) симметрично распределены относительно \(n(n + 1)/4\) и нуля соответственно, пока медиана в состоянии IC остаётся равной \(\mu_0\).} в состоянии IC, когда медиана равна \(\mu_0\). Следовательно, когда процесс находится в состоянии IC, имеем следующее:

\begin{itemize}
	\item распределения вероятностей, задаваемые функциями вероятностей, называются распределениями вероятностей в состоянии IC;
	\item поскольку распределение статистики для нанесения на карту \(SR_i\) в состоянии IC симметрично относительно 0, контрольные границы будут находиться на равных расстояниях от 0, если предполагается, что важность обнаружения восходящего и нисходящего сдвигов одинакова.
\end{itemize}

На рисунке \ref{fig-4.3} показаны распределения вероятностей в состоянии IC статистик \(T_n^+\) и \(SR_i\) для \(n = 5\). Видно, что два дискретных распределения симметричны относительно своих средних, то есть \(T_n^+\) симметрична относительно своего среднего \(\frac{n(n + 1)}{4} = \frac{(5)(6)}{4} = 7.5\), а \(SR_i\) симметрична относительно своего среднего 0.

\begin{figure}[!ht]
	\centering
	\includegraphics[width=0.6\linewidth]{"./images/fig-4.3.png"}
	\caption[]{Распределение вероятностей в состоянии IC статистик \(T_n^+\) и \(SR_i\) для \(n = 5\).}
	\label{fig-4.3}
\end{figure}

Точное распределение длины серии карт Shewhart-SR можно найти теоретически, но это сложнее, чем в случае знаковой карты, потому что распределение статистики для нанесения на карту не имеет замкнутой формы, и поэтому удобнее использовать компьютерные симуляции. Это описано в следующем разделе. Попутно отметим, что точное распределение длины серии непараметрических контрольных карт Shewhart-SR с правилами серий можно также найти с помощью подхода на основе марковских цепей (MC). Это сделано в Chakraborti and Eryilmaz (2007); читателю рекомендуется обратиться к их работе за подробностями.

\paragraph{Симуляционный подход.} Для получения характеристик длины серии карты Shewhart-SR можно использовать следующий алгоритм симуляции. Соответствующая SAS-программа приведена в Приложении B (SAS-программа № 2), а вывод показан ниже на рисунке \ref{fig-4.4}.

\begin{enumerate}
	\item[\textbf{Шаг 1:}] После задания объёма подгруппы (\(n > 1\)) и величины сдвига (\(\delta\)), который нужно обнаружить, генерируем случайные подгруппы из некоторого симметричного непрерывного распределения, например стандартного нормального, Стьюдента \(t\), Лапласа, логистического и т. д. Можно использовать любое симметричное непрерывное распределение.
	\item[\textbf{Шаг 2:}] Вычисляем статистику \(SR_t\) для каждой подгруппы.
	\item[\textbf{Шаг 3:}] Вычисляем контрольные границы, то есть константу для нанесения на карту \(a\), и сравниваем \(SR_t\) с контрольными границами. Вычисление константы для нанесения на карту \(a\) подробно объяснено в разделе 4.2.1.2.
	\item[\textbf{Шаг 4:}] Число подгрупп, необходимое до первого попадания \(SR_t\) на контрольные границы или за них, фиксируется как наблюдение из распределения длины серии.
	\item[\textbf{Шаг 5:}] Повторяем шаги 1--4 в общей сложности 100\,000 раз.
	\item[\textbf{Шаг 6:}] Как только получен «набор данных» из 100\,000 наблюдений из распределения длины серии, для получения характеристик длины серии можно использовать, например, \texttt{proc univariate} SAS. Для \(\delta = 0\) это будут характеристики длины серии в состоянии IC, тогда как для \(\delta \ne 0\) --- характеристики длины серии в состоянии OOC.
\end{enumerate}

\begin{figure}[!ht]
	\centering
	\includegraphics[width=0.7\linewidth]{"./images/fig-4.4.png"}
	\caption[]{Вывод SAS распределения RL при параметрах проектирования (\(n = 13\), \(a = 81\)) карты Shewhart-SR, выбранных для номинальной \(ARL_{IC} = 500\).}
	\label{fig-4.4}
\end{figure}

\subsection{Контрольные карты CUSUM}
Хотя карты Шухарта широко известны и чаще всего используются на практике благодаря своей простоте и общей эффективности, другие классы карт, такие как карты CUSUM, полезны, а иногда более естественно подходят к среде контроля процессов в силу последовательного характера сбора данных. Эти карты, обычно основанные на кумулятивных суммах статистики для нанесения на карту, получаемых по мере накопления данных во времени, известны как более эффективные при обнаружении определённых типов сдвигов процесса.

\subsubsection{Контрольные карты CUSUM на основе знаков}
Идея, лежащая в основе непараметрической, или свободной от распределения, карты CUSUM, в сущности такая же, как и для параметрической карты. Таким образом, карту CUSUM, основанную на статистике знакового критерия (обозначаемую CUSUM-SN), можно получить, заменив статистику \(X_i - \mu_0\) в уравнениях 3.3 и 3.4 на знаковую статистику \(SN_t\). Иными словами, для верхней односторонней карты CUSUM-SN, предназначенной для обнаружения положительных разностей с целевым значением \(\mu_0\) в состоянии IC или известным \(\mu_0\), используем
\begin{equation}
	SN_t^+ = \max(0, SN_{t-1}^+ + SN_t - k) \quad \text{для } t = 1, 2, 3, \ldots
	\label{eq:4.7}
\end{equation}
с начальным значением \(SN_0^+ = 0\) и которая подаёт сигнал в первый момент \(t\), для которого \(SN_t^+ \ge h\).

Аналогично для нижней односторонней карты CUSUM-SN, предназначенной для обнаружения отрицательных разностей с целевым значением \(\mu_0\) в состоянии IC или известным \(\mu_0\), используем
\begin{equation}
	SN_t^- = \min(0, SN_{t-1}^- + SN_t + k) \quad \text{для } t = 1, 2, 3, \ldots
	\label{eq:4.8}
\end{equation}
с начальным значением \(SN_0^- = 0\) и которая подаёт сигнал в первый момент \(t\), для которого \(SN_t^- \le -h\).

Соответствующая двусторонняя карта CUSUM-SN подаёт сигнал в первый момент \(t\), для которого сигнализирует любая из односторонних карт, то есть всякий раз, когда \(SN_t^+ \ge h\) или \(SN_t^- \le -h\). Отметим, что здесь интервал решения \(h > 0\) и референсное значение \(k > 0\) являются параметрами процедуры контрольной карты.

Отметим, что карта CUSUM-SN строится по идее параметрической карты CUSUM \(\overline{X}\). Однако если карта CUSUM \(\overline{X}\) предназначена для среднего и предполагает нормальность, то карта CUSUM-SN предназначена для медианы (или для любого параметра положения) и не требует никаких предположений о лежащем в основе процессе, кроме непрерывности. Когда распределение процесса симметрично, среднее и медиана совпадают, и параметрическая и непараметрическая карты становятся прямыми конкурентами.

\subsubsection{Знаковая контрольная карта CUSUM с правилами сигналов на основе серий}
Хотя правила сигналов на основе серий обычно используются для повышения эффективности контрольных карт Шухарта, Riaz, Abbas and Does (2011) предложили две схемы правил серий для параметрических карт CUSUM. Однако на сегодняшний день дополнение карты CUSUM-SN схемами правил серий не рассматривалось, и это могло бы стать темой дальнейших исследований.

\subsubsection{Методы вычисления распределения длины серии}
\paragraph{Подход на основе марковских цепей.} Момент, в который процедура подаёт сигнал, --- это первый момент, когда марковская цепь с конечным числом состояний \(S_t^+\) попадает в одно из поглощающих состояний, где пространство состояний задаётся как \(\Omega^+ = \{\varsigma_0, \varsigma_1, \ldots, \varsigma_{r+s-1}\}\), \(S_0^+ = 0\) и
\begin{equation}
	S_t^+ = \min\{h, \max\{0, S_{t-1}^+ + SN_t - k\}\}.
	\label{eq:4.9}
\end{equation}

Состояние, соответствующее сигналу карты CUSUM-SN, называется поглощающим. Очевидно, что поглощающее состояние только одно, поскольку карта подаёт сигнал, когда \(S_t^+\) оказывается на \(h\) или выше неё.

Распределение \(SN_t\) легко получить из биномиального распределения; напомним, что \(SN_i = 2T_i - n\) для всех \(i\), где \(T_i\) биномиально распределена с параметрами \(n\) и \(p = P(X_{ij} \ge \mu_0)\). Иллюстративный пример приведён в разделе 4.3.

Иллюстрация использования подхода MC приведена далее. Хотя точное распределение длины серии и различные связанные с ним представляющие интерес характеристики можно найти с помощью подхода MC, обычно удобнее использовать компьютерные симуляции. Это описано ниже.

\paragraph{Симуляционный подход.} Для получения характеристик длины серии карты CUSUM-SN можно использовать следующий алгоритм симуляции. Соответствующая SAS-программа приведена в Приложении B (SAS-программа № 3), а вывод показан на рисунке \ref{fig-4.5}. Для симуляции можно использовать и другое программное обеспечение, например R.

\begin{figure}[!ht]
	\centering
	\includegraphics[width=0.7\linewidth]{"./images/fig-4.5.png"}
	\caption[]{Вывод SAS распределения RL при параметрах проектирования (\(n = 10\), \(k = 6\), \(h = 4\)) карты CUSUM-SN, выбранных для номинальной \(ARL_{IC} = 500\).}
	\label{fig-4.5}
\end{figure}

\begin{enumerate}
	\item[\textbf{Шаг 1:}] После задания объёма подгруппы (\(n \ge 1\)) и величины сдвига (\(\delta\)), который нужно обнаружить, генерируем случайные подгруппы из некоторого непрерывного распределения, например стандартного нормального, Стьюдента \(t\), Лапласа, логистического, гамма и т. д. Можно использовать любое непрерывное распределение, симметричное или асимметричное.
	\item[\textbf{Шаг 2:}] Задаём два параметра проектирования, \(k\) и \(h\) (рекомендации по выбору параметров проектирования читатель найдёт в разделе 3.2.2).
	\item[\textbf{Шаг 3:}] Вычисляем статистики \(SN_t\) и соответствующие статистики для нанесения на карту \(SN_t^-\) и \(SN_t^+\) для каждой подгруппы.
	\item[\textbf{Шаг 4:}] Сравниваем статистики для нанесения на карту с интервалом решения \(h\).
	\item[\textbf{Шаг 5:}] Число подгрупп, необходимое до первого попадания \(SN_t^-\) на \(-h\) или ниже неё либо до первого попадания \(SN_t^+\) на \(h\) или выше неё, фиксируется как наблюдение из распределения длины серии.
	\item[\textbf{Шаг 6:}] Повторяем шаги 1--5 в общей сложности 100\,000 раз.
	\item[\textbf{Шаг 7:}] Как только получен «набор данных» из 100\,000 наблюдений из распределения длины серии, для получения характеристик длины серии можно использовать \texttt{proc univariate} SAS. Для \(\delta = 0\) это будут характеристики длины серии в состоянии IC, тогда как для \(\delta \ne 0\) --- характеристики длины серии в состоянии OOC.
\end{enumerate}

\subsubsection{Контрольные карты CUSUM на основе знаковых рангов}
Как упоминалось ранее при определении карты CUSUM-SN, карты CUSUM полезны и иногда более естественно подходят к среде контроля процессов в силу последовательного характера сбора данных. Эти карты, обычно основанные на кумулятивных суммах статистики для нанесения на карту, получаемых по мере накопления данных во времени, известны как более эффективные при обнаружении определённых типов сдвигов процесса. Здесь мы описываем карты знаковых рангов CUSUM (обозначаемые CUSUM-SR), предложенные Bakir and Reynolds (1979).

\paragraph{Контрольная карта знаковых рангов CUSUM.} Верхняя односторонняя карта CUSUM-SR для обнаружения положительных отклонений от целевого значения \(\mu_0\) в состоянии IC или известного \(\mu_0\) определяется как
\begin{equation}
	SR_t^+ = \max(0, SR_{t-1}^+ + SR_t - k) \quad \text{для } t = 1, 2, 3, \ldots
	\label{eq:4.10}
\end{equation}
с начальным значением \(SR_0^+ = 0\), которая подаёт сигнал в первый момент \(t\), для которого \(SR_t^+ \ge h\).

Аналогично для нижней односторонней карты CUSUM-SR, предназначенной для обнаружения отрицательных отклонений от целевого значения \(\mu_0\) в состоянии IC или известного \(\mu_0\), определяем
\begin{equation}
	SR_t^- = \min(0, SR_{t-1}^- + SR_t - k) \quad \text{для } t = 1, 2, 3, \ldots
	\label{eq:4.11}
\end{equation}
с начальным значением \(SR_0^- = 0\), которая подаёт сигнал в первый момент \(t\), для которого \(SR_t^- \le -h\).

Соответствующая двусторонняя карта CUSUM-SR подаёт сигнал в первый момент \(t\), для которого сигнализирует любая из односторонних карт, то есть всякий раз, когда \(SR_t^+ \ge h\) или \(SR_t^- \le -h\). Отметим, что здесь интервал решения \(h > 0\) и референсное значение \(k > 0\) являются (как и прежде) параметрами процедуры контрольной карты. Эта карта используется для обнаружения сдвига среднего в любом направлении, вверх или вниз.

\paragraph{Контрольная карта знаковых рангов CUSUM с правилами сигналов на основе серий.} Хотя правила сигналов на основе серий обычно используются для повышения эффективности контрольных карт Шухарта, Riaz, Abbas and Does (2011) предложили две схемы правил серий для параметрических карт CUSUM. Однако на сегодняшний день дополнение карты CUSUM-SR схемами правил серий не рассматривалось.

\paragraph{Методы вычисления распределения длины серии.} \emph{Подход на основе марковских цепей.} Момент, в который процедура подаёт сигнал, --- это первый момент, когда марковская цепь с конечным числом состояний \(S_t^+\) попадает в состояние \(\varsigma_{r+s-1}\), где пространство состояний задаётся как \(\Omega^+ = \{\varsigma_0, \varsigma_1, \ldots, \varsigma_{r+s-1}\}\), \(S_0^+ = 0\) и
\begin{equation}
	S_t^+ = \min\{h, \max\{0, S_{t-1}^+ + SR_t - k\}\}.
	\label{eq:4.12}
\end{equation}

Отметим, что, хотя уравнение \eqref{eq:4.12} явно не приведено у Bakir and Reynolds (1979), оно выводится по тому же принципу, что и уравнение \eqref{eq:4.9} для карты CUSUM-SN, а эта формула взята из Amin, Reynolds and Bakir (1995), их уравнение 4.3.

Состояние, соответствующее сигналу карты CUSUM, называется поглощающим. Очевидно, что поглощающее состояние только одно, поскольку карта подаёт сигнал, когда \(S_t^+\) оказывается на \(h\) или выше неё. Распределение \(SR_t\) легко получить из распределения статистики знаковых рангов Уилкоксона \(T_n^+\) (напомним, что \(SR_i = 2T_n^+ - \frac{n(n + 1)}{2}\) для всех \(i\)). Подробный пример приведён в разделе 4.3.

Подробная иллюстрация применения уравнения \eqref{eq:4.12}, то есть использования подхода MC, приведена далее (см. пример 4.10 и, в частности, таблицу \ref{tab:4.13}). Хотя точное распределение длины серии и его представляющие интерес характеристики можно найти с помощью подхода MC, обычно удобнее использовать компьютерную симуляцию. Это описано ниже.

\begin{figure}[!ht]
	\centering
	\includegraphics[width=0.7\linewidth]{"./images/fig-4.6.png"}
	\caption[]{Вывод SAS распределения RL при параметрах проектирования (\(n = 10\), \(k = 41\), \(h = 24\)) карты CUSUM-SR, выбранных так, чтобы номинальная \(ARL_{IC} = 500\).}
	\label{fig-4.6}
\end{figure}

\emph{Симуляционный подход.} Для получения характеристик длины серии карты CUSUM-SR можно использовать следующий алгоритм симуляции. Соответствующая SAS-программа приведена в Приложении B (SAS-программа № 4), а вывод показан на рисунке \ref{fig-4.6}. Для симуляции можно использовать и другое программное обеспечение, например R.

\begin{enumerate}
	\item[\textbf{Шаг 1:}] После задания объёма подгруппы (\(n > 1\)) и величины сдвига (\(\delta\)), который нужно обнаружить, генерируем случайные подгруппы из некоторого симметричного непрерывного распределения, например стандартного нормального, Стьюдента \(t\), Лапласа, логистического и т. д. Можно использовать любое симметричное непрерывное распределение.
	\item[\textbf{Шаг 2:}] Задаём два параметра проектирования, \(k\) и \(h\) (рекомендации по выбору параметров проектирования читатель найдёт в разделе 3.2.2).
	\item[\textbf{Шаг 3:}] Вычисляем статистики \(SR_t\) и соответствующие статистики для нанесения на карту \(SR_t^-\) и \(SR_t^+\) для каждой подгруппы.
	\item[\textbf{Шаг 4:}] Сравниваем статистики для нанесения на карту с интервалом решения \(h\).
	\item[\textbf{Шаг 5:}] Число подгрупп, необходимое до первого попадания \(SR_t^-\) на \(-h\) или ниже неё либо до первого попадания \(SR_t^+\) на \(h\) или выше неё, фиксируется как наблюдение из распределения длины серии.
	\item[\textbf{Шаг 6:}] Повторяем шаги 1--5 в общей сложности 100\,000 раз.
	\item[\textbf{Шаг 7:}] Как только получен «набор данных» из 100\,000 наблюдений из распределения длины серии, для получения характеристик длины серии можно использовать \texttt{proc univariate} SAS. Для \(\delta = 0\) это будут характеристики длины серии в состоянии IC, тогда как для \(\delta \ne 0\) --- характеристики длины серии в состоянии OOC.
\end{enumerate}

В дополнение к картам Шухарта и CUSUM, как было отмечено для параметрических карт, ещё одним популярным классом используемых на практике контрольных карт являются карты EWMA. Карты EWMA также используют последовательный (упорядоченный во времени) накопительный характер данных, возникающих в типичной среде SPC, и известны как эффективные при обнаружении меньших сдвигов, но при этом проще в реализации, чем карты CUSUM. Ниже мы обсуждаем свободные от распределения карты EWMA, основанные на статистиках SN и SR.

\subsection{Контрольные карты EWMA}
\subsubsection{Контрольные карты EWMA на основе знаков}
Сначала обсудим знаковую контрольную карту EWMA (обозначаемую EWMA-SN), рассматривая отдельные наблюдения. Опять же, следуя идее параметрической карты EWMA, процедура контрольной карты EWMA, основанная на знаковой статистике \(SN_i\), накапливает знаковые статистики \(SN_1, SN_2, SN_3, \ldots\). Напомним, что для отдельных измерений \(SN_i\) может принимать только значения \(-1\) или \(+1\). Статистика для нанесения на карту EWMA-SN определяется как
\begin{equation}
	Z_t = \lambda SN_t + (1 - \lambda) Z_{t-1}
	\label{eq:4.13}
\end{equation}
с начальным значением \(Z_0 = 0\) и весовой константой \(0 < \lambda \le 1\).

Центральная линия и точные контрольные границы двусторонней карты EWMA-SN, основанной на отдельных измерениях, задаются как
\begin{equation}
	\begin{aligned}
		UCL &= L\sqrt{\frac{\lambda}{2 - \lambda} \left(1 - (1 - \lambda)^{2t}\right)} \\
		CL &= 0 \\
		LCL &= -L\sqrt{\frac{\lambda}{2 - \lambda} \left(1 - (1 - \lambda)^{2t}\right)}
	\end{aligned}
	\label{eq:4.14}
\end{equation}

Контрольные границы упрощаются при \(t \to \infty\), то есть для больших значений \(t\), и имеют вид
\begin{equation}
	\begin{aligned}
		UCL &= +L\sqrt{\frac{\lambda}{2 - \lambda}} \\
		LCL &= -L\sqrt{\frac{\lambda}{2 - \lambda}}
	\end{aligned}
	\label{eq:4.15}
\end{equation}

Эти контрольные границы называются асимптотическими, или стационарными, контрольными границами; они применимы, когда процесс достиг стационарного состояния после работы в течение некоторого времени. Процесс считается находящимся в состоянии IC, когда статистика для нанесения на карту \(Z_t\), \(t = 1, 2, 3, \ldots\), находится между двумя контрольными границами, но как только статистика для нанесения на карту оказывается на UCL или LCL либо за ними, процесс объявляется вышедшим из-под контроля, и обычно начинается поиск назначаемых причин. Вывод точных контрольных границ приведён ниже для более общего случая, когда объём подгруппы \(n\) можно положить равным 1, чтобы получить приведённое выше выражение.

Двустороннюю EWMA можно модифицировать, чтобы образовать одностороннюю статистику, примерно так же, как двустороннюю CUSUM можно превратить в одностороннюю статистику. Верхняя односторонняя EWMA, например, задаётся как \(Z_t^+ = \max[0, \lambda SN_t + (1 - \lambda)Z_{t-1}]\) для \(t = 1, 2, 3, \ldots\) с начальным значением \(z_0^+ = 0\). Если статистика для нанесения на карту \(Z_t\) оказывается на \(UCL\) или выше неё, процесс считается вышедшим из-под контроля и необходим поиск назначаемых причин. Отметим, что статистики для нанесения на карту для нижней односторонней карты EWMA можно построить аналогичным образом. Это оставляется читателю в качестве упражнения.

Далее рассмотрим случай сгруппированных в подгруппы данных. Статистика для нанесения на карту EWMA-SN, основанная на сгруппированных данных, аналогична статистике карты EWMA-SN, основанной на отдельных измерениях. Однако теперь \(SN_t\) может принимать целые значения от \(-n\) до \(+n\) (включительно). Центральная линия (CL) и точные контрольные границы двусторонней карты EWMA-SN, основанной на сгруппированных данных, задаются как
\begin{equation}
	\begin{aligned}
		UCL &= L\sqrt{\frac{\lambda n}{2 - \lambda} \left(1 - (1 - \lambda)^{2t}\right)} \\
		CL &= 0 \\
		LCL &= -L\sqrt{\frac{\lambda n}{2 - \lambda} \left(1 - (1 - \lambda)^{2t}\right)}
	\end{aligned}
	\label{eq:4.16}
\end{equation}

Контрольные границы упрощаются при \(t \to \infty\), то есть для больших значений \(t\), и имеют вид
\begin{equation}
	\begin{aligned}
		UCL &= +L\sqrt{\frac{\lambda n}{2 - \lambda}} \\
		LCL &= -L\sqrt{\frac{\lambda n}{2 - \lambda}}
	\end{aligned}
	\label{eq:4.17}
\end{equation}

Эти контрольные границы называются асимптотическими, или стационарными, контрольными границами; они применимы, когда процесс достиг стационарного состояния после работы в течение некоторого времени. Как и в случае отдельных измерений, процесс считается находящимся в состоянии IC, когда статистика для нанесения на карту \(Z_t\), \(t = 1, 2, 3, \ldots\), находится между двумя контрольными границами. Однако как только статистика для нанесения на карту оказывается на UCL или LCL либо за ними, процесс объявляется вышедшим из-под контроля, и обычно начинается поиск назначаемых причин.

\subsubsection{Контрольные карты EWMA на основе знаков с правилами сигналов на основе серий}
Хотя правила сигналов на основе серий обычно используются для повышения эффективности контрольных карт Шухарта, Abbas, Riaz and Does (2011) предложили две схемы правил серий для параметрических карт EWMA. Однако на сегодняшний день дополнение карты EWMA-SN схемами правил серий не рассматривалось, и это тема для дальнейших исследований.

\subsubsection{Методы вычисления распределения длины серии}
\paragraph{Подход на основе марковских цепей.} Напомним, что нас интересует вычисление одношаговых переходных вероятностей \(M = [p_{ij}]\) для \(i, j = -s, -s + 1, \ldots, s - 1, s\). Для вычисления этих вероятностей предполагаем, что статистика для нанесения на карту равна \(S_i\) всякий раз, когда она находится в состоянии \(i\). Для не поглощающего \(j\) получаем
\begin{equation}
	\begin{aligned}
		p_{ij} &= P(S_j - \tau < Z_k \le S_j + \tau \mid Z_{k-1} = S_i) \quad \forall j\, (j \ne s) \\
		\text{и}\quad p_{ij} &= P(S_j - \tau < Z_k < S_j + \tau \mid Z_{k-1} = S_i) \quad \text{для } j = s.
	\end{aligned}
	\label{eq:4.18}
\end{equation}

Отметим, что подстрочный индекс \(k\) не означает, что статистика для нанесения на карту зависит от времени; он просто относится к рассматриваемой статистике для нанесения на карту. Для поглощающего \(j\) используем тот факт, что каждая строка матрицы \(M\) должна в сумме давать единицу, и поэтому вероятности перехода из невозвратного состояния в поглощающее находятся вычитанием. Нахождение существенной матрицы переходных вероятностей \(Q\) включает нахождение переходных вероятностей из одного невозвратного состояния в другое. Для вычисления переходных вероятностей карты EWMA-SN используем статистику для нанесения на карту и подставляем её в уравнение \eqref{eq:4.18}, чтобы получить
\begin{equation}
	\begin{aligned}
		p_{ij} &= P\left( \frac{(S_j - \tau) - (1 - \lambda)S_i}{\lambda} + n \Big/ 2 < T_k \right. \\
		&\qquad \left. \le \frac{(S_j + \tau) - (1 - \lambda)S_i}{\lambda} + n \Big/ 2 \right) \quad \forall j\, (j \ne s) \\
		\text{и}\quad p_{ij} &= P\left( \frac{(S_j - \tau) - (1 - \lambda)S_i}{\lambda} + n \Big/ 2 < T_k \right. \\
		&\qquad \left. < \frac{(S_j + \tau) - (1 - \lambda)S_i}{\lambda} + n \Big/ 2 \right) \quad \text{для } j = s.
	\end{aligned}
	\label{eq:4.19}
\end{equation}

Поскольку значения \(n\), \(\tau\), \(\lambda\), \(S_i\) и \(S_j\) являются известными постоянными (\(\lambda\) и \(n\) выбираются практиком, а \(\tau\), \(S_i\) и \(S_j\) вычисляются), биномиальные вероятности в уравнении \eqref{eq:4.19} легко вычислить. Доказательство уравнения \eqref{eq:4.19} приведено ниже, а иллюстративный пример дан в разделе 4.3.

\emph{Доказательство.} Для не поглощающего \(j\), то есть \(j = -s, -s + 1, \ldots, s - 1\), получаем одношаговые переходные вероятности \(p_{ij} = P(S_j - \tau < Z_k \le S_j + \tau \mid Z_{k-1} = S_i)\). Это вероятность того, что \(Z_k\) находится внутри состояния \(j\) при условии, что \(Z_{k-1}\) равна середине состояния \(i\). Используем определение статистики для нанесения на карту и тот факт, что \(SN_k = 2T_k - n\), где \(T_k\) биномиально распределена с параметрами \(n\) и \(p = P(X_{ij} > \mu_0)\). Эту переходную вероятность можно записать как
\begin{align*}
	p_{ij} &= P(S_j - \tau < \lambda SN_k + (1 - \lambda)Z_{k-1} \le S_j + \tau \mid Z_{k-1} = S_i) \\
	&= P(S_j - \tau < \lambda SN_k + (1 - \lambda)S_i \le S_j + \tau) \\
	&= P\left( \frac{(S_j - \tau) - (1 - \lambda)S_i}{\lambda} < SN_k \le \frac{(S_j + \tau) - (1 - \lambda)S_i}{\lambda} \right) \\
	&= P\left( \frac{(S_j - \tau) - (1 - \lambda)S_i}{\lambda} < 2T_k - n \le \frac{(S_j + \tau) - (1 - \lambda)S_i}{\lambda} \right) \\
	&= P\left( \frac{(S_j - \tau) - (1 - \lambda)S_i}{\lambda} + n \Big/2 < T_k \le \frac{(S_j + \tau) - (1 - \lambda)S_i}{\lambda} + n \Big/2 \right).
\end{align*}

Как только \(p_{ij}\) вычислены, подход MC можно использовать для вычисления точного распределения длины серии; иллюстрация приведена в разделе 4.3. Однако, как мы отмечали ранее, обычно удобнее использовать компьютерные симуляции, что описано ниже. \(\square\)

\paragraph{Симуляционный подход.} Для получения характеристик длины серии карты EWMA-SN можно использовать следующий алгоритм симуляции. Соответствующая SAS-программа приведена в Приложении B (SAS-программа № 5), а вывод показан на рисунке \ref{fig-4.7}.

\begin{figure}[!ht]
	\centering
	\includegraphics[width=0.7\linewidth]{"./images/fig-4.7.png"}
	\caption[]{Вывод SAS распределения RL при параметрах проектирования (\(n = 1\), \(\lambda = 0.10\), \(L = 2.667\)) карты EWMA-SN, выбранных для номинальной \(ARL_{IC} = 500\).}
	\label{fig-4.7}
\end{figure}

\begin{enumerate}
	\item[\textbf{Шаг 1:}] После задания объёма подгруппы (\(n \ge 1\)) и величины сдвига (\(\delta\)), который нужно обнаружить, генерируем случайные подгруппы из некоторого непрерывного распределения, например стандартного нормального, Стьюдента \(t\), Лапласа, логистического, гамма и т. д. Можно использовать любое непрерывное распределение, симметричное или асимметричное.
	\item[\textbf{Шаг 2:}] Задаём два параметра проектирования, \(\lambda\) и \(L\) (рекомендации по выбору параметров проектирования читатель найдёт в разделе 3.2.3).
	\item[\textbf{Шаг 3:}] Вычисляем статистики \(SN_t\) и соответствующие статистики для нанесения на карту \(Z_t\) для каждой подгруппы.
	\item[\textbf{Шаг 4:}] Вычисляем точные контрольные границы (см. уравнение \eqref{eq:4.16}) или стационарные контрольные границы (см. уравнение \eqref{eq:4.17}) и сравниваем \(Z_t\) с контрольными границами.
	\item[\textbf{Шаг 5:}] Число подгрупп, необходимое до первого попадания \(Z_t\) на контрольные границы или за них, фиксируется как наблюдение из распределения длины серии.
	\item[\textbf{Шаг 6:}] Повторяем шаги 1--5 в общей сложности 100\,000 раз.
	\item[\textbf{Шаг 7:}] Как только получен «набор данных» из 100\,000 наблюдений из распределения длины серии, для получения характеристик длины серии можно использовать \texttt{proc univariate} SAS. Для \(\delta = 0\) это будут характеристики длины серии в состоянии IC, тогда как для \(\delta \ne 0\) --- характеристики длины серии в состоянии OOC.
\end{enumerate}

\noindent\textbf{Замечание 4.3. Число подынтервалов между контрольными границами.}

При реализации подхода MC для карты EWMA-SN возникает вопрос о числе подынтервалов, то есть о значении \(\nu\), которое следует использовать. Когда \(\nu\) достаточно велико, подход MC даёт эффективный метод, позволяющий точно аппроксимировать характеристики длины серии схемы EWMA. На практике значения \(\nu\) около 100 дают удовлетворительные приближения (см. Yu, 2007, с. 20). Здесь мы проверяем это утверждение на примере карты EWMA-SN.

Проверка подхода MC с помощью симуляции Монте-Карло показывает, что при \(\nu = 1001\) расхождения лежат в пределах 2\,\% от симулированных значений (см. таблицу \ref{tab:4.4}). Большие значения \(\nu\) давали бы более точные ответы, но при этом некоторые характеристики длины серии нельзя было бы вычислить за практически приемлемое время.

Эти результаты аналогичны найденным для параметрической (нормальной теории) карты EWMA (см. таблицу 3.12), где было найдено, что чем больше значение \(\nu\), тем меньше процентное различие между симулированными значениями и значениями, полученными с помощью подхода MC.

Из таблицы \ref{tab:4.4} видно, что, как и следовало ожидать, чем больше значение \(\nu\), тем меньше процентное различие между симулированными значениями и значениями, полученными с помощью подхода MC. Когда параметры проектирования, например, взяты равными \(\lambda = 0.05\) и \(L = 2.472\), имеем большое процентное различие 18.62\,\%, когда \(\nu\) мало (\(\nu = 51\)), в отличие от малого процентного различия 0.18\,\%, когда \(\nu\) велико (\(\nu = 1001\)). По-видимому, \(\nu \ge 501\) должно быть вполне удовлетворительным на практике.

\begin{table}[!ht]
	\centering
	\caption{Значения \(ARL_{IC}\) и абсолютное процентное различие между симулированными значениями и значениями, полученными с помощью подхода MC, для карты EWMA-SN при различных значениях \(\nu\).}
	\label{tab:4.4}
	\par\smallskip
	\footnotesize Для справедливости сравнений как в симуляциях, так и в подходе MC использовались стационарные контрольные границы.
	\par\smallskip
	\begin{tabular}{@{}lcccc@{}}
		\toprule
		& \multicolumn{2}{c}{\(\lambda = 0.05\), \(L = 2.472\)} & \multicolumn{2}{c}{\(\lambda = 0.10\), \(L = 2.588\)} \\
		& Достигнутая \(ARL\) & Разл. (\%) & Достигнутая \(ARL\) & Разл. (\%) \\
		\midrule
		51   & 439.07 & (18.62\%) & 384.92 & (4.06\%) \\
		59   & 386.23 & (4.34\%)  & 420.36 & (13.64\%) \\
		67   & 379.34 & (2.48\%)  & 337.55 & (8.75\%) \\
		75   & 410.76 & (10.97\%) & 360.21 & (2.62\%) \\
		83   & 410.48 & (10.90\%) & 423.25 & (14.42\%) \\
		101  & 396.12 & (7.02\%)  & 357.42 & (3.38\%) \\
		151  & 380.61 & (2.83\%)  & 358.61 & (3.05\%) \\
		201  & 391.25 & (5.70\%)  & 330.66 & (10.61\%) \\
		251  & 368.50 & (0.45\%)  & 354.17 & (4.26\%) \\
		301  & 369.42 & (0.20\%)  & 379.63 & (2.63\%) \\
		401  & 365.71 & (1.20\%)  & 370.68 & (0.21\%) \\
		501  & 377.47 & (1.98\%)  & 377.33 & (2.01\%) \\
		1001 & 369.49 & (0.18\%)  & 373.80 & (1.05\%) \\
		\midrule
		100\,000 симуляций & 370.15 & & 369.91 & \\
		\bottomrule
	\end{tabular}
\end{table}

\subsubsection{Контрольные карты EWMA на основе знаковых рангов}
Как и карту EWMA-SN, можно рассмотреть карту EWMA, основанную на знаковых рангах и обозначаемую EWMA-SR. Карта EWMA-SR впервые была введена Amin and Searcy (1991); эта карта обсуждается далее в этом разделе.

Карта EWMA-SR основана на знаково-ранговой статистике \(SR_i\), которая определена в уравнении \eqref{eq:4.5}, и в общей формуле EWMA накапливает статистики \(SR_1, SR_2, SR_3, \ldots\). Таким образом, карта EWMA-SR определяется как
\begin{equation}
	Z_t = \lambda SR_t + (1 - \lambda) Z_{t-1}
	\label{eq:4.20}
\end{equation}
с начальным значением \(Z_0 = 0\) и весовой константой \(0 < \lambda \le 1\).

Центральная линия и точные контрольные границы задаются как
\begin{equation}
	\begin{aligned}
		UCL &= L\sigma_{SR_i}\sqrt{\frac{\lambda}{2 - \lambda} \left(1 - (1 - \lambda)^{2t}\right)} \\
		CL &= 0 \\
		LCL &= -L\sigma_{SR_i}\sqrt{\frac{\lambda}{2 - \lambda} \left(1 - (1 - \lambda)^{2t}\right)}
	\end{aligned}
	\label{eq:4.21}
\end{equation}
тогда как асимптотические, или стационарные, контрольные границы, то есть при \(t \to \infty\), задаются как
\begin{equation}
	\begin{aligned}
		UCL &= L\sigma_{SR_i}\sqrt{\frac{\lambda}{2 - \lambda}} \\
		LCL &= -L\sigma_{SR_i}\sqrt{\frac{\lambda}{2 - \lambda}}
	\end{aligned}
	\label{eq:4.22}
\end{equation}
где \(\sigma_{SR_t} = \sqrt{\frac{n(n + 1)(2n + 1)}{6}}\) --- стандартное отклонение статистики \(SR_i\) в состоянии IC, если внутри выборки или подгруппы нет связок (обсуждение свойств статистики знаковых рангов Уилкоксона и статистики \(SR_i\) см. в разделе 4.2.1.2, а о связках --- в разделе 4.2.1.1). Отметим, что \(\sigma_{SR_i}\) получается с использованием линейного соотношения между \(SR_i\) и \(T_n^+\):
\begin{align*}
	\sigma_{SR_t}^2 &= VAR(SR_t) \\
	&= VAR\left(2T_n^+ - \frac{n(n + 1)}{2}\right) = 4\left(\frac{n(n + 1)(2n + 1)}{24}\right) \\
	&= \frac{n(n + 1)(2n + 1)}{6}.
\end{align*}
Процесс считается находящимся в состоянии IC, пока все статистики для нанесения на карту \(Z_t\), \(t = 1, 2, 3, \ldots\), лежат между двумя контрольными границами, но как только какая-либо статистика для нанесения на карту оказывается на UCL или LCL либо за ними, процесс объявляется вышедшим из-под контроля, и обычно начинается поиск назначаемых причин.

Читателю рекомендуется обратиться к разделу 3.2.3 за подробным обсуждением выбора параметров проектирования \(\lambda\) и \(L\). Коротко рекомендация состоит в том, чтобы выбирать малое \(\lambda\), скажем равное 0.05, когда представляют интерес малые сдвиги; выбирать \(\lambda = 0.10\), если речь идёт об умеренных сдвигах; и выбирать \(\lambda = 0.20\), если интересны большие сдвиги. Таким образом, выбор \(\lambda = 0.05\) означает, что представляют интерес малые сдвиги. После того как \(\lambda\) выбрана, второй шаг заключается в нахождении \(L\) так, чтобы достигалась желаемая \(ARL_{IC}\).

Одностороннюю карту EWMA-SR можно построить из двусторонней карты EWMA примерно так же, как для карты CUSUM; например, верхняя односторонняя статистика для нанесения на карту EWMA задаётся как \(Z_i^+ = \max[0, \lambda SR_t + (1 - \lambda)Z_{t-1}]\) для \(t = 1, 2, 3, \ldots\) с начальным значением \(Z_0^+ = 0\). Если статистика для нанесения на карту \(Z_t\) оказывается на \(UCL\) или выше неё, процесс считается вышедшим из-под контроля и необходим поиск назначаемых причин. Отметим, что UCL для верхней односторонней карты EWMA, вероятно, будет отличаться от UCL двусторонней карты EWMA. Нижняя односторонняя статистика для нанесения на карту EWMA определяется аналогично. Вывод точных контрольных границ приведён в Приложении 4.3.

\subsubsection{Контрольная карта EWMA-SR с правилами сигналов на основе серий}
Хотя правила сигналов на основе серий обычно используются для повышения эффективности контрольных карт Шухарта, Abbas, Riaz and Does (2011) предложили две схемы правил серий для параметрических карт EWMA. Однако на сегодняшний день дополнение карты EWMA-SR схемами правил серий не рассматривалось. Это ещё одна тема для дальнейших исследований.

\subsubsection{Методы вычисления распределения длины серии}
\paragraph{Подход на основе марковских цепей.} Опять же, как и в случае карты CUSUM, можно использовать подход MC для вычисления точного распределения длины серии карты EWMA; для этого напомним, что нас интересует вычисление одношаговых переходных вероятностей \(M = [p_{ij}]\) для \(i, j = -s, -s + 1, \ldots, s - 1, s\). Для вычисления этих вероятностей предполагаем, что статистика для нанесения на карту равна \(S_i\) всякий раз, когда она находится в состоянии \(i\). Для не поглощающего \(j\) получаем
\begin{equation}
	\begin{aligned}
		p_{ij} &= P(S_j - \tau < Z_k \le S_j + \tau \mid Z_{k-1} = S_i) \quad \forall j\, (j \ne s) \\
		\text{и}\quad p_{ij} &= P(S_j - \tau < Z_k < S_j + \tau \mid Z_{k-1} = S_i) \quad \text{для } j = s.
	\end{aligned}
	\label{eq:4.23}
\end{equation}
Отметим, что подстрочный индекс \(k\) не означает, что статистика для нанесения на карту зависит от времени; он просто относится к рассматриваемой статистике для нанесения на карту. Для поглощающего \(j\) используем тот факт, что каждая строка матрицы \(M\) должна в сумме давать единицу, и поэтому вероятности перехода из невозвратного состояния в поглощающее находятся вычитанием. Нахождение существенной матрицы переходных вероятностей \(Q\) включает нахождение переходных вероятностей из одного невозвратного состояния в другое. Для вычисления переходных вероятностей карты EWMA-SR используем статистику для нанесения на карту и подставляем её в уравнение \eqref{eq:4.23}, чтобы получить
\begin{equation}
	\begin{aligned}
		p_{ij} &= P\left( \frac{(S_j - \tau) - (1 - \lambda)S_i}{\lambda} + \frac{n(n + 1)}{2} \Big/ 2 < T_n^+ \right. \\
		&\qquad \left. \le \frac{(S_j + \tau) - (1 - \lambda)S_i}{\lambda} + \frac{n(n + 1)}{2} \Big/ 2 \right) \quad \forall j\, (j \ne s) \\
		\text{и}\quad p_{ij} &= P\left( \frac{(S_j - \tau) - (1 - \lambda)S_i}{\lambda} + \frac{n(n + 1)}{2} \Big/ 2 < T_n^+ \right. \\
		&\qquad \left. < \frac{(S_j + \tau) - (1 - \lambda)S_i}{\lambda} + \frac{n(n + 1)}{2} \Big/ 2 \right) \quad \text{для } j = s.
	\end{aligned}
	\label{eq:4.24}
\end{equation}

Поскольку значения \(n\), \(\tau\), \(\lambda\), \(S_i\) и \(S_j\) являются известными постоянными (\(\lambda\) и \(n\) выбираются практиком, а \(\tau\), \(S_i\) и \(S_j\) вычисляются), вероятности знаковых рангов Уилкоксона в уравнении \eqref{eq:4.24} легко вычислить. Доказательство уравнения \eqref{eq:4.24} приведено ниже, а иллюстративный пример дан в разделе 4.3.

\emph{Доказательство.} Для не поглощающего \(j\), то есть \(j = -s, -s + 1, \ldots, s - 1\), получаем одношаговые переходные вероятности \(p_{ij} = P(S_j - \tau < Z_k \le S_j + \tau \mid Z_{k-1} = S_i)\). Это вероятность того, что \(Z_k\) находится внутри состояния \(j\) при условии, что \(Z_{k-1}\) равна середине состояния \(i\). Используя определение статистики для нанесения на карту и тот факт, что \(SR_k = 2T_n^+ - \frac{n(n + 1)}{2}\), где \(T_n^+\) --- известная статистика знаковых рангов Уилкоксона, эту переходную вероятность можно записать как
\begin{align*}
	p_{ij} &= P(S_j - \tau < \lambda SR_k + (1 - \lambda)Z_{k-1} \le S_j + \tau \mid Z_{k-1} = S_i) \\
	&= P(S_j - \tau < \lambda SR_k + (1 - \lambda)S_i \le S_j + \tau) \\
	&= P\left( \frac{(S_j - \tau) - (1 - \lambda)S_i}{\lambda} < SR_k \le \frac{(S_j + \tau) - (1 - \lambda)S_i}{\lambda} \right) \\
	&= P\left( \frac{(S_j - \tau) - (1 - \lambda)S_i}{\lambda} < 2T_n^+ - \frac{n(n+1)}{2} \le \frac{(S_j + \tau) - (1 - \lambda)S_i}{\lambda} \right) \\
	&= P\left( \frac{(S_j - \tau) - (1 - \lambda)S_i}{\lambda} + \frac{n(n+1)}{2}\Big/2 < T_n^+ \right. \\
	&\qquad \left. \le \frac{(S_j + \tau) - (1 - \lambda)S_i}{\lambda} + \frac{n(n+1)}{2}\Big/2 \right).
\end{align*}

\begin{figure}[!ht]
	\centering
	\includegraphics[width=0.7\linewidth]{"./images/fig-4.8.png"}
	\caption[]{Вывод SAS распределения RL при параметрах проектирования (\(n = 10\), \(\lambda = 0.1\), \(L = 2.794\)) карты EWMA-SR, выбранных для номинальной \(ARL_{IC} = 500\).}
	\label{fig-4.8}
\end{figure}

Как только \(p_{ij}\) вычислены, подход MC можно использовать для вычисления точного распределения длины серии; иллюстрация приведена в разделе 4.3. Однако, как мы отмечали ранее, обычно удобнее использовать компьютерные симуляции, что описано ниже. \(\square\)

\paragraph{Симуляционный подход.} Для получения характеристик длины серии карты EWMA-SR можно использовать следующий алгоритм симуляции. Соответствующая SAS-программа приведена в Приложении B (SAS-программа № 6), а вывод показан ниже на рисунке \ref{fig-4.8}.

\begin{enumerate}
	\item[\textbf{Шаг 1:}] После задания объёма подгруппы (\(n > 1\)) и величины сдвига (\(\delta\)), который нужно обнаружить, генерируем случайные подгруппы из некоторого симметричного непрерывного распределения, например стандартного нормального, Стьюдента \(t\), Лапласа, логистического и т. д. Можно использовать любое симметричное непрерывное распределение.
	\item[\textbf{Шаг 2:}] Задаём два параметра проектирования, \(\lambda\) и \(L\) (рекомендации по выбору параметров проектирования читатель найдёт в разделе 3.2.3).
	\item[\textbf{Шаг 3:}] Вычисляем статистики \(SR_t\) и соответствующие статистики для нанесения на карту \(Z_t\) для каждой подгруппы.
	\item[\textbf{Шаг 4:}] Вычисляем точные контрольные границы (см. уравнение \eqref{eq:4.21}) или стационарные контрольные границы (см. уравнение \eqref{eq:4.22}) и сравниваем \(Z_t\) с контрольными границами.
	\item[\textbf{Шаг 5:}] Число подгрупп, необходимое до первого попадания \(Z_t\) на контрольные границы или за них, фиксируется как наблюдение из распределения длины серии.
	\item[\textbf{Шаг 6:}] Повторяем шаги 1--5 в общей сложности 100\,000 раз.
	\item[\textbf{Шаг 7:}] Как только получен «набор данных» из 100\,000 наблюдений из распределения длины серии, для получения характеристик длины серии можно использовать \texttt{proc univariate} SAS. Для \(\delta = 0\) это будут характеристики длины серии в состоянии IC, тогда как для \(\delta \ne 0\) --- характеристики длины серии в состоянии OOC.
\end{enumerate}

\noindent\textbf{Замечание 4.4. Число подынтервалов между контрольными границами.}

При реализации подхода MC для карты EWMA-SR возникает вопрос о числе подынтервалов, то есть о значении \(\nu\). Когда число подынтервалов \(\nu\) достаточно велико, подход MC даёт эффективный метод, позволяющий точно аппроксимировать характеристики длины серии схемы EWMA. В разделе 4.2.3 было показано, что для карты EWMA-SN большие значения \(\nu\) давали бы более точные ответы, но при этом некоторые характеристики длины серии нельзя было бы вычислить за практически приемлемое время. Те же выводы справедливы и для карты EWMA-SR.

\section{Свободные от распределения контрольные карты в случае K: иллюстративные примеры}
\subsection{Контрольные карты Шухарта}

\noindent\textbf{Пример 4.1. Карта Shewhart-SN по отдельным измерениям.}

Столбец (a) таблицы \ref{tab:4.5} показывает набор данных из Montgomery (2009, табл.~9.1, с.~401). В нём 30 наблюдений, первые 20 из которых взяты случайно из нормального распределения со средним \(\mu_0 = 10\) и стандартным отклонением \(\sigma_0 = 1\). Остальные 10 наблюдений взяты из нормального распределения со средним \(\mu_1 = 11\) и стандартным отклонением \(\sigma_0 = 1\). Таким образом, \(\mu_1 = \mu_0 + \delta\sigma_0 = 10 + (1)(1)\), так что \(\delta = 1\). Столбец (b) таблицы \ref{tab:4.5} содержит значения статистики для нанесения на карту \(SN_i\) для \(i = 1, 2, \ldots, 30\). Первое наблюдение, например, равно 9.45, что меньше 10. Следовательно, \(SN_1 = -1\).

\begin{table}[!ht]
	\centering
	\caption{Данные для карты Shewhart-SN по отдельным измерениям и статистики \(SN\).}
	\label{tab:4.5}
	\begin{tabular}{@{}ccc@{}}\toprule
		Номер наблюдения & (a) & (b) \\ i & \(X_i\) & \(SN_i\) \\ \midrule
		1 & 9.45 & \(-1\) \\ 2 & 7.99 & \(-1\) \\ 3 & 9.29 & \(-1\) \\ 4 & 11.66 & \(1\) \\
		5 & 12.16 & \(1\) \\ 6 & 10.18 & \(1\) \\ 7 & 8.04 & \(-1\) \\ 8 & 11.46 & \(1\) \\
		9 & 9.20 & \(-1\) \\ 10 & 10.34 & \(1\) \\ 11 & 9.03 & \(-1\) \\ 12 & 11.47 & \(1\) \\
		13 & 10.51 & \(1\) \\ 14 & 9.40 & \(-1\) \\ 15 & 10.08 & \(1\) \\ 16 & 9.37 & \(-1\) \\
		17 & 10.62 & \(1\) \\ 18 & 10.31 & \(1\) \\ 19 & 8.52 & \(-1\) \\ 20 & 10.84 & \(1\) \\
		21 & 10.90 & \(1\) \\ 22 & 9.33 & \(-1\) \\ 23 & 12.29 & \(1\) \\ 24 & 11.50 & \(1\) \\
		25 & 10.60 & \(1\) \\ 26 & 11.08 & \(1\) \\ 27 & 10.38 & \(1\) \\ 28 & 11.62 & \(1\) \\
		29 & 11.31 & \(1\) \\ 30 & 10.52 & \(1\) \\ \bottomrule
	\end{tabular}
\end{table}

Для отдельных измерений (данных) по записи «правый хвост» для 1 в таблице A Приложения A имеем \(\frac{FAR}{2} = 0.5\), что является наименьшей вероятностью, связанной с \(n = 1\). Хотя эта вероятность не меньше и не равна 0.00135, это наименьшая вероятность, которую можно получить при работе с отдельными измерениями. Теперь, чтобы получить UCL, то есть \(a\), используем линейное соотношение \(SN_i = 2T_i - n\), так что \(a = 2(1) - 1 = 1\) при \(FAR = (2)(0.5) = 1\). Таким образом, для \(n = 1\) имеем \(LCL / UCL = \pm 1\) и \(CL = 0\) при соответствующей \(FAR = 1\), которая является наименьшей достижимой FAR. Поскольку значение медианы в состоянии IC известно или задано (случай K), \(ARL_{IC} = \frac{1}{FAR} = \frac{1}{1} = 1\). Опять же, это неприемлемо, поскольку далеко от типичного желаемого значения 370. Фактически мы получаем наибольшую возможную FAR и наименьшую возможную \(ARL_{IC}\).

Таким образом, контрольные границы и CL имеют вид
\[
\begin{gathered}
	UCL = +1 \\
	CL = 0 \\
	LCL = -1
\end{gathered}
\]
и соответствующая карта Shewhart-SN показана на рисунке \ref{fig-4.9}.

Поскольку мы получаем наибольшую возможную FAR (равную 1) и наименьшую возможную \(ARL_{IC}\) (также равную 1), нецелесообразно использовать карту Shewhart-SN при работе с отдельными измерениями. Далее будет показано, что карта EWMA, основанная на знаковой статистике (обозначаемая EWMA-SN), и карта CUSUM, основанная на знаковой статистике (обозначаемая CUSUM-SN), работают для отдельных измерений лучше.

\begin{figure}[!ht]
	\centering
	\includegraphics[width=0.6\linewidth]{"./images/fig-4.9.png"}
	\caption[]{Карта Shewhart-SN по отдельным измерениям.}
	\label{fig-4.9}
\end{figure}

\noindent\textbf{Пример 4.2. Карта Shewhart-SN по сгруппированным данным.}

Хотя карта Shewhart-SN практически непривлекательна для отдельных данных, как показано выше, её можно использовать для сгруппированных данных умеренного объёма. Столбец (a) таблицы \ref{tab:4.6} показывает наблюдения 15 независимых выборок или подгрупп по \(n = 10\) наблюдений в каждой из неизвестного непрерывного распределения процесса, которые используются для мониторинга известной медианы в состоянии IC \(\mu_0 = 0\). Столбец (b) таблицы \ref{tab:4.6} содержит значения статистики для нанесения на карту \(SN_i\) для \(i = 1, 2, \ldots, 15\). Например, в выборке с номером \(i = 4\) шесть из 10 значений больше нуля, четыре меньше нуля, ни одно не равно нулю. Таким образом, \(SN_4 = \sum_{j=1}^{10} \phi(X_{4j} - 0) = 2\).

Для объёма выборки 10 по записи «правый хвост» для 10 в таблице A Приложения A имеем \(\frac{FAR}{2} = 0.000977\), что меньше 0.00135. Теперь, чтобы получить UCL, то есть \(a\), используем линейное соотношение \(SN_i = 2T_i - n\), так что \(a = 2(10) - 10 = 10\) при \(FAR = (2)(0.000977) = 0.001994\), что меньше 0.0027. Таким образом, для \(n = 10\) имеем \(LCL / UCL = \pm 10\) и \(CL = 0\) при соответствующей \(FAR = 0.001994\) и \(ARL_{IC} = \frac{1}{FAR} = \frac{1}{0.001994} \approx 502\), что хорошо, поскольку близко к типичному желаемому значению 500; типичные промышленные значения включают \(ARL_{IC} = 370\) или 500.

Таким образом, контрольные границы и центральная линия имеют вид
\[
\begin{gathered}
	UCL = +10 \\
	CL = 0 \\
	LCL = -10
\end{gathered}
\]
и соответствующая карта Shewhart-SN показана на рисунке \ref{fig-4.10}.

Поскольку ни одна из статистик для нанесения на карту не попадает на контрольные границы или за них, процесс считается находящимся в состоянии IC относительно целевого значения медианы \(\mu_0 = 0\), и никакой корректирующей меры или поиска назначаемых причин не требуется. Отметим, что эта карта Shewhart-SN применяется без каких-либо предположений о форме лежащего в основе распределения. Например, не нужно предполагать и обосновывать, что данные следуют нормальному распределению, как это необходимо для параметрической карты Шухарта для среднего.

Подчёркивается, что свободная от распределения карта Shewhart-SN --- одна из простейших и наиболее универсальных контрольных карт, которые можно использовать для мониторинга положения любого непрерывного распределения без знания его дисперсии.

\begin{table}[!ht]
	\centering
	\caption{Данные для карты Shewhart-SN по сгруппированным данным и статистики \(SN\).}
	\label{tab:4.6}
	\footnotesize
	\setlength{\tabcolsep}{4pt}%
	\begin{tabular}{@{}crrrrrrrrrrr@{}}\toprule
		Номер & \multicolumn{10}{c}{(a)} & (b) \\
		\(i\) & \(X_1\) & \(X_2\) & \(X_3\) & \(X_4\) & \(X_5\) & \(X_6\) & \(X_7\) & \(X_8\) & \(X_9\) & \(X_{10}\) & \(SN_i\) \\ \midrule
		1 & \(-0.30\) & \(-1.28\) & \(0.24\) & \(1.28\) & \(1.20\) & \(1.73\) & \(-2.18\) & \(-0.23\) & \(1.10\) & \(-1.09\) & \(0\) \\
		2 & \(-0.69\) & \(-1.69\) & \(-1.85\) & \(-0.98\) & \(-0.77\) & \(-2.12\) & \(-0.57\) & \(-0.40\) & \(0.13\) & \(-0.37\) & \(-8\) \\
		3 & \(-0.33\) & \(-0.37\) & \(1.34\) & \(-0.09\) & \(-0.19\) & \(-0.51\) & \(1.97\) & \(0.87\) & \(2.38\) & \(-0.65\) & \(-2\) \\
		4 & \(1.66\) & \(-1.61\) & \(0.54\) & \(0.90\) & \(1.92\) & \(-0.08\) & \(-0.52\) & \(0.68\) & \(-0.38\) & \(0.76\) & \(2\) \\
		5 & \(-1.44\) & \(-0.85\) & \(-1.52\) & \(-0.36\) & \(-0.03\) & \(0.03\) & \(-0.32\) & \(2.19\) & \(-1.74\) & \(-0.74\) & \(-6\) \\
		6 & \(-2.58\) & \(1.45\) & \(-1.28\) & \(-0.65\) & \(0.76\) & \(0.47\) & \(0.87\) & \(0.60\) & \(-1.37\) & \(-1.12\) & \(0\) \\
		7 & \(0.69\) & \(0.32\) & \(-0.94\) & \(-0.24\) & \(0.13\) & \(0.56\) & \(0.14\) & \(-0.91\) & \(1.88\) & \(0.49\) & \(4\) \\
		8 & \(0.07\) & \(0.83\) & \(0.86\) & \(-0.64\) & \(-0.92\) & \(1.11\) & \(-1.20\) & \(-1.56\) & \(0.71\) & \(0.64\) & \(2\) \\
		9 & \(2.21\) & \(1.44\) & \(1.30\) & \(0.11\) & \(0.00\) & \(0.45\) & \(-0.03\) & \(-1.05\) & \(-1.77\) & \(0.83\) & \(3\) \\
		10 & \(0.44\) & \(0.62\) & \(0.21\) & \(-1.03\) & \(1.24\) & \(-0.31\) & \(-0.84\) & \(-0.82\) & \(-0.43\) & \(-0.45\) & \(-2\) \\
		11 & \(-0.52\) & \(0.85\) & \(0.51\) & \(-0.61\) & \(1.30\) & \(-1.76\) & \(0.55\) & \(-0.12\) & \(0.04\) & \(-0.65\) & \(0\) \\
		12 & \(-0.55\) & \(0.85\) & \(0.80\) & \(0.46\) & \(0.69\) & \(1.63\) & \(0.30\) & \(0.59\) & \(1.85\) & \(-0.34\) & \(6\) \\
		13 & \(1.04\) & \(0.14\) & \(1.14\) & \(-0.15\) & \(-0.78\) & \(1.08\) & \(-0.58\) & \(0.53\) & \(0.55\) & \(-0.32\) & \(2\) \\
		14 & \(-0.44\) & \(-1.37\) & \(1.99\) & \(-0.57\) & \(0.09\) & \(-0.23\) & \(2.84\) & \(1.25\) & \(0.88\) & \(1.33\) & \(2\) \\
		15 & \(0.19\) & \(0.54\) & \(-0.25\) & \(-1.22\) & \(1.27\) & \(-0.29\) & \(-1.31\) & \(0.76\) & \(0.78\) & \(0.43\) & \(2\) \\ \bottomrule
	\end{tabular}
\end{table}

\begin{figure}[!ht]
	\centering
	\includegraphics[width=0.7\linewidth]{"./images/fig-4.10.png"}
	\caption[]{Карта Shewhart-SN по сгруппированным данным.}
	\label{fig-4.10}
\end{figure}

\noindent\textbf{Пример 4.3. Карта Shewhart-SN с правилами сигналов на основе серий.}

Далее, как и для параметрической карты \(\overline{X}\) Шухарта, карта Shewhart-SN дополняется улучшенным правилом сигналов на основе серий «2 из 2» и показана здесь для иллюстрации. Смысл этого двоякий. Во-первых, сохраняется простота свободной от распределения карты Shewhart-SN, а во-вторых, эффективность этой простой карты повышается. Эта карта подаёт сигнал, когда либо одна точка нанесена на LCL/UCL или за них, либо две последовательные точки нанесены на LWL/UWL или за них. Как отмечалось ранее, мы рассматриваем только улучшенное правило, поскольку в литературе показано, что оно превосходит стандартное правило.

Для иллюстрации улучшенная карта Shewhart-SN «2 из 2» строится по тем же данным, что и в таблице \ref{tab:4.6}. Если теперь ввести предупреждающие границы, скажем, на уровне \(\pm 6\), то есть верхняя и нижняя предупреждающие границы равны \(UWL = +6\) и \(LWL = -6\) соответственно, то результирующая улучшенная карта Shewhart-SN «2 из 2» показана на рисунке \ref{fig-4.11}.

Следует отметить, что обычно предупреждающие границы находятся следующим образом. Сначала контрольные границы задаются на некотором фиксированном уровне, скажем \(\pm 10\), а затем предупреждающие границы находятся с помощью симуляции, например так, чтобы достигалось некоторое номинальное значение \(ARL_{IC}\).

Поскольку нет двух последовательных точек, нанесённых на LWL/UWL или за них, и нет ни одной точки на LCL/UCL или за ними, процесс считается находящимся в состоянии IC относительно целевого значения медианы \(\mu_0 = 0\), и никакой корректирующей меры или поиска назначаемых причин не требуется.

\begin{figure}[!ht]
	\centering
	\includegraphics[width=0.7\linewidth]{"./images/fig-4.11.png"}
	\caption[]{Карта Shewhart-SN с правилами сигналов на основе серий.}
	\label{fig-4.11}
\end{figure}

Далее рассмотрим свободную от распределения карту Shewhart-SR для мониторинга медианы непрерывного симметричного распределения, основанную на статистике знаковых рангов.

\noindent\textbf{Пример 4.4. Карта Шухарта по знаковым рангам.}

Столбец (a) таблицы \ref{tab:4.7a} содержит наблюдения 20 независимых выборок по \(n = 4\) наблюдения в каждой из неизвестного непрерывного распределения, симметричного относительно известного значения медианы в состоянии IC \(\mu_0\), которое без потери общности принимается равным 0. Таким образом, предполагается, что неизвестное непрерывное распределение симметрично относительно своей медианы 0. Соответствующие абсолютные ранги и знаковые ранги Уилкоксона показаны в столбце (c) таблицы \ref{tab:4.7a} и столбце (b) таблицы \ref{tab:4.7b} соответственно. Статистики для нанесения на карту для карты Shewhart-SR показаны в столбце (c) таблицы \ref{tab:4.7b}.

Если выбрать \(a\) равным наибольшему значению, которое может принимать статистика для нанесения на карту, то есть \(a = \frac{n(n + 1)}{2} = \frac{(4)(5)}{2} = 10\), то контрольные границы и центральная линия имеют вид
\[
\begin{gathered}
	UCL = 10 \\
	CL = 0 \\
	LCL = -10.
\end{gathered}
\]

Для \(n = 4\) и \(a = 10\) по записи «правый хвост» для 10 в таблице B Приложения A имеем \(\frac{FAR}{2} = 0.062\), что является наименьшей вероятностью, связанной с \(n = 4\). Хотя эта вероятность не меньше и не равна 0.00135, это наименьшая вероятность, которую можно получить при работе с выборкой объёма 4. Теперь, чтобы получить UCL, то есть \(a\), используем линейное соотношение \(SR_i = 2T_n^+ - n(n + 1)/2\), так что \(a = 2(10) - 10 = 10\) при \(FAR = (2)(0.062) = 0.124\). Таким образом, для \(n = 4\) имеем \(LCL/UCL = \pm 10\) и \(CL = 0\) при соответствующей \(FAR = 0.124\), которая является наименьшей достижимой FAR. Поскольку значение медианы в состоянии IC известно или задано (случай K), \(ARL_{IC} = \frac{1}{FAR} = \frac{1}{0.124} \approx 8\). Опять же, это низкое значение \(ARL\) в состоянии IC может быть неудовлетворительным с практической точки зрения, поскольку оно далеко от типичного желаемого значения 370. Однако это наибольшая \(ARL_{IC}\), достижимая для этой непараметрической карты при работе с выборкой объёма 4. Результирующая карта Shewhart-SR показана на рисунке \ref{fig-4.12}.

Из рисунка \ref{fig-4.12} видно, что процесс выходит из-под контроля начиная с выборки номер 2, поскольку статистика для нанесения на карту попадает на LCL. Таким образом, необходимы корректирующая мера и поиск назначаемых причин. Следует отметить, что большие значения \(ARL_{IC}\), более близкие к типичным значениям, таким как 370 или 500, достигаются при больших значениях объёма подгруппы \(n\). См., например, таблицы 2a и 2b у Bakir and Reynolds (1979).

\begin{table}[!ht]
	\centering
	\caption{Данные для карты Shewhart-SR.}
	\label{tab:4.7a}
	\small
	\begin{tabular}{@{}ccccccccccccc@{}}\toprule
		Номер & \multicolumn{4}{c}{(a)} & \multicolumn{4}{c}{(b)} & \multicolumn{4}{c}{(c)} \\
		выборки \(i\) & \(x_{i1}\) & \(x_{i2}\) & \(x_{i3}\) & \(x_{i4}\) & \(|x_{i1}|\) & \(|x_{i2}|\) & \(|x_{i3}|\) & \(|x_{i4}|\) & \(R_{i1}\) & \(R_{i2}\) & \(R_{i3}\) & \(R_{i4}\) \\ \midrule
		1 & \(-0.39\) & \(0.03\) & \(0.25\) & \(0.30\) & \(0.39\) & \(0.03\) & \(0.25\) & \(0.30\) & \(4\) & \(1\) & \(2\) & \(3\) \\
		2 & \(-0.61\) & \(-1.32\) & \(-2.84\) & \(-0.07\) & \(0.61\) & \(1.32\) & \(2.84\) & \(0.07\) & \(2\) & \(3\) & \(4\) & \(1\) \\
		3 & \(2.04\) & \(0.40\) & \(0.54\) & \(1.12\) & \(2.04\) & \(0.40\) & \(0.54\) & \(1.12\) & \(4\) & \(1\) & \(2\) & \(3\) \\
		4 & \(-0.51\) & \(-1.63\) & \(0.28\) & \(-0.96\) & \(0.51\) & \(1.63\) & \(0.28\) & \(0.96\) & \(2\) & \(4\) & \(1\) & \(3\) \\
		5 & \(-0.20\) & \(-0.01\) & \(-0.21\) & \(0.82\) & \(0.20\) & \(0.01\) & \(0.21\) & \(0.82\) & \(2\) & \(1\) & \(3\) & \(4\) \\
		6 & \(2.03\) & \(0.70\) & \(-1.21\) & \(-0.26\) & \(2.03\) & \(0.70\) & \(1.21\) & \(0.26\) & \(4\) & \(2\) & \(3\) & \(1\) \\
		7 & \(0.64\) & \(1.48\) & \(0.91\) & \(1.23\) & \(0.64\) & \(1.48\) & \(0.91\) & \(1.23\) & \(1\) & \(4\) & \(2\) & \(3\) \\
		8 & \(0.19\) & \(-2.74\) & \(-0.85\) & \(1.25\) & \(0.19\) & \(2.74\) & \(0.85\) & \(1.25\) & \(1\) & \(4\) & \(2\) & \(3\) \\
		9 & \(-0.51\) & \(-1.77\) & \(-1.69\) & \(0.24\) & \(0.51\) & \(1.77\) & \(1.69\) & \(0.24\) & \(2\) & \(4\) & \(3\) & \(1\) \\
		10 & \(-0.84\) & \(-0.71\) & \(-0.82\) & \(0.19\) & \(0.84\) & \(0.71\) & \(0.82\) & \(0.19\) & \(4\) & \(2\) & \(3\) & \(1\) \\
		11 & \(0.29\) & \(0.47\) & \(0.04\) & \(0.39\) & \(0.29\) & \(0.47\) & \(0.04\) & \(0.39\) & \(2\) & \(4\) & \(1\) & \(3\) \\
		12 & \(0.38\) & \(2.18\) & \(1.24\) & \(-0.80\) & \(0.38\) & \(2.18\) & \(1.24\) & \(0.80\) & \(1\) & \(4\) & \(3\) & \(2\) \\
		13 & \(-0.86\) & \(0.61\) & \(-0.70\) & \(-1.13\) & \(0.86\) & \(0.61\) & \(0.70\) & \(1.13\) & \(3\) & \(1\) & \(2\) & \(4\) \\
		14 & \(-0.31\) & \(1.09\) & \(-0.43\) & \(0.69\) & \(0.31\) & \(1.09\) & \(0.43\) & \(0.69\) & \(1\) & \(4\) & \(2\) & \(3\) \\
		15 & \(0.42\) & \(1.35\) & \(-0.51\) & \(-1.80\) & \(0.42\) & \(1.35\) & \(0.51\) & \(1.80\) & \(1\) & \(3\) & \(2\) & \(4\) \\
		16 & \(0.39\) & \(-0.65\) & \(1.57\) & \(-0.60\) & \(0.39\) & \(0.65\) & \(1.57\) & \(0.60\) & \(1\) & \(3\) & \(4\) & \(2\) \\
		17 & \(1.31\) & \(-0.59\) & \(-0.21\) & \(0.22\) & \(1.31\) & \(0.59\) & \(0.21\) & \(0.22\) & \(4\) & \(3\) & \(1\) & \(2\) \\
		18 & \(0.64\) & \(0.27\) & \(0.26\) & \(0.08\) & \(0.64\) & \(0.27\) & \(0.26\) & \(0.08\) & \(4\) & \(3\) & \(2\) & \(1\) \\
		19 & \(-0.33\) & \(-0.49\) & \(0.38\) & \(-1.54\) & \(0.33\) & \(0.49\) & \(0.38\) & \(1.54\) & \(1\) & \(3\) & \(2\) & \(4\) \\
		20 & \(-0.48\) & \(-0.54\) & \(1.73\) & \(-1.00\) & \(0.48\) & \(0.54\) & \(1.73\) & \(1.00\) & \(1\) & \(2\) & \(4\) & \(3\) \\ \bottomrule
	\end{tabular}
\end{table}

\begin{table}[!ht]
	\centering
	\caption{Расчёты для карты Shewhart-SR.}
	\label{tab:4.7b}
	\footnotesize
	\setlength{\tabcolsep}{4pt}%
	\begin{tabular}{@{}cccccccccc@{}}\toprule
		Шаг \(i\) & \multicolumn{4}{c}{(a)} & \multicolumn{4}{c}{(b)} & (c) \\
		& \(\phi(x_{i1} - \mu_0)\) & \(\phi(x_{i2} - \mu_0)\) & \(\phi(x_{i3} - \mu_0)\) & \(\phi(x_{i4} - \mu_0)\) & \(\phi(x_{i1})R_{i1}\) & \(\phi(x_{i2})R_{i2}\) & \(\phi(x_{i3})R_{i3}\) & \(\phi(x_{i4})R_{i4}\) & \(SR_i\) \\ \midrule
		1 & \(-1\) & \(1\) & \(1\) & \(1\) & \(-4\) & \(1\) & \(2\) & \(3\) & \(2\) \\
		2 & \(-1\) & \(-1\) & \(-1\) & \(-1\) & \(-2\) & \(-3\) & \(-4\) & \(-1\) & \(-10\) \\
		3 & \(1\) & \(1\) & \(1\) & \(1\) & \(4\) & \(1\) & \(2\) & \(3\) & \(10\) \\
		4 & \(-1\) & \(-1\) & \(1\) & \(-1\) & \(-2\) & \(-4\) & \(1\) & \(-3\) & \(-8\) \\
		5 & \(-1\) & \(-1\) & \(-1\) & \(1\) & \(-2\) & \(-1\) & \(-3\) & \(4\) & \(-2\) \\
		6 & \(1\) & \(1\) & \(-1\) & \(-1\) & \(4\) & \(2\) & \(-3\) & \(-1\) & \(2\) \\
		7 & \(1\) & \(1\) & \(1\) & \(1\) & \(1\) & \(4\) & \(2\) & \(3\) & \(10\) \\
		8 & \(1\) & \(-1\) & \(-1\) & \(1\) & \(1\) & \(-4\) & \(-2\) & \(3\) & \(-2\) \\
		9 & \(-1\) & \(-1\) & \(-1\) & \(1\) & \(-2\) & \(-4\) & \(-3\) & \(1\) & \(-8\) \\
		10 & \(-1\) & \(-1\) & \(-1\) & \(1\) & \(-4\) & \(-2\) & \(-3\) & \(1\) & \(-8\) \\
		11 & \(1\) & \(1\) & \(1\) & \(1\) & \(2\) & \(4\) & \(1\) & \(3\) & \(10\) \\
		12 & \(1\) & \(1\) & \(1\) & \(-1\) & \(1\) & \(4\) & \(3\) & \(-2\) & \(6\) \\
		13 & \(-1\) & \(1\) & \(-1\) & \(-1\) & \(-3\) & \(1\) & \(-2\) & \(-4\) & \(-8\) \\
		14 & \(-1\) & \(1\) & \(-1\) & \(1\) & \(-1\) & \(4\) & \(-2\) & \(3\) & \(4\) \\
		15 & \(1\) & \(1\) & \(-1\) & \(-1\) & \(1\) & \(3\) & \(-2\) & \(-4\) & \(-2\) \\
		16 & \(1\) & \(-1\) & \(1\) & \(-1\) & \(1\) & \(-3\) & \(4\) & \(-2\) & \(0\) \\
		17 & \(1\) & \(-1\) & \(-1\) & \(1\) & \(4\) & \(-3\) & \(-1\) & \(2\) & \(2\) \\
		18 & \(1\) & \(1\) & \(1\) & \(1\) & \(4\) & \(3\) & \(2\) & \(1\) & \(10\) \\
		19 & \(-1\) & \(-1\) & \(1\) & \(-1\) & \(-1\) & \(-3\) & \(2\) & \(-4\) & \(-6\) \\
		20 & \(-1\) & \(-1\) & \(1\) & \(-1\) & \(-1\) & \(-2\) & \(4\) & \(-3\) & \(-2\) \\ \bottomrule
	\end{tabular}
\end{table}

\begin{figure}[!ht]
	\centering
	\includegraphics[width=0.7\linewidth]{"./images/fig-4.12.png"}
	\caption[]{Карта Shewhart-SR.}
	\label{fig-4.12}
\end{figure}

\noindent\textbf{Пример 4.5. Карта Shewhart-SR с правилами сигналов на основе серий.}

Для иллюстрации улучшенная карта Shewhart-SR «2 из 2» строится по тем же данным, что и в таблице \ref{tab:4.7a}. Если теперь ввести предупреждающие границы, скажем, \(\pm 6\), то есть верхнюю предупреждающую границу \(UWL = +6\) и нижнюю предупреждающую границу \(LWL = -6\) соответственно, то результирующая улучшенная карта Shewhart-SR «2 из 2» показана на рисунке \ref{fig-4.13}. Следует отметить, что предупреждающие границы находятся следующим образом. Обычно контрольные границы задаются на некотором фиксированном уровне, скажем \(\pm 10\), а затем предупреждающие границы находятся с помощью симуляции так, чтобы достигалось некоторое номинальное значение \(ARL_{IC}\).

Карта Shewhart-SR «2 из 2» подаёт сигнал, когда две из двух последовательных статистик для нанесения на карту попадают на LWL/UWL или за них. Поскольку 9-я и 10-я статистики для нанесения на карту попадают между LWL и LCL, процесс объявляется вышедшим из-под контроля и необходимы поиск назначаемых причин. Кроме того, видно, что 11-я и 12-я статистики для нанесения на карту также дают сигнал о выходе из-под контроля. Больше примеров карты Shewhart-SR с правилами сигналов на основе серий см. в Chakraborti and Eryilmaz (2007).

\begin{figure}[!ht]
	\centering
	\includegraphics[width=0.7\linewidth]{"./images/fig-4.13.png"}
	\caption[]{Карта Shewhart-SR с правилами сигналов на основе серий.}
	\label{fig-4.13}
\end{figure}

\subsection{Контрольные карты CUSUM}

\noindent\textbf{Пример 4.6. Карта CUSUM-SN по отдельным измерениям.}

Для данных таблицы \ref{tab:4.5} вычислены две односторонние статистики для нанесения на карту CUSUM \(SN_t^+\) и \(SN_t^-\) с \(SN_0^+ = SN_0^- = 0\) и \(k = 0.5\); они показаны в таблице \ref{tab:4.8}. Интервал решения выбран равным \(h = 3\). Рекомендации по выбору параметров проектирования читатель найдёт в разделе 3.3.2. Для иллюстрации расчётов рассмотрим наблюдение номер 1. Уравнения для статистик для нанесения на карту \(SN_1^+\) и \(SN_1^-\) имеют вид
\[
	SN_1^+ = \max(0, SN_0^+ + SN_1 - k) = \max(0, 0 + (-1) - 0.5) = 0
\]
и
\[
	SN_1^- = \min(0, SN_0^- + SN_1 + k) = \min(0, 0 + (-1) + 0.5) = -0.5.
\]

\begin{table}[!ht]
	\centering
	\caption{Карта CUSUM-SN по отдельным измерениям.}
	\label{tab:4.8}
	\small
	\begin{tabular}{@{}cccccc@{}}\toprule
		Номер & (a) & (b) & (c) & (d) & (e) \\
		наблюдения \(t\) & \(SN_t\) & \(SN_t^+\) & \(N^+\) & \(SN_t^-\) & \(N^-\) \\ \midrule
		1 & \(-1\) & \(0.0\) & \(0\) & \(-0.5\) & \(1\) \\
		2 & \(-1\) & \(0.0\) & \(0\) & \(-1.0\) & \(2\) \\
		3 & \(-1\) & \(0.0\) & \(0\) & \(-1.5\) & \(3\) \\
		4 & \(1\) & \(0.5\) & \(1\) & \(0.0\) & \(0\) \\
		5 & \(1\) & \(1.0\) & \(2\) & \(0.0\) & \(0\) \\
		6 & \(1\) & \(1.5\) & \(3\) & \(0.0\) & \(0\) \\
		7 & \(-1\) & \(0.0\) & \(0\) & \(-0.5\) & \(1\) \\
		8 & \(1\) & \(0.5\) & \(1\) & \(0.0\) & \(0\) \\
		9 & \(-1\) & \(0.0\) & \(0\) & \(-0.5\) & \(1\) \\
		10 & \(1\) & \(0.5\) & \(1\) & \(0.0\) & \(0\) \\
		11 & \(-1\) & \(0.0\) & \(0\) & \(-0.5\) & \(1\) \\
		12 & \(1\) & \(0.5\) & \(1\) & \(0.0\) & \(0\) \\
		13 & \(1\) & \(1.0\) & \(2\) & \(0.0\) & \(0\) \\
		14 & \(-1\) & \(0.0\) & \(0\) & \(-0.5\) & \(1\) \\
		15 & \(1\) & \(0.5\) & \(1\) & \(0.0\) & \(0\) \\
		16 & \(-1\) & \(0.0\) & \(0\) & \(-0.5\) & \(1\) \\
		17 & \(1\) & \(0.5\) & \(1\) & \(0.0\) & \(0\) \\
		18 & \(1\) & \(1.0\) & \(2\) & \(0.0\) & \(0\) \\
		19 & \(-1\) & \(0.0\) & \(0\) & \(-0.5\) & \(1\) \\
		20 & \(1\) & \(0.5\) & \(1\) & \(0.0\) & \(0\) \\
		21 & \(1\) & \(1.0\) & \(2\) & \(0.0\) & \(0\) \\
		22 & \(-1\) & \(0.0\) & \(0\) & \(-0.5\) & \(1\) \\
		23 & \(1\) & \(0.5\) & \(1\) & \(0.0\) & \(0\) \\
		24 & \(1\) & \(1.0\) & \(2\) & \(0.0\) & \(0\) \\
		25 & \(1\) & \(1.5\) & \(3\) & \(0.0\) & \(0\) \\
		26 & \(1\) & \(2.0\) & \(4\) & \(0.0\) & \(0\) \\
		27 & \(1\) & \(2.5\) & \(5\) & \(0.0\) & \(0\) \\
		28 & \(1\) & \(3.0\) & \(6\) & \(0.0\) & \(0\) \\
		29 & \(1\) & \(3.5\) & \(7\) & \(0.0\) & \(0\) \\
		30 & \(1\) & \(4.0\) & \(8\) & \(0.0\) & \(0\) \\ \bottomrule
	\end{tabular}
\end{table}

\begin{figure}[!ht]
	\centering
	\includegraphics[width=0.7\linewidth]{"./images/fig-4.14.png"}
	\caption[]{Карта CUSUM-SN по отдельным измерениям.}
	\label{fig-4.14}
\end{figure}

Остальные расчёты сведены в столбцы (b) и (d) таблицы \ref{tab:4.8}. Наряду с \(SN_t^+\) и \(SN_t^-\), вычислены и показаны величины \(N^+\), которая считает число последовательных периодов с момента, когда \(SN_t^+\) стала отличной от нуля, и \(N^-\), которая считает число последовательных периодов с момента, когда \(SN_t^-\) стала отличной от нуля. Графическое представление карты CUSUM-SN показано на рисунке \ref{fig-4.14}. Карта подаёт сигнал на наблюдении номер 28 при \(N^+ = 6\), что указывает на то, что сдвиг, скорее всего, произошёл на наблюдении номер 22 (поскольку \(28 - 6 = 22\)). Необходимы поиск назначаемых причин и корректирующая мера.

Следует отметить, что для того же набора данных параметрическая (основанная на нормальной теории) карта CUSUM подаёт сигнал на наблюдении номер 29 при \(N^+ = 7\), что указывает на то, что сдвиг, скорее всего, произошёл на наблюдении номер 22 (поскольку \(29 - 7 = 22\)). Таким образом, параметрический аналог подаёт сигнал на одно наблюдение позже, и обе карты указывают на то, что сдвиг положения, скорее всего, произошёл на наблюдении номер 22.

Пример иллюстрирует ценность непараметрической контрольной карты в целом и непараметрической карты CUSUM-SN в частности для отдельных данных. Если использовать традиционную карту CUSUM, приходится делать предположение о нормальности, что может быть проблемой, особенно в ситуации с отдельными данными, поскольку на центральную предельную теорему здесь рассчитывать нельзя. Более того, известно, что карта CUSUM не очень робастна к предположению о нормальности, что может поставить под сомнение, например, фактическую FAR карты. Исследователи показали, что непараметрическая карта CUSUM-SN может превосходить карту CUSUM-SN нормальной теории в ситуациях, когда лежащее в основе распределение имеет более тяжёлые хвосты, чем нормальное. Расчёты для непараметрической карты несложны. В этом примере непараметрическая карта подала сигнал на один период раньше параметрической. Короче говоря, непараметрическая карта CUSUM-SN — более безопасная и работоспособная альтернативная контрольная карта для практического применения.

Далее иллюстрируем карту CUSUM-SN для сгруппированных данных.

\noindent\textbf{Пример 4.7. Карта CUSUM-SN по сгруппированным данным.}

Для данных таблицы \ref{tab:4.6} при \(n = 10\) вычислены две односторонние статистики для нанесения на карту CUSUM \(SN_t^+\) и \(SN_t^-\), с \(SN_0^+ = SN_0^- = 0\) и \(k = 2\); они показаны в таблице \ref{tab:4.9}. Напомним, что референсное значение \(k\) выбирается первым и обычно является небольшим. После этого интервал решения \(h\) находится так, чтобы достигалось номинальное значение \(ARL_{IC}\). Читатель может обратиться к Amin, Reynolds и Bakir (1995), таблица 5, где приведено несколько комбинаций \((h, k)\) для различных номинальных значений \(ARL_{IC}\) при объёме выборки \(n = 10\). В этом примере выбрано \(h = 8\). Для иллюстрации расчётов рассмотрим выборку номер 1. Уравнения для статистик для нанесения на карту \(SN_1^+\) и \(SN_1^-\) имеют вид
\[
	SN_1^+ = \max(0, SN_0^+ + SN_1 - k) = \max(0, 0 + 0 - 2) = 0
\]
и
\[
	SN_1^- = \min(0, SN_0^- + SN_1 + k) = \min(0, 0 + 0 + 2) = 0.
\]

\begin{table}[!ht]
	\centering
	\caption{Карта CUSUM-SN по сгруппированным данным.}
	\label{tab:4.9}
	\begin{tabular}{@{}cccccc@{}}\toprule
		Выборка \(t\) & (a) \(SN_t\) & (b) \(SN_t^+\) & (c) \(N^+\) & (d) \(SN_t^-\) & (e) \(N^-\) \\ \midrule
		1 & \(0\) & \(0\) & \(0\) & \(0\) & \(0\) \\
		2 & \(-8\) & \(0\) & \(0\) & \(-6\) & \(1\) \\
		3 & \(-2\) & \(0\) & \(0\) & \(-6\) & \(2\) \\
		4 & \(2\) & \(0\) & \(0\) & \(-2\) & \(3\) \\
		5 & \(-6\) & \(0\) & \(0\) & \(-6\) & \(4\) \\
		6 & \(0\) & \(0\) & \(0\) & \(-4\) & \(5\) \\
		7 & \(4\) & \(2\) & \(1\) & \(0\) & \(0\) \\
		8 & \(2\) & \(2\) & \(2\) & \(0\) & \(0\) \\
		9 & \(3\) & \(3\) & \(3\) & \(0\) & \(0\) \\
		10 & \(-2\) & \(0\) & \(0\) & \(0\) & \(0\) \\
		11 & \(0\) & \(0\) & \(0\) & \(0\) & \(0\) \\
		12 & \(6\) & \(4\) & \(1\) & \(0\) & \(0\) \\
		13 & \(2\) & \(4\) & \(2\) & \(0\) & \(0\) \\
		14 & \(2\) & \(4\) & \(3\) & \(0\) & \(0\) \\
		15 & \(2\) & \(4\) & \(4\) & \(0\) & \(0\) \\ \bottomrule
	\end{tabular}
\end{table}

\begin{figure}[!ht]
	\centering
	\includegraphics[width=0.7\linewidth]{"./images/fig-4.15.png"}
	\caption[]{Карта CUSUM-SN по сгруппированным данным.}
	\label{fig-4.15}
\end{figure}

Остальные расчёты сведены в столбцы (b) и (d) таблицы \ref{tab:4.9}. Наряду с \(SN_t^+\) и \(SN_t^-\), вычислены и показаны величины \(N^+\), которая считает число последовательных периодов с момента, когда \(SN_t^+\) стала отличной от нуля, и \(N^-\), которая считает число последовательных периодов с момента, когда \(SN_t^-\) стала отличной от нуля. Графическое представление карты CUSUM-SN показано на рисунке \ref{fig-4.15}. Эта карта не подаёт сигнал и не обнаруживает никаких очевидных закономерностей, поэтому ни поиск назначаемых причин, ни корректирующая мера не требуются. Следует отметить, что для того же набора данных параметрическая карта CUSUM также не подаёт сигнал.

\noindent\textbf{Пример 4.8. Карта CUSUM-SN: подход на основе марковских цепей.}

Рассмотрим верхнюю одностороннюю карту CUSUM-SN с интервалом решения 4 (\(h = 4\)), референсным значением 1 (\(k = 1\)) и объёмом выборки 5 (\(n = 5\)). Относительно выбора комбинации \((h, k) = (4, 1)\) напомним, что референсное значение \(k\) выбирается первым и обычно является небольшим. После этого интервал решения \(h\) находится так, чтобы достигалось номинальное значение \(ARL_{IC}\). Отметим, что Amin, Reynolds и Bakir (1995) в таблице 5 приводят комбинации \((h, k)\) только для объёма выборки \(n = 10\). Поэтому для нахождения комбинации \((h, k)\) при \(n = 5\), дающей достигнутую \(ARL_{IC}\), подходящую для применений SPC, использовался алгоритм симуляции, например на уровне промышленных стандартов 370 или 500. Пошаговые детали симуляционного подхода для карты CUSUM-SN см. в разделе 4.2.2.3. Мы иллюстрируем вычисление характеристик распределения длины серии в состоянии IC этой карты с помощью подхода MC. Для нечётного \(n\) референсное значение берётся нечётным, поскольку тогда сумма \((SN_i - k)\) равна чётным значениям, что уменьшает размер пространства состояний для MC. Это вдвое уменьшит размер матрицы с переходными вероятностями. Для \(h = 4\) пространство состояний равно \(\Omega^+ = \{\varsigma_0, \varsigma_1, \varsigma_2\} = \{0, 2, 4\}\), где \(0 = \varsigma_0 < \varsigma_1 < \varsigma_2 = h\). Пространство состояний вычисляется по уравнению \eqref{eq:4.9}; расчёты приведены в таблице \ref{tab:4.10}.

\begin{table}[!ht]
	\centering
	\caption{Вычисление пространства состояний при \(h = 4\), \(k = 1\) и \(n = 5\).}
	\label{tab:4.10}
	\small
	\begin{tabular}{@{}cccc@{}}\toprule
		\(SN_t\) & \(S_{t-1}^+ + SN_t - k\) & \(\max\{0, S_{t-1}^+ + SN_t - k\}\) & \(S_t^+ = \min\{h, \max\{0, S_{t-1}^+ + SN_t - k\}\}\) \\ \midrule
		\(-5\) & \(-6^{a}\) & \(0\) & \(0\) \\
		\(-3\) & \(-4\) & \(0\) & \(0\) \\
		\(-1\) & \(-2\) & \(0\) & \(0\) \\
		\(1\) & \(0\) & \(0\) & \(0\) \\
		\(3\) & \(2\) & \(2\) & \(2\) \\
		\(5\) & \(4\) & \(4\) & \(4\) \\ \bottomrule
	\end{tabular}
	\par\smallskip
	\footnotesize \(^{a}\) Поскольку нужно описать только пространство состояний, \(S_{t-1}^+\) может быть любым значением из \(\Omega^+\), и поэтому без потери общности полагаем \(S_{t-1}^+ = 0\). Любое другое возможное значение \(S_{t-1}^+\) привело бы к тому же \(\Omega^+\).
\end{table}

\begin{table}[!ht]
	\centering
	\caption{Классификация и описание состояний.}
	\label{tab:4.11}
	\begin{tabular}{@{}ccc@{}}\toprule
		Номер состояния & Описание состояния & Поглощающее (A) / непоглощающее (NA) \\ \midrule
		\(0\) & \(S_t^+ = 0\) & NA \\
		\(1\) & \(S_t^+ = 2\) & NA \\
		\(2\) & \(S_t^+ = 4\) & A \\ \bottomrule
	\end{tabular}
\end{table}

Из таблицы \ref{tab:4.11} видно, что есть два непоглощающих состояния, то есть \(v = 2\), и одно поглощающее. Поэтому соответствующая матрица переходных вероятностей имеет размер \(3 \times 3\). Используя два состояния в состоянии IC, матрица переходных вероятностей в состоянии IC имеет вид
\begin{align*}
	M_{3\times3}
	&= \left(\begin{array}{c|c} Q_{2\times2} & p_{2\times1} \\ \hline 0_{1\times2} & 1_{1\times1} \end{array}\right)
	= \left(\begin{array}{cc|c} p_{00} & p_{01} & p_{02} \\ p_{10} & p_{11} & p_{12} \\ \hline p_{20} & p_{21} & p_{22} \end{array}\right) \\
	&= \frac{1}{32}\left(\begin{array}{cc|c} 26 & 5 & 1 \\ 16 & 10 & 6 \\ \hline 0 & 0 & 32 \end{array}\right),
\end{align*}
где переходные вероятности вычислены по уравнению \eqref{eq:4.9}. Для иллюстрации ниже показано вычисление \(p_{01}\), когда процесс находится в состоянии IC:
\begin{align*}
	p_{01} &= P(S_t = 2 \mid S_{t-1} = 0) \\
	      &= P(\min\{4, \max\{0, 0 + SN_t - 1\}\} = 2) \\
	      &= P(\max\{0, SN_t - 1\} = 2) \\
	      &= P(SN_t - 1 = 2) \\
	      &= P(SN_t = 3) \\
	      &= P(2T_t - 5 = 3) \\
	      &= P(T_t = 4) \\
	      &= \frac{5}{32}.
\end{align*}
Предпоследнее равенство следует из того, что случайная величина \(T_t\), имеющая в состоянии IC распределение \(BIN(n = 5, p = 0.5)\), принимает только неотрицательные целые значения (\(T_t \in \{0, 5\}\)).

Как только одношаговые переходные вероятности вычислены и найдена существенная матрица переходных вероятностей \(Q\), характеристики длины серии в состоянии IC получаются по уравнениям 3.22--3.24. Это даёт \(ARL_{IC} = 16.62\) и \(SDRL_{IC} = 15.51\) при значениях 5-го, 25-го, 50-го, 75-го и 95-го процентилей в состоянии IC, равных 2, 6, 12, 23 и 48 соответственно.

Значение \(ARL_{IC}\) означает, что, когда процесс находится в состоянии IC, первый сигнал будет наблюдаться в среднем примерно на каждой 17-й нанесённой точке этой карты. Первый квартиль равен 6, поэтому мы знаем, что ложная тревога не возникнет на первых шести нанесённых точках с вероятностью не более 75\%. Другой пример: \(MRL_{IC} = 12\), что означает, что первая ложная тревога будет наблюдаться в пределах первых 12 нанесённых точек как минимум в 50\% случаев.

Отметим, что, хотя здесь иллюстрируются характеристики в состоянии IC, характеристики в состоянии OOC легко получить, изменив значение \(p\) так, чтобы \(p \ne 0.5\), вычислив значения \(p_{ij}\), образовав существенную матрицу переходных вероятностей \(Q\) и затем используя те же формулы MC. Детали мы оставляем читателю для проверки.

\noindent\textbf{Пример 4.9. Карта CUSUM-SR.}

Для данных таблицы \ref{tab:4.7a} вычислены две односторонние статистики для нанесения на карту CUSUM \(SR_t^+\) и \(SR_t^-\), с \(SR_0^+ = SR_0^- = 0\) и \(k = 3\), а также два счётчика \(N^+\) и \(N^-\); они показаны в таблице \ref{tab:4.12}. Интервал решения выбран равным \(h = 8\).

Для иллюстрации расчётов рассмотрим выборку номер 1. Выражения для статистик для нанесения на карту CUSUM \(SR_1^+\) и \(SR_1^-\) имеют вид
\begin{align*}
	SR_1^+ &= \max(0, SR_0^+ + SR_1 - k) \\
	       &= \max(0, 0 + 2 - 3) \\
	       &= 0
\end{align*}
и
\begin{align*}
	SR_1^- &= \min(0, SR_0^- + SR_1 + k) \\
	       &= \min(0, 0 + 2 + 3) \\
	       &= 0.
\end{align*}

\begin{table}[!ht]
	\centering
	\caption{Карта CUSUM-SR.}
	\label{tab:4.12}
	\begin{tabular}{@{}ccccc@{}}\toprule
		\multicolumn{1}{c}{Выборка \(i\)} & \multicolumn{2}{c}{(a)} & \multicolumn{2}{c}{(b)} \\ \cmidrule(lr){2-3}\cmidrule(lr){4-5}
		& \(SR_t^+\) & \(N^+\) & \(SR_t^-\) & \(N^-\) \\ \midrule
		1 & \(0\) & \(0\) & \(0\) & \(0\) \\
		2 & \(0\) & \(0\) & \(-7\) & \(1\) \\
		3 & \(7\) & \(1\) & \(0\) & \(0\) \\
		4 & \(0\) & \(0\) & \(-5\) & \(1\) \\
		5 & \(0\) & \(0\) & \(-4\) & \(2\) \\
		6 & \(0\) & \(0\) & \(0\) & \(0\) \\
		7 & \(7\) & \(1\) & \(0\) & \(0\) \\
		8 & \(2\) & \(2\) & \(0\) & \(0\) \\
		9 & \(0\) & \(0\) & \(-5\) & \(1\) \\
		10 & \(0\) & \(0\) & \(-10^{a}\) & \(2\) \\
		11 & \(7\) & \(1\) & \(0\) & \(0\) \\
		12 & \(10^{a}\) & \(2\) & \(0\) & \(0\) \\
		13 & \(0\) & \(0\) & \(-5\) & \(1\) \\
		14 & \(1\) & \(1\) & \(0\) & \(0\) \\
		15 & \(0\) & \(0\) & \(0\) & \(0\) \\
		16 & \(0\) & \(0\) & \(0\) & \(0\) \\
		17 & \(0\) & \(0\) & \(0\) & \(0\) \\
		18 & \(7\) & \(1\) & \(0\) & \(0\) \\
		19 & \(0\) & \(0\) & \(-3\) & \(1\) \\
		20 & \(0\) & \(0\) & \(-2\) & \(2\) \\ \bottomrule
	\end{tabular}
	\par\smallskip
	\footnotesize \(^{a}\) Двусторонняя карта CUSUM-SR подаёт сигнал.
\end{table}

\begin{figure}[!ht]
	\centering
	\includegraphics[width=0.7\linewidth]{"./images/fig-4.16.png"}
	\caption[]{Карта CUSUM-SR для сгруппированных данных.}
	\label{fig-4.16}
\end{figure}

Остальные расчёты сведены в столбцы (a) и (b) таблицы \ref{tab:4.12}. Карта CUSUM-SR, показанная на рисунке \ref{fig-4.16}, подаёт сигнал дважды. Первый сигнал даётся на выборке номер 10, когда \(SR_{10}^- = -10\) попадает ниже нижнего интервала решения, то есть \(-h = -8\), при счётчике \(N^- = 2\), что указывает на то, что сдвиг, скорее всего, произошёл на выборке номер 8. Второй сигнал даётся вскоре после этого на выборке номер 12, когда \(SR_{12}^+ = 10\) попадает выше верхнего интервала решения, то есть \(h = 8\), при счётчике \(N^+ = 2\), что снова указывает на то, что сдвиг, скорее всего, произошёл на выборке номер 8.

Если последующие исследования выявят какую-либо назначаемую причину для первого сигнала, нижнюю одностороннюю статистику для нанесения на карту CUSUM следовало бы перезапустить с нуля, тогда как верхние односторонние статистики для нанесения на карту CUSUM продолжали бы работать «как обычно». Обратное следовало бы сделать, если бы назначаемые причины были найдены для второго сигнала. Однако в рассматриваемой ситуации, а именно сигнал на нижней стороне (указывающий на возможный нисходящий сдвиг) на выборке номер 10, а вскоре после этого сигнал на верхней стороне (указывающий на возможный восходящий сдвиг) на выборке номер 12, эти два сигнала были сочтены ложными тревогами, и ни одна из статистик для нанесения на карту CUSUM не была сброшена в ноль. В конце концов, хотя это и не является невозможным, маловероятно, что за нисходящим сдвигом процесса «так скоро» последует восходящий, если только не было «перерегулирования» в положительном направлении со стороны кого-либо из персонала, работающего с контрольной картой.

Следует отметить, что для того же набора данных параметрическая карта CUSUM подаёт сигнал на нижней CUSUM на выборке номер 2 при счётчике \(N^+ = 1\), что указывает на то, что сдвиг, скорее всего, произошёл в самом начале. Таким образом, карта CUSUM-X превосходит карту CUSUM-SR, что согласуется с тем фактом, что для этих данных критерий согласия для нормальности не отвергается. Однако на практике нормальность обычно может вызывать сомнения или не быть обоснованной из-за недостатка информации или данных, и непараметрический метод может оказаться более предпочтительным.

Далее иллюстрируем вычисление некоторых характеристик распределения длины серии верхней односторонней карты CUSUM-SR.

\noindent\textbf{Пример 4.10. Карта CUSUM-SR: подход на основе марковских цепей.}

Рассмотрим верхнюю одностороннюю карту CUSUM-SR с интервалом решения 6 (\(h = 6\)), референсным значением 2 (\(k = 2\)) и объёмом выборки 4 (\(n = 4\)). Опять же, относительно выбора комбинации \((h, k) = (6, 2)\) напомним, что референсное значение \(k\) выбирается первым и обычно является небольшим. После этого интервал решения \(h\) находится так, чтобы достигалось номинальное значение \(ARL_{IC}\). Читатель может обратиться к Bakir and Reynolds (1979), где приведено несколько комбинаций \((h, k)\) для различных номинальных значений \(ARL_{IC}\) при объёме выборки \(n = 4, 6\) и 10. Как и для карты CUSUM-SN, мы иллюстрируем вычисление характеристик распределения длины серии в состоянии IC карты CUSUM-SR с помощью подхода MC.

Для \(h = 6\) имеем \(\Omega^+ = \{\varsigma_0, \varsigma_1, \varsigma_2, \varsigma_3\} = \{0, 2, 4, 6\}\), где \(0 = \varsigma_0 < \varsigma_1 < \varsigma_2 < \varsigma_3 = h\). Пространство состояний вычисляется по уравнению \eqref{eq:4.12}; расчёты приведены в таблице \ref{tab:4.13}.

\begin{table}[!ht]
	\centering
	\caption{Вычисление пространства состояний при \(h = 6\), \(k = 2\) и \(n = 4\).}
	\label{tab:4.13}
	\small
	\begin{tabular}{@{}cccc@{}}\toprule
		\(SR_t\) & \(S_{t-1}^+ + SR_t - k\) & \(\max\{0, S_{t-1}^+ + SR_t - k\}\) & \(S_t^+ = \min\{h, \max\{0, S_{t-1}^+ + SR_t - k\}\}\) \\ \midrule
		\(-10\) & \(-12^{a}\) & \(0\) & \(0\) \\
		\(-8\) & \(-10\) & \(0\) & \(0\) \\
		\(-6\) & \(-8\) & \(0\) & \(0\) \\
		\(-4\) & \(-6\) & \(0\) & \(0\) \\
		\(-2\) & \(-4\) & \(0\) & \(0\) \\
		\(0\) & \(-2\) & \(0\) & \(0\) \\
		\(2\) & \(0\) & \(0\) & \(0\) \\
		\(4\) & \(2\) & \(2\) & \(2\) \\
		\(6\) & \(4\) & \(4\) & \(4\) \\
		\(8\) & \(6\) & \(6\) & \(6\) \\
		\(10\) & \(8\) & \(8\) & \(6\) \\ \bottomrule
	\end{tabular}
	\par\smallskip
	\footnotesize \(^{a}\) Поскольку нужно описать только пространство состояний, \(S_{t-1}^+\) может быть любым значением из \(\Omega^+\), и поэтому без потери общности полагаем \(S_{t-1}^+ = 0\). Любое другое возможное значение \(S_{t-1}^+\) привело бы к тому же \(\Omega^+\).
\end{table}

\begin{table}[!ht]
	\centering
	\caption{Классификация и описание состояний.}
	\label{tab:4.14}
	\begin{tabular}{@{}ccc@{}}\toprule
		Номер состояния & Описание состояния & Поглощающее (A) / непоглощающее (NA) \\ \midrule
		\(0\) & \(S_t^+ = 0\) & NA \\
		\(1\) & \(S_t^+ = 2\) & NA \\
		\(2\) & \(S_t^+ = 4\) & NA \\
		\(3\) & \(S_t^+ = 6\) & A \\ \bottomrule
	\end{tabular}
\end{table}

Из таблицы \ref{tab:4.14} видно, что есть три непоглощающих состояния, то есть \(v = 3\), и одно поглощающее. Поэтому соответствующая матрица переходных вероятностей будет иметь размер \(4 \times 4\). Используя три состояния в состоянии IC, матрица переходных вероятностей в состоянии IC имеет вид
\begin{align*}
	M_{4\times4}
	&= \left(\begin{array}{c|c} Q_{3\times3} & p_{3\times1} \\ \hline 0_{1\times3} & 1_{1\times1} \end{array}\right)
	= \left(\begin{array}{ccc|c} p_{00} & p_{02} & p_{04} & p_{06} \\ p_{20} & p_{22} & p_{24} & p_{26} \\ p_{40} & p_{42} & p_{44} & p_{46} \\ \hline p_{60} & p_{62} & p_{64} & p_{66} \end{array}\right) \\
	&= \frac{1}{16}\left(\begin{array}{ccc|c} 11 & 2 & 1 & 2 \\ 9 & 2 & 2 & 3 \\ 7 & 2 & 2 & 5 \\ \hline 0 & 0 & 0 & 16 \end{array}\right),
\end{align*}
где переходные вероятности вычислены по уравнению \eqref{eq:4.12}. Для иллюстрации ниже показано вычисление \(p_{02}\), когда процесс находится в состоянии IC:
\begin{align*}
	p_{02} &= P(S_t = 2 \mid S_{t-1} = 0) \\
	       &= P(\min\{6, \max\{0, 0 + SR_t - 2\}\} = 2) \\
	       &= P(\max\{0, SR_t - 2\} = 2) \\
	       &= P(SR_t - 2 = 2) \\
	       &= P(SR_t = 4) \\
	       &= P(2T_n^+ - 10 = 4) \\
	       &= P(T_n^+ = 7) \\
	       &= \frac{2}{16}.
\end{align*}
Предпоследнее равенство следует из распределения вероятностей в состоянии IC известной статистики знаковых рангов Уилкоксона \(T_n^+\), подробности которого приведены в таблице \ref{tab:4.3} раздела 4.2.1.2.

Как только одношаговые переходные вероятности вычислены и найдена существенная матрица переходных вероятностей \(Q\), характеристики длины серии в состоянии IC получаются по уравнениям 3.22--3.24. Это даёт \(ARL_{IC} = 6.81\) и \(SDRL_{IC} = 6.11\) при значениях 5-го, 25-го, 50-го, 75-го и 95-го процентилей в состоянии IC, равных 1, 2, 5, 9 и 19 соответственно. Эти значения помогают понять эффективность контрольной карты через её распределение длины серии; например, значение \(ARL_{IC}\) означает, что, когда процесс находится в состоянии IC, первый сигнал будет наблюдаться в среднем на каждой седьмой нанесённой точке. Первый квартиль равен 2, поэтому мы знаем, что ложная тревога возникнет в пределах первых двух нанесённых точек с вероятностью не менее 75\%. Другой пример: \(MRL_{IC} = 5\), что означает, что первая ложная тревога будет наблюдаться в пределах первых пяти нанесённых точек как минимум в 50\% случаев.

Опять же, следует отметить, что, хотя здесь показаны только характеристики распределения длины серии в состоянии IC, можно выполнить аналогичные шаги для получения характеристик длины серии в состоянии OOC, для чего нужно вычислить распределение в состоянии OOC и переходные вероятности \(p_{ij}\) для статистики знаковых рангов Уилкоксона \(T_n^+\).

Далее иллюстрируем карту EWMA-SN для отдельных данных.

\noindent\textbf{Пример 4.11. Карта EWMA-SN по отдельным измерениям.}

Непараметрическая двусторонняя карта CUSUM-SN для данных таблицы \ref{tab:4.8} была показана в примере 4.6. В качестве альтернативы теперь иллюстрируется непараметрическая контрольная карта EWMA. Для этого используем \(\lambda = 0.10\) для вычисления статистик для нанесения на карту и \(L = 2.820\) для нахождения контрольных границ. Подробное обсуждение выбора параметров проектирования \(\lambda\) и \(L\) см. в разделе 3.2.3. Коротко рекомендация состоит в том, чтобы выбирать малое \(\lambda\), скажем равное 0.05, когда представляют интерес малые сдвиги; выбирать \(\lambda = 0.10\), если речь идёт об умеренных сдвигах; и выбирать \(\lambda = 0.20\), если интересны большие сдвиги. Таким образом, выбор \(\lambda = 0.10\) означает, что представляют интерес умеренные сдвиги. После того как \(\lambda\) выбрана, второй шаг заключается в нахождении \(L\) так, чтобы достигалась желаемая \(ARL_{IC}\). Для отдельных измерений и \(\lambda = 0.10\) показано, что \(L = 2.820\) даёт достигнутую \(ARL_{IC}\), близкую к 500 (см. Graham, Human and Chakraborti, 2009)\footnote{Graham, M.~A., Human, S.~W. и Chakraborti, S. (2009). A nonparametric EWMA control chart based on the sign statistic. Technical report 09/04, Department of Statistics, University of Pretoria, Pretoria, South Africa, ISBN 978-1-86854-777-7. Печатный экземпляр этого технического отчёта доступен в University of Pretoria (Претория, ЮАР). Электронную копию можно запросить, написав автору по адресу marien.graham@up.ac.za.}.

Для иллюстрации расчётов рассмотрим наблюдение номер 1 с \(SN_1 = -1\) и предположим, что начальное значение \(Z_0 = 0\). Тогда первая статистика для нанесения на карту на карте EWMA-SN (рисунок \ref{fig-4.17}) равна
\begin{align*}
	Z_1 &= \lambda SN_1 + (1 - \lambda)Z_0 \\
	    &= 0.10(-1) + (1 - 0.10)0 = -0.1.
\end{align*}
Значение второй статистики для нанесения на карту равно
\begin{align*}
	Z_2 &= \lambda SN_2 + (1 - \lambda)Z_1 \\
	    &= 0.10(-1) + (1 - 0.10)(-0.1) \\
	    &= -0.19,
\end{align*}
а остальные расчёты сведены в столбец (b) таблицы \ref{tab:4.15}.

Для простоты используем стационарные контрольные границы на карте EWMA-SN, которые находятся по уравнению \eqref{eq:4.15}, то есть
\begin{align*}
	LCL &= -L\sqrt{\frac{\lambda}{2 - \lambda}} \\
	    &= -2.820\sqrt{\frac{0.10}{2 - 0.10}} \\
	    &= -0.647
\end{align*}
и
\begin{align*}
	UCL &= L\sqrt{\frac{\lambda}{2 - \lambda}} \\
	    &= 2.820\sqrt{\frac{0.10}{2 - 0.10}} \\
	    &= 0.647
\end{align*}
соответственно.

\begin{table}[!ht]
	\centering
	\caption{Расчёты для карты EWMA-SN по отдельным измерениям.}
	\label{tab:4.15}
	\begin{tabular}{@{}ccc@{}}\toprule
		Выборка \(t\) & (a) \(SN_t\) & (b) \(Z_t = \lambda SN_t + (1 - \lambda)Z_{t-1}\) \\ \midrule
		1 & \(-1\) & \(-0.100\) \\
		2 & \(-1\) & \(-0.190\) \\
		3 & \(-1\) & \(-0.271\) \\
		4 & \(1\) & \(-0.144\) \\
		5 & \(1\) & \(-0.030\) \\
		6 & \(1\) & \(0.073\) \\
		7 & \(-1\) & \(-0.034\) \\
		8 & \(1\) & \(0.069\) \\
		9 & \(-1\) & \(-0.037\) \\
		10 & \(1\) & \(0.066\) \\
		11 & \(-1\) & \(-0.040\) \\
		12 & \(1\) & \(0.064\) \\
		13 & \(1\) & \(0.157\) \\
		14 & \(-1\) & \(0.042\) \\
		15 & \(1\) & \(0.137\) \\
		16 & \(-1\) & \(0.024\) \\
		17 & \(1\) & \(0.121\) \\
		18 & \(1\) & \(0.209\) \\
		19 & \(-1\) & \(0.088\) \\
		20 & \(1\) & \(0.179\) \\
		21 & \(1\) & \(0.261\) \\
		22 & \(-1\) & \(0.135\) \\
		23 & \(1\) & \(0.222\) \\
		24 & \(1\) & \(0.300\) \\
		25 & \(1\) & \(0.370\) \\
		26 & \(1\) & \(0.433\) \\
		27 & \(1\) & \(0.489\) \\
		28 & \(1\) & \(0.540\) \\
		29 & \(1\) & \(0.586\) \\
		30 & \(1\) & \(0.628\) \\ \bottomrule
	\end{tabular}
\end{table}

\begin{figure}[!ht]
	\centering
	\includegraphics[width=0.7\linewidth]{"./images/fig-4.17.png"}
	\caption[]{Карта EWMA-SN по отдельным измерениям.}
	\label{fig-4.17}
\end{figure}

Результирующая карта EWMA-SN показана на рисунке \ref{fig-4.17}. Поскольку ни одна из статистик для нанесения на карту не попадает на LCL или UCL либо за них, процесс считается находящимся в состоянии IC, и последующих (корректирующих) мер не требуется. Однако характер изменения статистик для нанесения на карту вызывает некоторое беспокойство по поводу того, что положение может увеличиваться.

Следует отметить, что для того же набора данных карта CUSUM-SN подала сигнал на наблюдении номер 28 при \(N^+ = 6\), что указывает на то, что сдвиг, скорее всего, произошёл на наблюдении номер 22. Параметрическая карта CUSUM, однако, подала сигнал на наблюдении номер 29 при \(N^+ = 7\), что указывает на то, что сдвиг, скорее всего, произошёл на наблюдении номер 22. Хотя карта EWMA-SN не подаёт сигнал, начиная с наблюдения номер 22 виден явный восходящий тренд, и карта почти подаёт сигнал на 30-м (последнем) наблюдении.

\subsection{Контрольные карты EWMA}

\noindent\textbf{Пример 4.12. Карта EWMA-SN по сгруппированным данным.}

Двусторонняя карта CUSUM-SN для данных таблицы \ref{tab:4.9} была показана в примере 4.7. В качестве альтернативы теперь иллюстрируется контрольная карта EWMA. Для этого используем \(\lambda = 0.05\) для вычисления статистик для нанесения на карту и \(L = 2.612\) для нахождения контрольных границ (см. Graham, Chakraborti and Human, 2011a). Подробное обсуждение выбора параметров проектирования \(\lambda\) и \(L\) см. в разделе 3.2.3. Для выборок объёма 10 и \(\lambda = 0.05\) показано, что \(L = 2.612\) даёт достигнутое значение \(ARL_{IC}\), близкое к 500.

Для иллюстрации расчётов рассмотрим выборку номер 1 с \(SN_1 = 0\) и предположим, что начальное значение \(Z_0 = 0\). Тогда первая статистика для нанесения на карту на карте EWMA-SN (рисунок \ref{fig-4.18}) равна
\begin{align*}
	Z_1 &= \lambda SN_1 + (1 - \lambda)Z_0 \\
	    &= 0.05(0) + (1 - 0.05)0 \\
	    &= 0.
\end{align*}
Значение второй статистики для нанесения на карту равно
\begin{align*}
	Z_2 &= \lambda SN_2 + (1 - \lambda)Z_1 \\
	    &= 0.05(-8) + (1 - 0.05)0 \\
	    &= -0.4,
\end{align*}
а остальные расчёты сведены в столбец (b) таблицы \ref{tab:4.16}.

\begin{table}[!ht]
	\centering
	\caption{Расчёты для карты EWMA-SN по сгруппированным данным.}
	\label{tab:4.16}
	\begin{tabular}{@{}ccc@{}}\toprule
		Выборка \(t\) & (a) \(SN_t\) & (b) \(Z_t = \lambda SN_t + (1 - \lambda)Z_{t-1}\) \\ \midrule
		1 & \(0\) & \(0.000\) \\
		2 & \(-8\) & \(-0.400\) \\
		3 & \(-2\) & \(-0.480\) \\
		4 & \(2\) & \(-0.356\) \\
		5 & \(-6\) & \(-0.638\) \\
		6 & \(0\) & \(-0.606\) \\
		7 & \(4\) & \(-0.376\) \\
		8 & \(2\) & \(-0.257\) \\
		9 & \(3\) & \(-0.094\) \\
		10 & \(-2\) & \(-0.190\) \\
		11 & \(0\) & \(-0.180\) \\
		12 & \(6\) & \(0.129\) \\
		13 & \(2\) & \(0.222\) \\
		14 & \(2\) & \(0.311\) \\
		15 & \(2\) & \(0.396\) \\ \bottomrule
	\end{tabular}
\end{table}

Для простоты используем стационарные контрольные границы на карте EWMA-SN, которые находятся по уравнению \eqref{eq:4.17}:
\[
	LCL = -L\sqrt{\frac{\lambda n}{2 - \lambda}} = -2.612\sqrt{\frac{(0.05)(10)}{2 - 0.05}} = -1.323
\]
и
\[
	UCL = L\sqrt{\frac{\lambda n}{2 - \lambda}} = 2.612\sqrt{\frac{(0.05)(10)}{2 - 0.05}} = 1.323.
\]

\begin{figure}[!ht]
	\centering
	\includegraphics[width=0.7\linewidth]{"./images/fig-4.18.png"}
	\caption[]{Карта EWMA-SN по сгруппированным данным.}
	\label{fig-4.18}
\end{figure}

Результирующая карта EWMA-SN показана на рисунке \ref{fig-4.18}. Поскольку ни одна из статистик для нанесения на карту не попадает на LCL и/или UCL или за них и не обнаруживается никакой очевидной вызывающей беспокойство закономерности, процесс считается находящимся в состоянии IC, и корректирующих мер не требуется.

Следует отметить, что для того же набора данных параметрическая карта CUSUM и свободная от распределения карта CUSUM-SN также не подают сигнал.

\noindent\textbf{Пример 4.13. Карта EWMA-SN: подход на основе марковских цепей.}

Здесь мы рассматриваем карту EWMA-SN для отдельных данных с параметрами проектирования \(\lambda = 0.05\), \(L = 2\) и \(v = 2s + 1 = 9\), так что \(s = 4\), то есть делим интервал между контрольными границами на девять подынтервалов (см. рисунок \ref{fig-4.19}); в результате получается «малая» матрица переходных вероятностей, имеющая всего 10 строк и 10 столбцов. Затем иллюстрируем, как можно вычислить распределение длины серии методом MC. Отметим, что для этой карты стационарные контрольные границы вычисляются по уравнению \eqref{eq:4.15} и равны \(LCL = -0.320\) и \(UCL = 0.320\).

\begin{figure}[!ht]
	\centering
	\includegraphics[width=0.45\linewidth]{"./images/fig-4.19.png"}
	\caption[]{Разбиение интервала между LCL и UCL на \(\nu = 9\) подынтервалов для \(\lambda = 0.05\) и \(L = 2\).}
	\label{fig-4.19}
\end{figure}

Используя эти девять состояний, матрица переходных вероятностей имеет вид
\[
	M_{10\times10} = \left(\begin{array}{c|c} Q_{9\times9} & p_{9\times1} \\ \hline 0_{1\times9} & 1_{1\times1} \end{array}\right),
\]
где в состоянии IC она равна
\[
	\frac{1}{2}\left(\begin{array}{ccccccccc|c}
		0 & 1 & 0 & 0 & 0 & 0 & 0 & 0 & 0 & 1 \\
		1 & 0 & 1 & 0 & 0 & 0 & 0 & 0 & 0 & 0 \\
		0 & 1 & 0 & 1 & 0 & 0 & 0 & 0 & 0 & 0 \\
		0 & 0 & 1 & 0 & 1 & 0 & 0 & 0 & 0 & 0 \\
		0 & 0 & 0 & 1 & 0 & 1 & 0 & 0 & 0 & 0 \\
		0 & 0 & 0 & 0 & 1 & 0 & 1 & 0 & 0 & 0 \\
		0 & 0 & 0 & 0 & 0 & 1 & 0 & 1 & 0 & 0 \\
		0 & 0 & 0 & 0 & 0 & 0 & 1 & 0 & 1 & 0 \\
		0 & 0 & 0 & 0 & 0 & 0 & 0 & 1 & 0 & 1 \\
		\hline
		0 & 0 & 0 & 0 & 0 & 0 & 0 & 0 & 0 & 2
	\end{array}\right),
\]
а переходные вероятности вычислены по уравнению \eqref{eq:4.19}. Для иллюстрации ниже приведено вычисление \(p_{01}\):
\begin{align*}
	p_{01} &= P\left( \frac{(S_1 - \tau) - (1 - \lambda)S_0}{\lambda} + \frac{n}{2} < T_k \le \frac{(S_1 + \tau) - (1 - \lambda)S_0}{\lambda} + \frac{n}{2} \right) \\
	      &= P\left( \frac{(0.071 - 0.036) - (1 - 0.05)0}{0.05} + \frac{1}{2} < T_k \le \frac{(0.071 + 0.036) - (1 - 0.05)0}{0.05} + \frac{1}{2} \right) \\
	      &= P(0.856 < T_k \le 1.568) = P(T_k = 1) = \frac{1}{2}.
\end{align*}
Предпоследнее равенство следует из того, что случайная величина \(T_k\), имеющая в состоянии IC распределение \(BIN(n = 1, p = 0.5)\), принимает только неотрицательные целые значения (\(T_k \in \{0, 1\}\)). Как только одношаговые переходные вероятности вычислены и найдена существенная матрица переходных вероятностей \(Q\), характеристики длины серии в состоянии IC получаются по уравнениям 3.22--3.24. Это даёт \(ARL_{IC} = 25.00\) и \(SDRL_{IC} = 20.00\) при значениях 5-го, 25-го, 50-го, 75-го и 95-го процентилей в состоянии IC, равных 5, 11, 19, 33 и 65 соответственно.

Значение \(ARL_{IC}\) означает, что, когда процесс находится в состоянии IC, первый сигнал будет наблюдаться в среднем на каждой 25-й нанесённой точке. Первый квартиль равен 11, поэтому мы знаем, что ложная тревога не возникнет на первых 11 нанесённых точках с вероятностью не более 75\%. Другой пример: \(MRL_{IC} = 19\), что означает, что первая ложная тревога будет наблюдаться в пределах первых 19 нанесённых точек как минимум в 50\% случаев.

Опять же, следует отметить, что, хотя здесь показаны только характеристики в состоянии IC, характеристики в состоянии OOC легко получить, положив \(p \ne 0.5\) в приведённых выше формулах.

\noindent\textbf{Пример 4.14. Карта EWMA-SR.}

Двусторонняя карта CUSUM-SR была показана в примере 4.9 по данным таблицы \ref{tab:4.12}. В качестве альтернативы теперь иллюстрируется контрольная карта EWMA-SR. Для этого используем \(\lambda = 0.15\) для вычисления статистик для нанесения на карту и \(L = 2.25\) для нахождения контрольных границ (см. Graham, Chakraborti and Human, 2011b). Для иллюстрации расчётов рассмотрим выборку номер 1 с \(SR_1 = 2\) и предположим, что начальное значение \(Z_0 = 0\). Тогда первая статистика для нанесения на карту на карте EWMA-SR (рисунок \ref{fig-4.20}) равна
\begin{align*}
	Z_1 &= \lambda SR_1 + (1 - \lambda)Z_0 \\
	    &= 0.15(2) + (1 - 0.15)0 \\
	    &= 0.30.
\end{align*}
Значение второй статистики для нанесения на карту равно
\begin{align*}
	Z_2 &= \lambda SR_2 + (1 - \lambda)Z_1 \\
	    &= 0.15(-10) + (1 - 0.15)0.30 \\
	    &= -1.245,
\end{align*}
а остальные расчёты сведены в столбец (b) таблицы \ref{tab:4.17}.

\begin{table}[!ht]
	\centering
	\caption{Расчёты для карты EWMA-SR.}
	\label{tab:4.17}
	\begin{tabular}{@{}ccc@{}}\toprule
		Номер выборки \(i\) & (a) \(SR_i\) & (b) \(Z_t = \lambda SR_t + (1 - \lambda)Z_{t-1}\) \\ \midrule
		1 & \(2\) & \(0.300\) \\
		2 & \(-10\) & \(-1.245\) \\
		3 & \(10\) & \(0.442\) \\
		4 & \(-8\) & \(-0.825\) \\
		5 & \(-2\) & \(-1.001\) \\
		6 & \(2\) & \(-0.551\) \\
		7 & \(10\) & \(1.032\) \\
		8 & \(-2\) & \(0.577\) \\
		9 & \(-8\) & \(-0.709\) \\
		10 & \(-8\) & \(-1.803\) \\
		11 & \(10\) & \(-0.033\) \\
		12 & \(6\) & \(0.872\) \\
		13 & \(-8\) & \(-0.459\) \\
		14 & \(4\) & \(0.210\) \\
		15 & \(-2\) & \(-0.121\) \\
		16 & \(0\) & \(-0.103\) \\
		17 & \(2\) & \(0.212\) \\
		18 & \(10\) & \(1.681\) \\
		19 & \(-6\) & \(0.528\) \\
		20 & \(-2\) & \(0.149\) \\ \bottomrule
	\end{tabular}
\end{table}

Для простоты используем стационарные контрольные границы на карте EWMA-SR, которые находятся по уравнению \eqref{eq:4.22}:
\begin{align*}
	LCL &= -L\sigma_{SR_i}\sqrt{\frac{\lambda}{2 - \lambda}} \\
	    &= -2.25\sqrt{\frac{4(4 + 1)(2 \times 5 + 1)}{24}}\sqrt{\frac{0.15}{2 - 0.15}} \\
	    &= -1.94
\end{align*}
и
\begin{align*}
	UCL &= L\sigma_{SR_i}\sqrt{\frac{\lambda}{2 - \lambda}} \\
	    &= 2.25\sqrt{\frac{4(4 + 1)(2 \times 5 + 1)}{24}}\sqrt{\frac{0.15}{2 - 0.15}} \\
	    &= 1.94
\end{align*}
соответственно.

\begin{figure}[!ht]
	\centering
	\includegraphics[width=0.7\linewidth]{"./images/fig-4.20.png"}
	\caption[]{Карта EWMA-SR.}
	\label{fig-4.20}
\end{figure}

Результирующая карта EWMA-SR показана на рисунке \ref{fig-4.20}. Поскольку ни одна из статистик для нанесения на карту не попадает на LCL или UCL либо за них, процесс считается находящимся в состоянии IC, и корректирующих мер не требуется.

Следует отметить, что для тех же данных карта CUSUM-SR подала сигнал на выборке номер 10 при счётчике \(N^- = 2\), что указывает на то, что сдвиг, скорее всего, произошёл на выборке номер 8. Хотя карта EWMA-SR не подаёт сигнал, из рисунка \ref{fig-4.20} видно, что карта EWMA-SR почти подаёт сигнал на выборке номер 10. Похожая картина наблюдается и для параметрической карты EWMA на этих данных, то есть она почти подаёт сигнал на выборке номер 10.

\begin{figure}[!ht]
	\centering
	\includegraphics[width=0.45\linewidth]{"./images/fig-4.21.png"}
	\caption[]{Разбиение интервала между LCL и UCL на \(\nu = 9\) подынтервалов для \(\lambda = 0.1\) и \(L = 3\).}
	\label{fig-4.21}
\end{figure}

\noindent\textbf{Пример 4.15. Карта EWMA-SR: подход на основе марковских цепей.}

Наконец, рассмотрим карту EWMA-SR для \(n = 6\) с параметрами проектирования \(\lambda = 0.1\), \(L = 3\) и \(v = 2s + 1 = 9\), так что \(s = 4\), то есть делим интервал между контрольными границами на девять подынтервалов (см. рисунок \ref{fig-4.21}); в результате получается «малая» матрица переходных вероятностей, имеющая всего 10 строк и 10 столбцов. Затем иллюстрируем, как можно вычислить характеристики распределения длины серии методом MC. Стационарные контрольные границы вычисляются по уравнению \eqref{eq:4.22} и равны \(LCL = -6.565\) и \(UCL = 6.565\).

Используя эти девять состояний, матрица переходных вероятностей имеет вид
\[
	M_{10\times10} = \left(\begin{array}{c|c} Q_{9\times9} & p_{9\times1} \\ \hline 0_{1\times9} & 1_{1\times1} \end{array}\right),
\]
где в состоянии IC она равна
\[
	\frac{1}{64}\left(\begin{array}{ccccccccc|c}
		32 & 24 & 3 & 0 & 0 & 0 & 0 & 0 & 0 & 5 \\
		7 & 30 & 25 & 2 & 0 & 0 & 0 & 0 & 0 & 0 \\
		0 & 10 & 32 & 20 & 2 & 0 & 0 & 0 & 0 & 0 \\
		0 & 0 & 14 & 32 & 17 & 1 & 0 & 0 & 0 & 0 \\
		0 & 0 & 0 & 14 & 36 & 14 & 0 & 0 & 0 & 0 \\
		0 & 0 & 0 & 1 & 17 & 32 & 14 & 0 & 0 & 0 \\
		0 & 0 & 0 & 0 & 2 & 20 & 32 & 10 & 0 & 0 \\
		0 & 0 & 0 & 0 & 0 & 2 & 25 & 30 & 7 & 0 \\
		0 & 0 & 0 & 0 & 0 & 0 & 3 & 24 & 32 & 5 \\
		\hline
		0 & 0 & 0 & 0 & 0 & 0 & 0 & 0 & 0 & 64
	\end{array}\right),
\]
а переходные вероятности вычислены в состоянии IC по уравнению \eqref{eq:4.24}. Для иллюстрации ниже приведено вычисление \(p_{01}\):
\begin{align*}
	p_{01} &= P\left( \frac{(S_1 - \tau) - (1 - \lambda)S_0}{\lambda} + \frac{n(n + 1)}{2}\Big/2 < T_n^+ \right. \\
	      &\qquad \left. \le \frac{(S_1 + \tau) - (1 - \lambda)S_0}{\lambda} + \frac{n(n + 1)}{2}\Big/2 \right) \\
	      &= P\left( \frac{(1.459 - 0.729) - (1 - 0.1)0}{0.1} + \frac{6(6 + 1)}{2}\Big/2 < T_n^+ \right. \\
	      &\qquad \left. \le \frac{(1.459 + 0.729) - (1 - 0.1)0}{0.1} + \frac{6(6 + 1)}{2}\Big/2 \right) \\
	      &= P(14.145 < T_n^+ \le 21.440) \\
	      &= P(T_n^+ = 15) + P(T_n^+ = 16) + P(T_n^+ = 17) + P(T_n^+ = 18) \\
	      &\qquad + P(T_n^+ = 19) + P(T_n^+ = 20) + P(T_n^+ = 21) \\
	      &= \frac{17}{64}.
\end{align*}
Предпоследнее равенство следует из того, что случайная величина \(T_n^+\) принимает только неотрицательные целые значения. Кроме того, отметим, что вероятности вычислены с использованием распределения \(T_n^+\) в состоянии IC, приведённого в таблице B Приложения A.

Как и прежде, как только одношаговые переходные вероятности вычислены и найдена существенная матрица переходных вероятностей \(Q\), характеристики длины серии в состоянии IC получаются по уравнениям 3.22--3.24. Это даёт \(ARL_{IC} = 699.49\) и \(SDRL_{IC} = 690.55\) при значениях 5-го, 25-го, 50-го, 75-го и 95-го процентилей в состоянии IC, равных 44, 208, 488, 966 и 2078 соответственно.

Значение \(ARL_{IC}\) означает, что, когда процесс находится в состоянии IC, первый сигнал будет наблюдаться в среднем примерно на каждой 700-й нанесённой точке. Первый квартиль равен 44, поэтому мы знаем, что ложная тревога не возникнет на первых 44 нанесённых точках с вероятностью не более 75\%. Другой пример: \(MRL_{IC} = 488\), что означает, что первая ложная тревога будет наблюдаться в пределах первых 488 нанесённых точек как минимум в 50\% случаев.

Опять же, следует отметить, что, хотя здесь показаны только характеристики в состоянии IC, характеристики в состоянии OOC можно аналогично найти с использованием распределения \(T_n^+\) в состоянии OOC.

До сих пор мы обсуждали некоторые свободные от распределения контрольные карты в случае K, когда значение параметра в состоянии IC известно. Однако на практике бывают ситуации, когда значения параметров в состоянии IC неизвестны. В этом случае мы имеем дело со случаем U, который рассматривается далее.

\section{Свободные от распределения контрольные карты по количественным признакам в случае U}

\subsection{Контрольные карты Шухарта}
\subsubsection{Контрольные карты Шухарта на основе статистики теста предшествования}
Как показано выше, карты Шухарта и карты CUSUM и EWMA, основанные на знаковой и знаково-ранговой статистиках, полезны при мониторинге положения процесса в случае K. Однако во многих практических ситуациях значение параметра положения в состоянии IC не известно, и его приходится оценивать по эталонной выборке. В этом случае мы имеем дело со случаем U. Как и ранее, возможность использовать непараметрические карты остаётся, поскольку их распределение в состоянии IC свободно от распределения. В этом разделе рассматриваются некоторые из таких карт: карта Shewhart-Prec (основанная на тесте предшествования), карта Shewhart-MW (основанная на статистике Манна--Уитни), карта CUSUM-EX (основанная на статистике превышения), карта CUSUM-rank, карта EWMA-EX и карта EWMA-rank.

Начнём с двухвыборочного теста предшествования (precedence test). Пусть имеются две непрерывные функции распределения \(F\) и \(G\), соответствующие эталонной выборке и выборке, из которой извлекается тестовая выборка (то есть выборка Фазы II). Пусть \(X_1, X_2, \ldots, X_m\) --- эталонная выборка объёма \(m\) из распределения \(F\), а \(Y_1, Y_2, \ldots, Y_n\) --- независимая тестовая выборка объёма \(n\) из распределения \(G\). Определим порядковые статистики \(X_{(1:m)} < X_{(2:m)} < \cdots < X_{(m:m)}\) эталонной выборки и \(Y_{(1:n)} < Y_{(2:n)} < \cdots < Y_{(n:n)}\) тестовой выборки. Тест предшествования основывается на том, сколько \(X\)-наблюдений предшествуют (то есть меньше) \(j\)-й порядковой статистики \(Y_{(j:n)}\) тестовой выборки. Обозначим это число через \(W_j\), то есть
\begin{equation}
	W_j = \sum_{i=1}^{m} I\left( X_i < Y_{(j:n)} \right), \quad j = 1, 2, \ldots, n,
	\label{eq:4.25}
\end{equation}
где \(I(\cdot)\) --- обычная индикаторная функция. Распределение статистики \(W_j\) в состоянии IC подробно выводится в приложении 4.1.2. Оно имеет вид
\begin{equation}
	P_{IC}(W_j = w) = \frac{\binom{w + j - 1}{w}\binom{m + n - w - j}{m - w}}{\binom{m + n}{m}},
	\label{eq:4.26}
\end{equation}
для \(w = 0, 1, \ldots, m\). Это распределение не зависит от \(F\), и поэтому статистика \(W_j\) свободна от распределения. Для иллюстрации на рисунках \ref{fig-4.22a}--\ref{fig-4.22c} показаны распределения вероятностей \(W_j\) в состоянии IC для некоторых значений \(m\) (объём эталонной выборки), \(n\) (объём тестовой выборки) и некоторых значений \(j\).

\begin{figure}[!ht]
	\centering
	\includegraphics[width=0.7\linewidth]{"./images/fig-4.22a.png"}
	\caption{Распределение \(W_j\) в состоянии IC при \(m = 9\), \(n = 7\) и \(j = 4\).}
	\label{fig-4.22a}
\end{figure}

\begin{figure}[!ht]
	\centering
	\includegraphics[width=0.7\linewidth]{"./images/fig-4.22b.png"}
	\caption{Распределение \(W_j\) в состоянии IC при \(m = 12\), \(n = 9\) и \(j = 5\).}
	\label{fig-4.22b}
\end{figure}

\begin{figure}[!ht]
	\centering
	\includegraphics[width=0.7\linewidth]{"./images/fig-4.22c.png"}
	\caption{Распределение \(W_j\) в состоянии IC при \(m = 12\), \(n = 9\) и \(j = 2\).}
	\label{fig-4.22c}
\end{figure}


На рисунке \ref{fig-4.22b}, например, показано распределение \(W_j\) в состоянии IC при \(m = 12\), \(n = 9\) и \(j = 5\), то есть когда \(j\)-я порядковая статистика тестовой выборки является выборочной медианой. Подстановка этих значений в уравнение \eqref{eq:4.26} даёт
\[
	P_{IC}(W_5 = w) = \frac{\binom{w + 4}{w}\binom{16 - w}{12 - w}}{\binom{21}{12}}
\]
для \(w = 0, 1, \ldots, 12\). На рисунке \ref{fig-4.22a} показано распределение \(W_j\) в состоянии IC при \(m = 9\), \(n = 7\) и \(j = 4\). Отметим, что когда \(n\) нечётно и порядковая статистика тестовой выборки выбрана в качестве медианы, то есть \(j = (n + 1)/2\) для карты Med, распределение вероятностей \(W_j\) в состоянии IC симметрично относительно своего среднего, которое в случаях рисунков \ref{fig-4.22a} и \ref{fig-4.22b} равно 4.5 и 6 соответственно. Читатель может найти формулы для среднего и стандартного отклонения распределения \(W_j\) в состоянии IC. На рисунке \ref{fig-4.22c} показано распределение \(W_j\) в состоянии IC при \(m = 12\), \(n = 9\) и \(j = 2\), то есть когда статистикой для нанесения на карту служит вторая порядковая статистика тестовой выборки. Это распределение имеет правостороннюю асимметрию.

Как уже отмечалось, \(W_j\) --- свободная от распределения статистика, то есть её распределение в состоянии IC не зависит от лежащих в основе распределений \(F\) и \(G\). Будет показано, что распределение длины серии контрольной карты, основанной на \(W_j\), можно получить из распределения статистики предшествования \(W_j\), и, следовательно, распределение длины серии этой карты в состоянии IC свободно от распределения. Такую карту обычно называют картой Shewhart-Prec.

Отметим, что распределение вероятностей \(W_j\) в состоянии IC легко допускает тестовые выборки неравного объёма. Например, если каждая тестовая выборка имеет объём \(n_i\), то уравнение \eqref{eq:4.26} сохраняет силу с заменой \(n\) на \(n_i\) и \(j = 1, 2, \ldots, n_i\). Таким образом, карта Shewhart-Prec достаточно гибка, чтобы допускать переменный объём тестовых выборок.

Чтобы показать связь с контрольной картой, определённой в уравнении \eqref{eq:4.26}, заметим, что в силу предположения о непрерывности \(F\) и \(G\) (совпадений нет) \(Y_{(j:n)} \le X_{(b:m)}\) тогда и только тогда, когда не более \(b - 1\) наблюдений \(X\) предшествуют \(Y_{(j:n)}\), то есть когда \(W_j \le b - 1\). Аналогично \(X_{(a:m)} \le Y_{(j:n)}\) тогда и только тогда, когда по крайней мере \(a\) наблюдений \(X\) предшествуют \(Y_{(j:n)}\), то есть когда \(W_j \ge a\). Таким образом, событие отсутствия сигнала \(\{X_{(a:m)} \le Y_{(j:n)} \le X_{(b:m)}\}\) эквивалентно событию \(\{a \le W_j \le b - 1\}\), и, следовательно, контрольную карту можно выразить через статистику предшествования \(W_j\). Процесс следует объявить находящимся в состоянии IC, если \(\{a \le W_j \le b - 1\}\), и в состоянии OOC, если \(\{W_j \le a - 1\}\) или \(\{W_j \ge b\}\). Поэтому эта карта называется картой Shewhart-Prec. В силу эквивалентности событий отсутствия сигнала вероятность того, что карта Shewhart-Prec не подаёт сигнал,
\[
	p = P\left( X_{(a:m)} \le Y_{(j:n)} \le X_{(b:m)} \right) = P(\text{нет сигнала})
\]
может быть выражена через распределение \(W_j\) как
\[
	P\left( X_{(a:m)} \le Y_{(j:n)} \le X_{(b:m)} \right) = P(a \le W_j \le b - 1) = p(m, n, j; F, G),
\]
скажем. Далее будет показано, что распределение длины серии карты Shewhart-Prec в состоянии IC можно выразить через распределение статистики предшествования в состоянии IC, и, следовательно, карта Shewhart-Prec свободна от распределения, поскольку \(W_j\), как показано, свободна от распределения. Это также будет показано далее при изучении распределения длины серии.

\paragraph{Методы вычисления распределения длины серии карты Shewhart-Prec.} Чтобы оценить эффективность карты предшествования, нужно вывести и изучить её распределение длины серии. В случае U, поскольку для оценивания контрольных границ используется эталонная выборка, контрольные границы являются случайными величинами, которые вносят дополнительную изменчивость, и это нужно учитывать при выводе распределения длины серии. Поэтому, как и в случае U для карты \(\overline{X}\) нормальной теории, используем метод обусловливания/разу-словливания (CUC), рассмотренный в главе 3. Для этого есть два основных шага. Сначала мы обусловливаем на двух порядковых статистиках эталонной выборки \(X_{(a:m)}\) и \(X_{(b:m)}\), то есть на двух контрольных границах, чтобы найти условное распределение длины серии. Затем находим безусловное распределение длины серии, беря математическое ожидание (усредняя) по распределению двух порядковых статистик. Однако имеется существенное различие между случаем нормальной теории и непараметрическим случаем: в нормальном случае оценки среднего и стандартного отклонения, определяющие контрольные границы, независимы, а две порядковые статистики, определяющие контрольные границы карты Shewhart-Prec, --- нет. Поэтому шаг взятия математического ожидания (усреднения или разобусловливания) при нахождении безусловного распределения длины серии в непараметрическом случае включает интегрирование по совместному распределению двух порядковых статистик (контрольных границ), тогда как в нормальном случае он включает произведение индивидуальных плотностей вероятности двух независимых оценок.

Начнём с условного распределения длины серии.

\paragraph{Условное распределение длины серии двусторонней карты Shewhart-Prec.} Две порядковые статистики эталонной выборки \(X_{(a:m)}\) и \(X_{(b:m)}\) определяют контрольные границы карты Shewhart-Prec. Условно на этих порядковых статистиках, то есть когда они фиксированы, события сигнала в Фазе II взаимно независимы, поскольку статистики для нанесения на карту \(Y_{(j:n)}\) происходят из независимых выборок Фазы II. Таким образом, условно на порядковых статистиках эталонной выборки распределение длины серии карты Shewhart-Prec является геометрическим распределением с параметром \(1 - p_C\), где \(p_C\) --- условная вероятность отсутствия сигнала, \(p_C = P_C\left( x_{(a:m)} \le Y_{(j:n)} \le x_{(b:m)} \mid X_{(a:m)} = x_{(a:m)}, X_{(b:m)} = x_{(b:m)} \right)\), которая задаётся как
\begin{equation}
\begin{aligned}
	p_C &= P_C\left( F^{-1}\left(u_{(a:m)}\right) \le G^{-1}\left(U_{(j:n)}\right) \le F^{-1}\left(u_{(b:m)}\right) \right) \\
	    &= P_C\left( GF^{-1}\left(u_{(a:m)}\right) \le U_{(j:n)} \le GF^{-1}\left(u_{(b:m)}\right) \right),
\end{aligned}
\label{eq:4.27}
\end{equation}
где \(u_{(e:l)}\) обозначает наблюдённое значение \(e\)-й порядковой статистики \(U_{(e:l)}\) в выборке объёма \(l\) из равномерного распределения \(Uniform(0, 1)\). Отметим использование вероятностного интегрального преобразования (PIT) (см. раздел 1.4), согласно которому, например, \(G\left(Y_{(j:n)}\right)\) распределена как \(U_{(j:n)}\), \(j\)-я порядковая статистика в выборке объёма \(n\) из равномерного \(Uniform(0, 1)\) распределения, которая, как можно показать, имеет бета-распределение с параметрами \(j\) и \(n - j + 1\) (см. раздел 1.4) с плотностью \(f(u) = \frac{1}{B(j, n - j + 1)}u^{j-1}(1 - u)^{n-j}\).

Таким образом,
\[
	p_C = \int_{GF^{-1}(u_{(a:m)})}^{GF^{-1}(u_{(b:m)})} f(u)\,du = F_2\left( u_{(a:m)}, u_{(b:m)}, j, n, F, G \right),
\]
например.

Следовательно, условное распределение случайной величины длины серии \(N\), то есть \(\left( N \mid X_{(a:m)}, X_{(b:m)} \right)\), является \(GEO(1 - p_C)\), и поэтому функция вероятности условного распределения длины серии имеет вид
\begin{equation}
	P_C\left( N = k \mid X_{(a:m)}, X_{(b:m)} \right) = p_C^{k-1}(1 - p_C) = p_C^{k-1} - p_C^{k}
	\label{eq:4.28}
\end{equation}
для \(k = 1, 2, 3, \ldots\). Отметим, что полная запись условия имеет вид \(X_{(a:m)} = x_{(a:m)}\), \(X_{(b:m)} = x_{(b:m)}\), однако для краткости пишем \(X_{(a:m)}, X_{(b:m)}\). Все свойства условного распределения длины серии следуют из свойств геометрического распределения. Например, функция распределения условной длины серии равна
\begin{equation}
	P_C\left( N \le k \mid X_{(a:m)}, X_{(b:m)} \right) = 1 - p_C^{k}
	\label{eq:4.29}
\end{equation}
для \(k = 1, 2, 3, \ldots\). Условная средняя длина серии равна
\begin{equation}
	ARL_C = E\left( N \mid X_{(a:m)}, X_{(b:m)} \right) = \frac{1}{1 - p_C}.
	\label{eq:4.30}
\end{equation}
Аналогично условная дисперсия длины серии равна
\begin{equation}
	VRL_C = \operatorname{var}\left( N \mid X_{(a:m)}, X_{(b:m)} \right) = \frac{p_C}{(1 - p_C)^2}.
	\label{eq:4.31}
\end{equation}
Отсюда следует, что условное стандартное отклонение длины серии равно
\begin{equation}
	SDRL_C = SD\left( N \mid X_{(a:m)}, X_{(b:m)} \right) = \sqrt{\frac{p_C}{(1 - p_C)^2}}.
	\label{eq:4.32}
\end{equation}
Отметим, что в уравнениях \eqref{eq:4.29}--\eqref{eq:4.32} индекс \(C\) означает «условное» (conditional), как и для параметрических карт в главе 3.

Когда процесс находится в состоянии IC, \(p_C\) сводится к
\begin{equation}
	p_{C, IC} = P_{C, IC}\left( u_{(a:m)} \le U_{(j:n)} \le u_{(b:m)} \mid U_{(a:m)}, U_{(b:m)} \right) = \int_{u_{(a:m)}}^{u_{(b:m)}} f(u)\,du.
	\label{eq:4.33}
\end{equation}
Опять же отметим, что полная запись условия имеет вид \(U_{(a:m)} = u_{(a:m)}\), \(U_{(b:m)} = u_{(b:m)}\), однако для краткости пишем \(U_{(a:m)}, U_{(b:m)}\).

В уравнении \eqref{eq:4.33} индекс \(C, IC\) означает «условное в состоянии IC» (conditional in-control). Таким образом, \(p_{C, IC}\) обозначает условную вероятность отсутствия сигнала в состоянии IC. Поэтому функцию распределения, среднее, дисперсию и стандартное отклонение условного распределения длины серии в состоянии IC можно получить из уравнений \eqref{eq:4.29}--\eqref{eq:4.32} заменой \(p_C\) на \(p_{C, IC}\).

Кроме того, разлагая член \((1 - u)^{n-j}\) в бета-плотности в выражении для \(p_C\) в биномиальную сумму и затем интегрируя, можно показать, что условную вероятность отсутствия сигнала можно выразить как
\begin{equation}
\begin{aligned}
	p_C &= F_2\left( u_{(a:m)}, u_{(b:m)}, j, n, F, G \right) \\
	    &= \frac{1}{B(j, n - j + 1)}\sum_{h=0}^{n-j} \frac{(-1)^h \binom{n - j}{h}}{j + h} \\
	    &\qquad \times \left[ \left(GF^{-1}\left(u_{(b:m)}\right)\right)^{j+h} - \left(GF^{-1}\left(u_{(a:m)}\right)\right)^{j+h} \right].
\end{aligned}
\label{eq:4.34}
\end{equation}
Последнее выражение может оказаться полезным для вычислений. Условное распределение длины серии и его различные характеристики, такие как среднее, стандартное отклонение, процентили и т.~д., являются случайными величинами, и они могут дать полезное представление о влиянии оценивания параметров на контрольную карту, в частности на карту Shewhart-Prec. Само условное распределение длины серии, например, можно проанализировать, чтобы изучить влияние объёма эталонной выборки на эффективность контрольной карты. Ряд исследователей в недавней литературе по SPC предлагали уделять больше внимания условному распределению длины серии при проектировании карты Фазы II, но мы здесь эти детали не рассматриваем. Далее мы выводим безусловное распределение длины серии карты Shewhart-Prec. Безусловное распределение длины серии и его характеристики более традиционно используются в литературе для изучения эффективности карт.

\paragraph{Безусловное распределение длины серии двусторонней карты Shewhart-Prec.} Функцию вероятности безусловного распределения длины серии, то есть \(P_U(N = k)\), и его различные характеристики можно легко найти с помощью метода CUC. Отметим, что индекс \(U\) означает «безусловное» (unconditional). Например, безусловная функция вероятности имеет вид
\begin{equation}
	P_U(N = k) = E_{X_{(a:m)}, X_{(b:m)}}\left( P_C\left( N = k \mid X_{(a:m)}, X_{(b:m)} \right) \right),
	\label{eq:4.35}
\end{equation}
которая с использованием уравнения \eqref{eq:4.28} упрощается к
\begin{equation}
	P_U(N = k) = E_{X_{(a:m)}, X_{(b:m)}}\left( p_C^{k-1} \right) - E_{X_{(a:m)}, X_{(b:m)}}\left( p_C^{k} \right) = D_U^*(k - 1) - D_U^*(k)
	\label{eq:4.36}
\end{equation}
для \(k = 1, 2, 3, \ldots\), где
\begin{equation}
\begin{aligned}
	D_U(k) &= E_{X_{(a:m)}, X_{(b:m)}}\left( p_C^{k} \right) \\
	       &= \int_{0}^{1}\int_{0}^{u_{(b:m)}} \left( \int_{GF^{-1}(u_{(a:m)})}^{GF^{-1}(u_{(b:m)})} f(u)\,du \right)^{k} f\left(u_{(a:m)}, u_{(b:m)}\right) du_{(a:m)} du_{(b:m)} \\
	       &= \int_{0}^{1}\int_{0}^{u_{(b:m)}} \left( F_2\left( u_{(a:m)}, u_{(b:m)}, j, n, F, G \right) \right)^{k} \\
	       &\qquad \times f\left(u_{(a:m)}, u_{(b:m)}\right) du_{(a:m)} du_{(b:m)},
\end{aligned}
\label{eq:4.37}
\end{equation}
при \(D_U(0) = 1\) и
\[
	f\left(u_{(a:m)}, u_{(b:m)}\right) = C\left(u_{(a:m)}\right)^{a-1}\left(u_{(b:m)} - u_{(a:m)}\right)^{b-a-1}\left(1 - u_{(b:m)}\right)^{m-b},
\]
где \(C = \frac{m!}{(a - 1)!(b - a - 1)!(m - b)!}\), а \(f\left(u_{(a:m)}, u_{(b:m)}\right)\) --- совместная плотность \(U_{(a:m)}\) и \(U_{(b:m)}\) в выборке объёма \(m\) из равномерного \(Uniform(0, 1)\) распределения.

Безусловное распределение длины серии карты Shewhart-Prec в состоянии IC получается при \(G = F\), и мы обозначаем его \(P_{U, IC}\). Таким образом, безусловная функция вероятности длины серии в состоянии IC имеет вид
\begin{equation}
	P_{U, IC}(N = k) = D_U(k - 1) - D_U(k)
	\label{eq:4.38}
\end{equation}
для \(k = 1, 2, 3, \ldots\), где
\begin{equation}
	D_U(k) = \int_{0}^{1}\int_{0}^{u_{(b:m)}} \left( F_2\left( u_{(a:m)}, u_{(b:m)}, j, n, F, F \right) \right)^{k} f\left(u_{(a:m)}, u_{(b:m)}\right) du_{(a:m)} du_{(b:m)}
	\label{eq:4.39}
\end{equation}
при \(D_U(0) = 1\). Эту формулу можно использовать рекурсивно, чтобы найти безусловную функцию вероятности длины серии в состоянии IC.

Функцию распределения безусловной длины серии в общем случае можно получить из уравнений \eqref{eq:4.35}--\eqref{eq:4.39}. Так, например, можно получить безусловную функцию распределения длины серии в состоянии IC; она равна
\begin{equation}
	P_{U, IC}(N \le k) = 1 - D_U(k)
	\label{eq:4.40}
\end{equation}
для \(k = 1, 2, \ldots\)

\paragraph{Безусловные характеристики эффективности длины серии двусторонней карты Shewhart-Prec.} Различные характеристики эффективности карты, такие как безусловная средняя длина серии, безусловная медианная длина серии, безусловное стандартное отклонение длины серии и т.~д., можно найти по функции вероятности или функции распределения длины серии. В качестве альтернативы для безусловного среднего и стандартного отклонения длины серии можно напрямую использовать метод CUC, воспользовавшись тем, что условное распределение длины серии является геометрическим. Такой подход может оказаться выгодным в других ситуациях при вычислении первых двух моментов безусловного распределения длины серии. Чтобы проиллюстрировать это, сначала отметим, что, как показано ранее в уравнении \eqref{eq:4.30}, условная средняя длина серии равна
\begin{equation}
	ARL_C = \frac{1}{1 - p_C},
	\label{eq:4.41}
\end{equation}
и, следовательно, безусловную среднюю длину серии двусторонней карты Shewhart-Prec, \(ARL_U\), можно получить, взяв математическое ожидание \(ARL_C\) по совместному распределению \(U_{(a:m)}\) и \(U_{(b:m)}\), двух порядковых статистик в выборке объёма \(m\) из равномерного \(Uniform(0, 1)\) распределения. Следовательно,
\begin{equation}
\begin{aligned}
	ARL_U &= \int_{0}^{1}\int_{0}^{u_{(b:m)}} \frac{1}{1 - p_C}\, f\left(u_{(a:m)}, u_{(b:m)}\right) du_{(a:m)} du_{(b:m)} \\
	      &= \int_{0}^{1}\int_{0}^{u_{(b:m)}} \left[ \frac{1}{1 - F_2\left( u_{(a:m)}, u_{(b:m)}, j, n, F, G \right)} \right] \\
	      &\qquad \times f\left(u_{(a:m)}, u_{(b:m)}\right) du_{(a:m)} du_{(b:m)}.
\end{aligned}
\label{eq:4.42}
\end{equation}
Безусловная средняя длина серии в состоянии IC получается при \(G = F\). Таким образом, уравнение \eqref{eq:4.42} сводится к
\begin{equation}
\begin{aligned}
	ARL_{U, IC} &= \int_{0}^{1}\int_{0}^{u_{(b:m)}} \left[ \frac{1}{1 - F_2\left( u_{(a:m)}, u_{(b:m)}, j, n, F, F \right)} \right] \\
	            &\qquad \times f\left(u_{(a:m)}, u_{(b:m)}\right) du_{(a:m)} du_{(b:m)}.
\end{aligned}
\label{eq:4.43}
\end{equation}

Аналогично для дисперсии также можно использовать метод CUC. Для этого, как отмечено в главе 3, дисперсию любой случайной величины \(T\) можно получить по формуле
\[
	\operatorname{var}(T) = E\left( \operatorname{var}(T \mid W) \right) + \operatorname{var}\left( E(T \mid W) \right).
\]
Применяя эту формулу к переменной длины серии \(N\) в качестве \(T\) и к обусловливающей переменной \(W = \left( X_{(a:m)}, X_{(b:m)} \right)\), находим условную дисперсию и математическое ожидание в правой части этой формулы из уравнений \eqref{eq:4.31} и \eqref{eq:4.30} соответственно. Следовательно, безусловная дисперсия длины серии равна
\[
	VRL_U = E\left( VRL_C \right) + \operatorname{var}\left( ARL_C \right) = E\left( \frac{p_C}{(1 - p_C)^2} \right) + \operatorname{var}\left( \frac{1}{1 - p_C} \right),
\]
где математическое ожидание снова берётся по совместной плотности \(U_{(a:m)}\) и \(U_{(b:m)}\). Таким образом, метод CUC может быть очень полезным приёмом при вычислении распределения длины серии и его первых двух моментов.

Кроме того, как отмечалось ранее, медианная длина серии (MRL) является дополнительной важной характеристикой эффективности карты, которая должна играть более заметную роль в SPC, поскольку распределение длины серии имеет правостороннюю асимметрию. MRL двусторонней карты Shewhart-Prec в общем случае можно удобно вычислить по функции распределения длины серии. Например, безусловную MRL в состоянии IC можно найти, решая относительно наименьшего целого \(MRL_{IC}\) неравенство \(P_{U, IC}\left( N \le MRL_{IC} \right) \ge 0.5\), которое с использованием уравнения \eqref{eq:4.40} приводит к решению неравенства \(D_U\left( MRL_{IC} \right) \le 0.5\), где \(D_U(\cdot)\) определена в уравнении \eqref{eq:4.39}. Тогда можно показать (мы оставляем это читателю в качестве упражнения), что \(MRL_{IC}\) задаётся решением уравнения
\begin{equation}
\begin{aligned}
	&\int_{0}^{1}\int_{0}^{u_{(b:m)}} F_2\left( u_{(a:m)}, u_{(b:m)}, m, n, F, F \right)^{MRL_{IC}} \\
	&\qquad \times f\left(u_{(a:m)}, u_{(b:m)}\right) du_{(a:m)} du_{(b:m)} = 0.5,
\end{aligned}
\label{eq:4.44}
\end{equation}
где \(F_2\left( u_{(a:m)}, u_{(b:m)}, m, n, F, F \right)\) задана в уравнении \eqref{eq:4.34}.

Для решения уравнения \eqref{eq:4.44} относительно \(MRL_{IC}\) (или MRL в общем случае) потребуются некоторые нелинейные методы; детали мы оставляем читателю. Другие процентили распределения длины серии в состоянии IC можно получить аналогично. Можно отметить, что для нахождения процентилей безусловного распределения длины серии в состоянии OOC требуется задать функцию преобразования \(GF^{-1}\), необходимую для вычисления функции вероятности или функции распределения длины серии. Мы затронем этот вопрос в разделе сравнения эффективности далее.

Хотя у нас есть точные формулы для вычисления функции вероятности или функции распределения длины серии карты Shewhart-Prec, для оценки которых требуется численное интегрирование, как это часто бывает на практике со многими контрольными картами, симуляции могут дать более быструю альтернативу. Шаги описаны в следующем разделе.

\paragraph{Симуляционный подход.} Для получения характеристик длины серии карты Shewhart-Prec можно использовать следующий алгоритм симуляции. Соответствующая программа SAS приведена в приложении B, программа SAS номер 7, а вывод показан на рисунке \ref{fig-4.23}.

\noindent\textbf{Шаг 1.} После задания объёма эталонной выборки Фазы I (\(m \ge 1\)), объёма тестовой выборки Фазы II (\(n \ge 1\)) и величины сдвига (\(\delta\)), который нужно обнаружить, задаём два распределения для генерации выборок Фазы I и Фазы II соответственно. В случае IC оба распределения совпадают (говорят, что сдвиг равен нулю (\(\delta = 0\))), поэтому для генерации наблюдений можно использовать любое непрерывное распределение. В случае OOC напомним, что распределение выборки Фазы II предполагается того же вида, что и распределение выборки Фазы I, но со сдвигом среднего в единицах стандартного отклонения совокупности. Выборка Фазы I, например, может быть извлечена из нормального распределения со средним 0 и стандартным отклонением 1, а затем выборка Фазы II может быть извлечена из нормального распределения со средним \(\mu = \delta \ne 0\) и стандартным отклонением 1 (говорят, что сдвиг среднего \(\mu\) не равен нулю (\(\delta \ne 0\))). Для генерации наблюдений можно использовать любое общее непрерывное распределение, симметричное или асимметричное.

\begin{figure}[!ht]
	\centering
	\includegraphics[width=0.75\linewidth]{"./images/fig-4.23.png"}
	\caption[]{Вывод SAS распределения RL при параметрах проектирования (\(m = 1000\), \(n = 5\), \(j = 3\), \(a = 48\), \(b = 953\)) карты Shewhart-Prec, выбранных для номинальной \(ARL_{IC} = 500\).}
	\label{fig-4.23}
\end{figure}

\noindent\textbf{Шаг 2.} Выбираем объём тестовой выборки \(n\) и значение \(j\) для нужной порядковой статистики в тестовой выборке. Именно она служит статистикой для нанесения на карту; например, можно выбрать выборочную медиану --- это даст карту Med.

\noindent\textbf{Шаг 3.} Задаём нужную номинальную \(ARL_{IC}\), например 370 или 500.

\noindent\textbf{Шаг 4.} Константы для нанесения на карту \(a\) и \(b\) находятся одним из двух способов. Первый способ заключается в использовании таблиц из Chakraborti, van der Laan and Van de Wiel (2004). Однако они приводят константы только для \(m = 50, 100, 500\) и 1000 соответственно и для \(n = 5, 11\) и 25 соответственно при нужных номинальных значениях \(ARL_{IC}\), равных 370, 500 и 1000. Если у читателя другая комбинация \((m, n)\) или он хочет использовать другое нужное номинальное значение \(ARL_{IC}\), следует использовать второй способ, который заключается в «поисковом алгоритме». В этом методе \(ARL_{IC}\) вычисляется для различных комбинаций \((a, b)\) и выбирается комбинация, ближайшая к нужной номинальной \(ARL_{IC}\). Более подробно о «поисковом алгоритме» говорится в следующем разделе, где подробно обсуждается реализация карты.

\noindent\textbf{Шаг 5.} Оцениваемые контрольные границы Шухарта задаются соответствующими порядковыми статистиками эталонной выборки \(\widehat{LCL} = X_{(a:m)}\) и \(\widehat{UCL} = X_{(b:m)}\) соответственно.

\noindent\textbf{Шаг 6.} Число подгрупп, необходимых до первого попадания \(W_j\) на контрольные границы или за них, записывается как наблюдение из распределения длины серии.

\noindent\textbf{Шаг 7.} Повторяем шаги 1--6 в общей сложности 100\,000 раз.

\noindent\textbf{Шаг 8.} Получив «набор данных» из 100\,000 наблюдений распределения длины серии, можно воспользоваться процедурой \texttt{proc univariate} пакета SAS для получения характеристик длины серии. Для \(\delta = 0\) это будут характеристики длины серии в состоянии IC, а для \(\delta \ne 0\) --- характеристики длины серии в состоянии OOC.

\paragraph{Реализация двусторонней карты Shewhart-Prec.} Чтобы реализовать двустороннюю карту Shewhart-Prec, нужны значения индексов \(a\) и \(b\), позволяющие найти контрольные границы, то есть две порядковые статистики эталонной выборки \(X_{(a:m)}\) и \(X_{(b:m)}\), по эталонным данным. Как и в случае параметрических карт, эти контрольные границы обычно определяются для заданного номинального значения \(ARL_{IC}\). Таким образом, проектируя двустороннюю карту предшествования с заданным номинальным значением \(ARL_{IC}\), скажем 370, мы находим два положительных целых \(a\) и \(b\), \(1 \le a < b \le m\), при заданных \(m\), \(n\) и \(j\), такие что \(ARL_{IC} = 370\).

Напомним, что безусловная средняя длина серии в состоянии IC задана в уравнении \eqref{eq:4.43}. Решение уравнения \(ARL_{IC} = 370\) относительно \(a\) и \(b\) в общем случае можно найти с помощью «поискового алгоритма» (в котором \(a\) и \(b\) изменяются одновременно так, что \(1 \le a < b \le m\)), который останавливается, когда достигнутая средняя длина серии в состоянии IC близка к 370. Для этого полезно отметить, что в общем случае ARL возрастает по \(b\) и убывает по \(a\). Для карты Med требуемые вычисления относительно проще, поскольку здесь можно использовать симметрию распределения статистики предшествования в состоянии IC, положить \(b = m - a + 1\) и решить уравнение \(ARL_{IC} = 370\) относительно \(a\). Это обсуждается в следующем разделе.

\paragraph{Двусторонняя карта Med.} Чтобы реализовать двустороннюю карту Med, как отмечено ранее, полагаем \(b = m - a + 1\) и решаем уравнения \(ARL_{IC} = 370\), 500 и 1000 соответственно относительно \(a\). Значения, приведённые в таблице \ref{tab:4.18}, взяты из Chakraborti, van der Laan and Van de Wiel (2004).

Отметим, что поскольку распределение длины серии карты предшествования в состоянии IC дискретно, скорее всего, номинальная \(ARL_{IC}\) не будет достигнута точно. Мы используем консервативный подход, при котором соответствующая достигнутая \(ARL_{IC}\) «близка» к номинальному значению. Например, для \(m = 100\), \(n = 5\) и \(j = 3\) при номинальной \(ARL_{IC} = 370\) из таблицы \ref{tab:4.18} контрольные границы карты Med задаются 6-й и 95-й порядковыми статистиками эталонной выборки. Достигнутая \(ARL_{IC}\) для этой карты равна 359.6, а соответствующая достигнутая FAR равна 0.00554. Значение 359.6 хорошо согласуется с традиционным 370, но отметим, что достигнутая FAR равна не 0.0027, а почти вдвое больше. Это одно из проявлений оценивания параметров, присущего случаю U.

Статистические свойства карты Med, включая распределение длины серии и связанные с ним характеристики, можно получить из общих рассуждений о картах Shewhart-Prec, приведённых ранее. Характеристики длины серии карты Med можно найти, подставив \(2s + 1\) вместо \(n\), \(s + 1\) вместо \(j\) и \(m - a + 1\) вместо \(b\) в уравнения \eqref{eq:4.26}--\eqref{eq:4.44}; мы оставляем это читателю в качестве упражнения.

\begin{table}[!ht]
	\centering
	\caption{Индексы контрольной карты \((a, b)\), достигнутая \(ARL_{IC}\) и достигнутая FAR для двусторонней карты Med при заданной средней длине серии в состоянии IC, \(m\), \(n\) и \(j\).}
	\label{tab:4.18}
	\small
	\begin{tabular}{@{}rrrcccc@{}}
		\toprule
		& & & \multicolumn{4}{c}{\(m\)} \\
		\cmidrule(lr){4-7}
		Нужная \(ARL_{IC}\) & \(n\) & \(j\) & 100 & 200 & 500 & 1000 \\
		\midrule
		\multirow{9}{*}{370} & \multirow{3}{*}{5} & \multirow{3}{*}{3} & \((6,95)\) & \((11,190)\) & \((27,474)\) & \((53,948)\) \\
		& & & 359.6 & 385.5 & 365.5 & 374.4 \\
		& & & 0.00554 & 0.00375 & 0.00332 & 0.00287 \\
		& \multirow{3}{*}{11} & \multirow{3}{*}{6} & \((14,87)\) & \((27,174)\) & \((67,434)\) & \((132,869)\) \\
		& & & 367.8 & 372.4 & 355.3 & 377.3 \\
		& & & 0.00627 & 0.00412 & 0.00332 & 0.00288 \\
		& \multirow{3}{*}{25} & \multirow{3}{*}{13} & \((24,77)\) & \((46,155)\) & \((114,387)\) & \((227,774)\) \\
		& & & 316.5 & 375.7 & 367.3 & 377.0 \\
		& & & 0.01108 & 0.00526 & 0.00354 & 0.00301 \\
		\midrule
		\multirow{9}{*}{500} & \multirow{3}{*}{5} & \multirow{3}{*}{3} & \((6,95)\) & \((10,191)\) & \((24,477)\) & \((48,953)\) \\
		& & & 359.6 & 520.1 & 520.3 & 501.9 \\
		& & & 0.00554 & 0.00291 & 0.00227 & 0.00216 \\
		& \multirow{3}{*}{11} & \multirow{3}{*}{6} & \((13,88)\) & \((26,175)\) & \((63,438)\) & \((125,876)\) \\
		& & & 574.5 & 460.5 & 497.2 & 506.9 \\
		& & & 0.00449 & 0.00342 & 0.00242 & 0.00225 \\
		& \multirow{3}{*}{25} & \multirow{3}{*}{13} & \((23,78)\) & \((45,156)\) & \((111,390)\) & \((219,782)\) \\
		& & & 510.8 & 472.5 & 480.1 & 492.2 \\
		& & & 0.00804 & 0.00436 & 0.00276 & 0.00233 \\
		\midrule
		\multirow{9}{*}{1000} & \multirow{3}{*}{5} & \multirow{3}{*}{3} & \((5,96)\) & \((8,193)\) & \((19,482)\) & \((38,963)\) \\
		& & & 678.4 & 1067.5 & 1055.8 & 1005.9 \\
		& & & 0.00352 & 0.00161 & 0.00109 & 0.00111 \\
		& \multirow{3}{*}{11} & \multirow{3}{*}{6} & \((11,90)\) & \((23,178)\) & \((56,445)\) & \((110,891)\) \\
		& & & 940.6 & 927.8 & 953.2 & 1021.7 \\
		& & & 0.00212 & 0.00188 & 0.00131 & 0.00109 \\
		& \multirow{3}{*}{25} & \multirow{3}{*}{13} & \((22,79)\) & \((42,159)\) & \((103,398)\) & \((206,795)\) \\
		& & & 854.6 & 986.2 & 1033.6 & 1004.4 \\
		& & & 0.00573 & 0.00240 & 0.00136 & 0.00117 \\
		\bottomrule
	\end{tabular}
	\par\smallskip
	\footnotesize В каждой ячейке для данной \(m\) указаны индекс \((a,b)\), достигнутая \(ARL_{IC}\) и достигнутая FAR.
\end{table}

\paragraph{Эффективность двусторонней карты Shewhart-Prec.} Эффективность любой контрольной карты имеет первостепенное значение как для практика, так и для исследователя. Поэтому, как отмечалось ранее, при изучении эффективности карты предшествования представляет интерес распределение длины серии. Например, если положение процесса сдвигается от значения в состоянии IC, нас может интересовать вероятность того, что карта обнаружит сдвиг (или подаст сигнал об изменении) на первой тестовой выборке, на второй тестовой выборке или, более общо, на \(k\)-й последующей тестовой выборке. Аналогично нас интересует, например, средняя длина серии в состоянии OOC (\(ARL_{OOC}\)), то есть сколько в среднем (сколько выборок Фазы II) требуется, прежде чем карта обнаружит сдвиг положения. Чтобы ответить на эти вопросы, нужно распределение длины серии в состоянии OOC и связанные с ним характеристики. Здесь мы сосредоточимся в основном на безусловном распределении длины серии, как это принято в текущей литературе, но, как отмечалось ранее, возможен и условный анализ, который, возможно, даёт больше информации.

Как мы уже видели, безусловное распределение длины серии карты Shewhart-Prec не является геометрическим. В уравнении \eqref{eq:4.28} приведена формула функции вероятности длины серии. Её можно вычислить, чтобы ответить на подобные вопросы. Отметим, однако, что для вычисления распределения длины серии в случаях, отличных от ситуации IC, например при сдвиге, нужно задать функцию \(GF^{-1}\left( u_{(b:m)} \right)\). Пример приведён ниже.

Двумя (очень) популярными характеристиками распределения длины серии являются средняя длина серии, или среднее распределения длины серии, и вероятность сигнала или, эквивалентно, вероятность отсутствия сигнала. Поэтому мы сосредоточимся на этих характеристиках как мерах эффективности карты предшествования. Следует отметить, что практическое соображение при реализации контрольной карты состоит в том, что как для двусторонней, так и для односторонней контрольных карт должны выполняться определённые условия на константы карты \(a\) и \(b\), а также на \(m\), \(n\) и \(j\), чтобы \(ARL_{IC}\) была конечной. Эти условия приведены в Chakraborti, van der Laan and Van de Wiel (2004):

\begin{enumerate}
	\item для двусторонней контрольной карты:
	\[
		ARL_{IC} < \infty \iff (a - j)(n - j + 1) + j(m - b + 1) > 0;
	\]
	и
	\item для односторонней контрольной карты:
	\[
		ARL_{IC} < \infty \iff (m - b) - (n - j) > 0.
	\]
\end{enumerate}

Например, если мы намерены построить двустороннюю карту Med с тестовыми выборками объёма \(n = 5\), используя \(m = 50\) эталонных наблюдений и при заданной FAR 0.0027, то таблица \ref{tab:4.18} даёт значения \(a = 1\) и \(b = 50\) при точной FAR, равной 0.00076. Однако для этих значений получаем \((a - j)(n - j + 1) + j(m - b + 1) = -3\), что меньше нуля. Поэтому из таблицы 1 в Chakraborti, van der Laan and Van de Wiel (2004) находим, что \(ARL_{IC} = \infty\).

При изучении эффективности любой контрольной карты следует сначала изучить её свойства робастности при возможных нарушениях предположений, в рамках которых карта проектировалась. Для этого робастность контрольной карты в состоянии IC наиболее важна, поскольку без неё свойства обнаружения сдвига карты, какими бы привлекательными они ни были, оказались бы сомнительными. Можно провести аналогию с размером и мощностью критерия проверки гипотез. Если критерий не сохраняет свой размер, то есть не является робастным к ошибке первого рода относительно лежащих в основе предположений, его свойства мощности мало что значат. Свободные от распределения карты предшествования по определению робастны в состоянии IC в том смысле, что их свойства в состоянии IC остаются одинаковыми для всех непрерывных распределений. Теперь рассмотрим карту Med с этой точки зрения.

\paragraph{Робастность двусторонней карты Med в состоянии IC.} Как отмечалось ранее, одно из наиболее привлекательных свойств свободных от распределения контрольных карт --- устойчивость распределения длины серии в состоянии IC. В противоположность этому напомним, что средняя длина серии в состоянии IC параметрических контрольных карт, таких как стандартная параметрическая карта \(\overline{X}\) Шухарта, карты CUSUM и EWMA, может (иногда сильно) изменяться при изменениях лежащего в основе распределения процесса. Фактически характеристики эффективности в состоянии IC стандартных параметрических карт могут существенно изменяться при отклонениях от параметрического распределения (например, нормального), в рамках которого проектировалась карта.

Таблица \ref{tab:4.19}, например, показывает значения \(ARL_{IC}\) двусторонней карты Med. Для сравнения мы также включаем соответствующие результаты для двусторонних (параметрических) карт CUSUM и EWMA, все они основаны на нормальном распределении со средним \(\mu\) и дисперсией \(\sigma^2\). Эти результаты взяты из Chakraborti, van der Laan and Van de Wiel (2004). Константы карты CUSUM \(h\) и \(k\), а также константы карты EWMA \(\lambda\) и \(L\) выбирались так, чтобы карта одинаково хорошо обнаруживала либо малый сдвиг, то есть \(\tfrac{1}{2}\sigma\) в данном случае, либо большой сдвиг процесса, то есть \(2\sigma\). Читателю рекомендуется обратиться к разделам 3.2.2 и 3.2.3 для подробного обсуждения выбора параметров проектирования \((h, k)\) и \((\lambda, L)\) для карт CUSUM и EWMA соответственно.

В исследовании робастности Chakraborti, van der Laan and Van de Wiel (2004) включили распределения различной формы: нормальные, не нормальные, но близкие к нормальным с более тяжёлыми и более лёгкими хвостами, и асимметричные распределения. Для не нормальных, но близких к нормальным распределений использовались распределение Стьюдента \(t\) и распределение Лапласа (также известное как двойное экспоненциальное распределение). Эти два распределения отражают влияние тяжёлых хвостов (более тяжёлых, чем у нормального) на \(ARL_{IC}\), тогда как равномерное распределение использовалось для изучения влияния лёгких хвостов (более лёгких, чем у нормального). Кроме того, чтобы изучить влияние асимметрии на \(ARL_{IC}\), были включены два различных гамма-распределения. Наконец, чтобы изучить влияние изменения дисперсии процесса, средние длины серии для распределения \(N(0, 1)\) сравнивались со средними длинами серии для распределения \(N(0, 1.1)\).

Результаты таблицы \ref{tab:4.19} поучительны. Мы видим, что \(ARL_{IC}\) карты CUSUM особенно чувствительна к отклонениям от предположения о нормальности и к изменению дисперсии процесса, особенно когда карта CUSUM проектируется специально для обнаружения большого сдвига положения процесса. Напротив, карта EWMA относительно робастна к отклонениям от нормальности, особенно когда весовая константа \(\lambda\) мала. Однако карта EWMA не особенно робастна к изменению дисперсии распределения процесса. Во всех же случаях \(ARL_{IC}\) карты Shewhart-Prec остаётся постоянной и равна 501.89 для заданного номинального значения \(ARL_{IC}\), равного 500. В последнее время ряд исследований рассматривал вопрос робастности в состоянии IC многих параметрических контрольных карт и пришёл к выводу, что предположение о распределении может оказаться решающим при применении этих карт: например, при нарушении предположения о распределении по этим картам наблюдается гораздо больше ложных тревог. Это не согласуется с заявленной FAR и может, пожалуй, разрушить эффективность значительной части схемы SPC. В результате возник новый интерес к свободным от распределения контрольным картам. Далее мы рассмотрим сравнения эффективности в состоянии OOC.

\begin{table}[!ht]
	\centering
	\caption{Значения \(ARL_{IC}\) для нескольких распределений.}
	\label{tab:4.19}
	\small
	\begin{tabular}{@{}lccccc@{}}
		\toprule
		& Med & \multicolumn{2}{c}{CUSUM} & \multicolumn{2}{c}{EWMA} \\
		\cmidrule(lr){2-2}\cmidrule(lr){3-4}\cmidrule(lr){5-6}
		& \(m = 1000\) & Большой & Малый & Большой & Малый \\
		& \(n = 5\), \(j = 3\) & сдвиг & сдвиг & сдвиг & сдвиг \\
		Распределение & \(P_0 = 0.0078\) & \(h = 4.61\) & \(h = 0.86\) & \(L = 2.81\) & \(L = 3.05\) \\
		& \(a = 48\), \(b = 953\) & \(k = 0.56\) & \(k = 2.24\) & \(\lambda = 0.1\) & \(\lambda = 0.4\) \\
		& & \(n = 5\) & \(n = 5\) & \(n = 5\) & \(n = 5\) \\
		\midrule
		\(N(0, 1)\) & 501.89 & 500.00 & 500.00 & 500.00 & 500.00 \\
		\(Laplace(0, 1)\) & 501.89 & 390.64 & 196.18 & 418.77 & 289.82 \\
		\(Uniform(0, 1)\) & 501.89 & 549.16 & 1820.13 & 494.73 & 703.72 \\
		\(N(0, 1.1)\) & 501.89 & 206.38 & 196.69 & 214.49 & 199.79 \\
		\(Gamma(1, 1)\) & 501.89 & 310.74 & 121.53 & 404.27 & 207.62 \\
		\(Gamma(4, 1)\) & 501.89 & 462.43 & 345.68 & 472.33 & 370.20 \\
		\(t(4)\) & 501.89 & 266.86 & 131.90 & 380.98 & 190.39 \\
		\(t(40)\) & 501.89 & 456.76 & 413.44 & 470.87 & 380.63 \\
		\bottomrule
	\end{tabular}
\end{table}

\paragraph{Средняя длина серии двусторонней карты Med в состоянии OOC.} Для сравнения средней длины серии двусторонней карты предшествования в состоянии OOC с таковой для стандартной параметрической карты \(\overline{X}\) Шухарта, карты CUSUM и карты EWMA в таблицах \ref{tab:4.20a}--\ref{tab:4.20c} приведены средние длины серии для распределений \(N(0, 1)\), \(Gamma(1, 1)\) и \(t(4)\). Эти результаты снова взяты из Chakraborti, van der Laan and Van de Wiel (2004). Из таблиц \ref{tab:4.20a} и \ref{tab:4.20b} видно, что карты CUSUM и EWMA лучше карты предшествования при обнаружении (особенно малых) сдвигов в нормальном распределении и распределении \(Gamma(1, 1)\). Однако для нормального распределения различия между картой предшествования и параметрическими картами CUSUM и EWMA не очень велики для сдвигов, превышающих, скажем, \(1.5\sigma\). Тем не менее это неудивительно, учитывая, что карта Shewhart-Prec является картой Шухарта.

Из таблицы \ref{tab:4.20c} видно, что для распределения \(t(4)\), то есть для симметричного распределения с более тяжёлыми хвостами, чем у стандартного нормального распределения, малые сдвиги лучше обнаруживаются картой CUSUM (с малым значением константы карты \(k\)), а также картой EWMA (с малым значением весовой константы \(\lambda\)), но большие сдвиги снова обнаруживаются раньше картой предшествования. Отметим, однако, что для этих сравнений используется эталонная выборка из 1000 наблюдений. Вообще говоря, верно, что непараметрические методы требуют больше данных, но, разумеется, они не требуют предположения о параметрическом распределении модели. Выбор сводится к конкретным применениям, доступным ресурсам и приоритетам.

Хотя значительная часть нашего обсуждения была посвящена двусторонним контрольным картам в соответствии с традиционной практикой, отметим, что иногда требуются односторонние контрольные карты, которые полезны в ситуациях, когда представляет интерес только восходящий или только нисходящий сдвиг определённого параметра процесса. Например, мы можем контролировать прочность на разрыв троса, используемого для поддержания тяжёлого груза, или контролировать (остаточный или будущий) срок службы лампы накаливания. Многие двусторонние контрольные карты, как параметрические, так и непараметрические, можно модифицировать под эту ситуацию. Однако карты предшествования, по-видимому, особенно подходят для этого, поскольку они основаны на упорядоченных значениях. Мы кратко затронем это в следующем разделе.

\begin{table}[!ht]
	\centering
	\caption{Значения средней длины серии для распределения \(N(0, 1)\).}
	\label{tab:4.20a}
	\small
	\begin{tabular}{@{}lccccc@{}}
		\toprule
		& Med & \multicolumn{2}{c}{CUSUM} & \multicolumn{2}{c}{EWMA} \\
		\cmidrule(lr){2-2}\cmidrule(lr){3-4}\cmidrule(lr){5-6}
		& \(m = 1000\) & Большой & Малый & Большой & Малый \\
		Сдвиг & \(n = 5\), \(j = 3\) & \(h = 4.61\) & \(h = 0.86\) & \(L = 2.81\) & \(L = 3.05\) \\
		\(\sigma\) & \(FAR = 0.0078\) & \(k = 0.56\) & \(k = 2.24\) & \(\lambda = 0.1\) & \(\lambda = 0.4\) \\
		& \(a = 48\), \(b = 953\) & \(n = 5\) & \(n = 5\) & \(n = 5\) & \(n = 5\) \\
		\midrule
		0.00 & 501.89 & 500.00 & 500.00 & 500.00 & 500.00 \\
		0.25 & 240.93 & 33.35 & 162.72 & 25.75 & 56.05 \\
		0.50 & 71.70 & 8.78 & 36.09 & 8.86 & 11.00 \\
		0.75 & 24.22 & 4.86 & 10.60 & 5.33 & 4.76 \\
		1.00 & 9.79 & 3.40 & 4.21 & 3.87 & 2.97 \\
		1.50 & 2.70 & 2.22 & 1.56 & 2.57 & 1.77 \\
		2.00 & 1.37 & 1.77 & 1.09 & 2.04 & 1.25 \\
		2.50 & 1.07 & 1.34 & 1.01 & 1.81 & 1.04 \\
		3.00 & 1.01 & 1.06 & 1.00 & 1.40 & 1.00 \\
		4.00 & 1.00 & 1.00 & 1.00 & 1.01 & 1.00 \\
		\bottomrule
	\end{tabular}
\end{table}

\begin{table}[!ht]
	\centering
	\caption{Значения средней длины серии для распределения \(Gamma(1, 1)\).}
	\label{tab:4.20b}
	\small
	\begin{tabular}{@{}lccccc@{}}
		\toprule
		& Med & \multicolumn{2}{c}{CUSUM} & \multicolumn{2}{c}{EWMA} \\
		\cmidrule(lr){2-2}\cmidrule(lr){3-4}\cmidrule(lr){5-6}
		& \(m = 1000\) & Большой & Малый & Большой & Малый \\
		Сдвиг & \(n = 5\), \(j = 3\) & \(h = 5.18\) & \(h = 1.75\) & \(L = 2.81\) & \(L = 3.05\) \\
		\(\sigma\) & \(FAR = 0.0078\) & \(k = 0.56\) & \(k = 2.24\) & \(\lambda = 0.1\) & \(\lambda = 0.4\) \\
		& \(a = 48\), \(b = 953\) & \(n = 5\) & \(n = 5\) & \(n = 5\) & \(n = 5\) \\
		\midrule
		0.00 & 501.89 & 496.43 & 498.22 & 509.76 & 505.68 \\
		0.25 & 439.01 & 41.40 & 182.85 & 25.18 & 80.83 \\
		0.50 & 256.45 & 10.28 & 96.94 & 7.66 & 18.73 \\
		0.75 & 125.31 & 5.44 & 24.50 & 3.91 & 9.98 \\
		1.00 & 61.87 & 3.75 & 9.05 & 2.47 & 3.70 \\
		1.50 & 15.81 & 2.43 & 2.32 & 1.38 & 1.76 \\
		2.00 & 4.83 & 1.90 & 1.36 & 1.03 & 1.20 \\
		2.50 & 1.63 & 1.61 & 1.02 & 1.00 & 1.00 \\
		3.00 & 1.02 & 1.15 & 1.00 & 1.00 & 1.00 \\
		4.00 & 1.00 & 1.00 & 1.00 & 1.00 & 1.00 \\
		\bottomrule
	\end{tabular}
\end{table}

\begin{table}[!ht]
	\centering
	\caption{Значения средней длины серии для распределения \(t(4)\).}
	\label{tab:4.20c}
	\small
	\begin{tabular}{@{}lccccc@{}}
		\toprule
		& Med & \multicolumn{2}{c}{CUSUM} & \multicolumn{2}{c}{EWMA} \\
		\cmidrule(lr){2-2}\cmidrule(lr){3-4}\cmidrule(lr){5-6}
		& \(m = 1000\) & Большой & Малый & Большой & Малый \\
		Сдвиг & \(n = 5\), \(j = 3\) & \(h = 5.38\) & \(h = 1.81\) & \(L = 3.04\) & \(L = 3.65\) \\
		\(\sigma\) & \(FAR = 0.0078\) & \(k = 0.56\) & \(k = 2.24\) & \(\lambda = 0.1\) & \(\lambda = 0.4\) \\
		& \(a = 48\), \(b = 953\) & \(n = 5\) & \(n = 5\) & \(n = 5\) & \(n = 5\) \\
		\midrule
		0.00 & 501.89 & 495.23 & 499.47 & 507.91 & 504.44 \\
		0.25 & 300.09 & 41.88 & 340.58 & 28.28 & 128.87 \\
		0.50 & 32.70 & 10.24 & 136.21 & 7.86 & 22.45 \\
		0.75 & 11.16 & 5.54 & 35.50 & 3.98 & 7.11 \\
		1.00 & 2.26 & 3.86 & 9.84 & 2.47 & 3.77 \\
		1.50 & 1.16 & 2.49 & 2.31 & 1.38 & 1.20 \\
		2.00 & 1.02 & 1.98 & 1.34 & 1.06 & 1.03 \\
		2.50 & 1.00 & 1.63 & 1.06 & 1.01 & 1.00 \\
		3.00 & 1.00 & 1.20 & 1.01 & 1.00 & 1.00 \\
		4.00 & 1.00 & 1.01 & 1.00 & 1.00 & 1.00 \\
		\bottomrule
	\end{tabular}
\end{table}

\paragraph{Замечание 4.5. Односторонние контрольные карты Shewhart-Prec.} В первом сценарии, если, например, прочность троса на разрыв снижается под действием веса, трос может просто оборваться, поскольку не сможет выдержать или удержать груз, и впоследствии это может стать опасностью. Однако, напротив, если прочность троса на разрыв возрастает, то есть сдвигается вверх, это будет преимуществом, поскольку тогда трос сможет выдержать более тяжёлый груз, чем изначально требовалось. Аналогично, если средний срок службы определённого вида лампы накаливания меньше заявленного среднего срока службы, больше клиентов могут решить не покупать эту марку, и, следовательно, это может привести к потере выручки для производителя. Поэтому разумно защищаться от нисходящего сдвига среднего срока службы лампы накаливания. В обоих этих сценариях нас интересует только обнаружение нисходящего сдвига интересующего параметра процесса, в данном случае параметра положения, и поэтому желательна нижняя односторонняя контрольная карта. В качестве альтернативы в ситуации, когда только восходящий сдвиг может иметь реальное значение, потребуется верхняя односторонняя карта. Читателю рекомендуется обратиться к разделу 3.2.1 для обсуждения того, когда уместна нижняя односторонняя, верхняя односторонняя или двусторонняя карта. Вопрос о том, когда нужна односторонняя или двусторонняя контрольная карта, аналогичен вопросу о выборе между односторонним и двусторонним доверительным интервалом.

Корректировки, необходимые для перехода от двусторонней карты предшествования к односторонней, относительно просты. Например, если представляет интерес обнаружение только восходящего сдвига, то есть того, увеличился ли параметр положения (или сдвинулся вправо), мы используем только \(\widehat{UCL} = X_{(b:m)}\) в качестве оцениваемой UCL без LCL. Однако если представляет интерес обнаружение только нисходящего сдвига, то есть того, уменьшился ли параметр положения (или сдвинулся влево), мы используем только \(\widehat{LCL} = X_{(a:m)}\) в качестве оцениваемой LCL без UCL. Мы предлагаем читателю завершить обсуждение односторонних карт, включая константы для нанесения на карту и вывод свойств карты, следуя приведённым деталям для двусторонних карт. Некоторые дополнительные подробности можно найти в Chakraborti, van der Laan and Van de Wiel (2004).

\subsubsection{Контрольные карты Шухарта на основе статистики критерия Манна--Уитни}
Хотя тесты предшествования являются хорошо известными свободными от распределения тестами и полезны в некоторых ситуациях, критерий Манна--Уитни (обозначаемый MW; см. Mann and Whitney (1947)) --- другой, пожалуй, более популярный двухвыборочный свободный от распределения критерий, используемый на практике как альтернатива классическому двухвыборочному \(t\)-критерию. Известно, что критерий MW (см. Gibbons and Chakraborti, 2010, с.~261) более эффективен, чем критерий предшествования, для распределений с лёгкими хвостами. Литература по критерию MW обширна. Сначала мы опишем статистику MW, а затем обсудим контрольную карту Шухарта, основанную на статистике MW.

\paragraph{Статистика критерия Манна--Уитни.} Свободная от распределения карта Шухарта для положения в Фазе II, основанная на статистиках MW, была предложена Chakraborti и Van de Wiel (2008). В этой постановке, как и в случае критерия предшествования, рассматриваются два непрерывных распределения (распределения \(X\) и \(Y\)) с функциями распределения \(F\) и \(G\) соответственно, где \(F\) представляет функцию распределения эталонного распределения, а \(G\) --- функцию распределения распределения, из которого извлекается тестовая выборка. Из \(F\) извлекается случайная выборка объёма \(m\), то есть \(X_1, X_2, \ldots, X_m\), а из \(G\) --- независимая случайная выборка объёма \(n\), то есть \(Y_1, Y_2, \ldots, Y_n\). Критерий MW основан на относительном положении \(X\)- и \(Y\)-наблюдений в объединённой выборке. Для этого определим \(mn\) индикаторных случайных величин как
\[
	D_{ij} = \begin{cases} 1, & \text{если } X_i > Y_j, \\ 0, & \text{если } X_i < Y_j, \end{cases} \quad i = 1, 2, \ldots, m,\; j = 1, 2, \ldots, n.
\]
Критерий MW основан на статистике \(U\), которая определяется как число случаев, когда наблюдение \(Y\) предшествует наблюдению \(X\) в объединённом упорядоченном расположении \(X\)- и \(Y\)-наблюдений. Это означает, что каждое наблюдение \(X\) сравнивается с каждым из наблюдений \(Y\) попарно, и всякий раз, когда \(X\) больше \(Y\), это сравнение оценивается как 1; в противном случае --- как 0. Статистика MW подсчитывает общее число единиц в этих \(mn\) сравнениях и может быть записана как \(M_{YX} = \sum_{i=1}^{m}\sum_{j=1}^{n} D_{ij}\). Таким образом, \(M_{XY}\) представляет число пар \((X, Y)\), в которых наблюдение \(X\) больше \(Y\). Начнём с примера того, как вычислять статистику MW.

\noindent\textbf{Пример 4.16. Вычисление статистики MW.}

Предположим, что \(m = 3\) и \(n = 3\). Ранжировав набор из \(m + n\) наблюдений совместно от наименьшего к наибольшему, предположим, что мы получили следующую конфигурацию: \(x_1\, y_1\, y_2\, x_2\, x_3\, y_3\). Исходя из этого, вычисляем \(U\)-критерий MW следующим образом. Сначала отметим, что
\begin{align*}
	D_{11} &= 0, \text{ так как } X_1 < Y_1; & D_{21} &= 1, \text{ так как } X_2 > Y_1; & D_{31} &= 1, \text{ так как } X_3 > Y_1; \\
	D_{12} &= 0, \text{ так как } X_1 < Y_2; & D_{22} &= 1, \text{ так как } X_2 > Y_2; & D_{32} &= 1, \text{ так как } X_3 > Y_2; \\
	D_{13} &= 0, \text{ так как } X_1 < Y_3; & D_{23} &= 0, \text{ так как } X_2 < Y_3; & D_{33} &= 0, \text{ так как } X_3 < Y_3.
\end{align*}
Таким образом, \(M_{YX} = \sum_{i=1}^{m}\sum_{j=1}^{n} D_{ij} = 4\).

\paragraph{Контрольная карта Шухарта на основе статистики Манна--Уитни.} Для контрольной карты Шухарта, основанной на статистике MW (далее обозначаемой карта Shewhart-MW) Chakraborti и Van de Wiel (2008), пусть \(X_1, X_2, \ldots, X_m\) обозначают эталонную выборку Фазы I из процесса в состоянии IC, а \(Y_1, Y_2, \ldots, Y_n\) --- (тестовую) выборку Фазы II объёма \(n\). Статистикой для нанесения на карту служит общее число пар \((X, Y)\), в которых наблюдение \(Y\) больше наблюдения \(X\) (или \(X\) предшествует \(Y\)): \(M_{XY} = \sum_{i=1}^{m}\sum_{j=1}^{n} I\left( Y_j > X_i \right)\), где \(I(\cdot)\) обозначает обычную индикаторную функцию. То, что мы считаем \(Y\)-наблюдения, большие \(X\)-наблюдений (а не наоборот, как в предыдущем описании), сделано просто для удобства. Отметим, что \(M_{XY}\) принимает значения от 0 до \(mn\), и большие (меньшие) значения \(M_{XY}\) указывают на положительный (отрицательный) сдвиг положения. Наименьшее значение 0 достигается, когда все \(Y\) меньше \(X\), а наибольшее значение \(mn\) --- когда все \(Y\) больше \(X\). Если \(Y\)-наблюдения стохастически больше (меньше) \(X\)-наблюдений, то \(M_{XY}\) будет значимо большим (малым), и поэтому карта Shewhart-MW подаёт сигнал о ситуации OOC, когда \(M_{XY} > b_{mn}\) или \(M_{XY} < a_{mn}\), где \(b_{mn}\) и \(a_{mn}\) --- соответственно UCL и LCL. В модели положения \(Y\)-наблюдения стохастически больше (меньше) \(X\)-наблюдений в соответствии с условием \(\theta_i < \theta_1\) (или \(\theta_i > \theta_1\)) для соответствующего параметра положения. Иначе говоря, процесс находится в состоянии OOC, когда медиана \(Y\) уменьшилась (или увеличилась) по сравнению с медианой \(X\).

Хотя критерий MW --- хорошо известный свободный от распределения критерий, более привычным двухвыборочным свободным от распределения критерием для положения является критерий суммы рангов Уилкоксона (WRS), предложенный Wilcoxon (1945) и опубликованный раньше статьи Mann и Whitney (1947). Статистика WRS (\(W_{XY}\)) является одним из представителей класса так называемых «линейных ранговых статистик» (см. Gibbons and Chakraborti, 2010, с.~275). Статистика WRS --- сумма рангов \(Y\)-наблюдений в объединённой выборке \(X\)- и \(Y\)-наблюдений. Можно показать, что \(W_{XY} = M_{XY} + m(m + 1)/2\) (см., например, Gibbons and Chakraborti, 2010, с.~293), так что оба критерия эквивалентны.

Таким образом, двустороннюю карту Shewhart-MW Chakraborti и Van de Wiel (2008) можно также описать через статистику WRS \(W_{XY}\). Для этого отметим, что при отсутствии совпадений среднее и дисперсия \(W_{XY}\) в состоянии IC могут быть найдены (см., например, Gibbons and Chakraborti, 2010, с.~291) как
\begin{equation}
	E(W_{XY}) = \frac{m(m + n + 1)}{2} \quad \text{и} \quad VAR(W_{XY}) = \frac{mn(m + n + 1)}{12}.
	\label{eq:4.45}
\end{equation}
Следовательно, карта Shewhart-Wilcoxon подаёт сигнал OOC, когда
\[
	W_{XY} > b_{mn} + m(m + 1)/2 \quad \text{или} \quad W_{XY} < a_{mn} + \frac{m(m + 1)}{2}.
\]
Проиллюстрируем вычисление статистик MW и WRS на небольшом наборе данных.

\noindent\textbf{Пример 4.17. Вычисление статистик MW и WRS.}

Используя те же данные, что и в предыдущем примере, предположим, что \(m = 3\) и \(n = 3\) и совпадений нет. Ранжировав набор из \(m + n = 6\) наблюдений совместно от наименьшего к наибольшему, напомним, что мы получаем следующую упорядоченную объединённую выборку: \(x_1\, y_1\, y_2\, x_2\, x_3\, y_3\). Для этих данных легко видеть, что \(M_{XY} = \sum_{i=1}^{m}\sum_{j=1}^{n} I\left( Y_j > X_i \right) = 5\). Отметим, что ранее мы нашли \(M_{YX} = 4\) и \(M_{YX} + M_{XY} = mn = 9\). Следовательно, \(W_{XY} = \sum_{i=1}^{m}\sum_{j=1}^{n} I\left( Y_j > X_i \right) + \frac{m(m + 1)}{2} = 5 + 6 = 11\).

Далее мы изучаем распределение длины серии карты Shewhart-MW.

\paragraph{Методы вычисления распределения длины серии.} Следуя методологии Chakraborti и Van de Wiel (2008), пусть \(p(x)\) обозначает вероятность сигнала при любой тестовой (Фазе II) выборке при заданной эталонной выборке \(X = (X_1, X_2, \ldots, X_m) = (x_1, x_2, \ldots, x_m) = x\). Таким образом, \(p(x)\) является функцией функции распределения \(G\) и записывается как \(p_G(x)\). Условное распределение длины серии является геометрическим и задаётся как
\begin{equation}
	P_C\left( N = k \mid X = x \right) = \left( 1 - p_G(x) \right)^{k-1}\left( p_G(x) \right), \quad k = 1, 2, \ldots
	\label{eq:4.46}
\end{equation}
Следовательно, условная ARL равна
\begin{equation}
	ARL_C = E\left( N \mid X = x \right) = \frac{1}{p_G(x)}.
	\label{eq:4.47}
\end{equation}
Отсюда следует, что безусловная ARL равна
\begin{equation}
	ARL_U = \int_{-\infty}^{\infty}\cdots\int_{-\infty}^{\infty} \frac{1}{p_G(x)}\, dF(x_1)\cdots dF(x_m).
	\label{eq:4.48}
\end{equation}

\paragraph{Безусловная ARL карты Shewhart-MW в состоянии IC.} Безусловная ARL в состоянии IC находится подстановкой \(G = F\) в уравнение \eqref{eq:4.48}. Таким образом,
\begin{equation}
	ARL_{U, IC} = \int_{-\infty}^{\infty}\cdots\int_{-\infty}^{\infty} \frac{1}{p_F(x)}\, dF(x_1)\cdots dF(x_m),
	\label{eq:4.49}
\end{equation}
где \(p(x)\) теперь является функцией \(F\), поскольку \(G = F\). Используя PIT, уравнение \eqref{eq:4.49} можно записать как
\begin{equation}
	ARL_{U, IC} = \int_{0}^{1}\cdots\int_{0}^{1} \frac{1}{p_U(u)}\, du_1\cdots du_m,
	\label{eq:4.50}
\end{equation}
где \(u_i = F(x_i)\), \(i = 1, 2, \ldots, m\), \(u = (u_1, u_2, \ldots, u_m)\), а индекс \(U\) у \(p_U\) обозначает тот факт, что вероятность сигнала \(p(x)\) вычислена при общем распределении \(F = G\) на основе \(u_i\), которые являются реализациями \(U(0, 1)\)-распределения.

\paragraph{ARL в состоянии OOC.} В случае OOC безусловная ARL в состоянии OOC равна
\begin{equation}
	ARL_{U, OOC} = \int_{-\infty}^{\infty}\cdots\int_{-\infty}^{\infty} \frac{1}{p_G(x)}\, dF(x_2)\cdots dF(x_m),
	\label{eq:4.51}
\end{equation}
где \(\delta = \theta_i - \theta_1\) представляет разность (сдвиг) между параметрами положения \(G\) и \(F\). Таким образом, мы можем найти как \(ARL_{U, IC}\), так и \(ARL_{U, OOC}\), используя метод CUC.

Нам нужно вычислить \(ARL_{U, IC}\) и \(ARL_{U, OOC}\) соответственно, чтобы реализовать карту Shewhart-MW и оценить её эффективность. Для этого отметим, что у нас нет точной компактной формулы для ARL. В отличие от карты Shewhart-Prec, приведённые выше выражения для ARL представляют собой \(m\)-мерные интегралы. Соответственно, вычисление безусловной ARL аналитически по этим формулам трудно и требует много времени. В качестве практической альтернативы, как уже неоднократно упоминалось, мы снова рекомендуем использовать симуляции для оценивания и анализа ARL или любой другой характеристики распределения длины серии. Общие шаги симуляции характеристик длины серии для любой контрольной карты описаны ниже. Отметим, что эти шаги применимы к любой карте, параметрической или непараметрической; разница лишь в том, что конкретная статистика для нанесения на карту и её вычисление будут меняться.

\paragraph{Симуляционный подход.} Для получения характеристик длины серии карты Shewhart-MW можно использовать следующий алгоритм симуляции. Соответствующая программа SAS приведена в приложении B, программа SAS номер 8, а вывод показан ниже на рисунке \ref{fig-4.24}.

\noindent\textbf{Шаг 1.} Задаём объём эталонной выборки Фазы I \(m \ge 1\), объём тестовой выборки Фазы II \(n \ge 1\), нужную номинальную \(ARL_{IC}\) и величину сдвига (\(\delta\)), который нужно обнаружить. Затем задаём два распределения для генерации выборок Фазы I и Фазы II соответственно. В случае IC оба распределения совпадают (говорят, что сдвиг равен нулю (\(\delta = 0\))). В случае OOC распределение выборки Фазы II имеет тот же вид, что и у выборки Фазы I, но со сдвигом положения в единицах стандартного отклонения совокупности. Выборка Фазы I, например, может быть извлечена из нормального распределения со средним 0 и стандартным отклонением 1, тогда как в случае OOC выборки Фазы II могут быть извлечены из нормального распределения с некоторым средним \(\delta \ne 0\) и стандартным отклонением 1 (говорят, что сдвиг среднего равен \(\delta \ne 0\)). Для двух лежащих в основе распределений можно использовать любое непрерывное распределение, симметричное или асимметричное.

\noindent\textbf{Шаг 2.} Вычисляем контрольные границы, то есть константу для нанесения на карту \(b_{mn}\); \(a_{mn}\) можно получить из соотношения \(a_{mn} = mn - b_{mn}\). Входными являются \(m\), \(n\) и номинальная \(ARL_{IC}\). Обычно для получения константы \(b_{mn}\) используется поисковый алгоритм (с применением программного пакета, такого как SAS или R).

\noindent\textbf{Шаг 3.} Вычисляем статистику для нанесения на карту MW (\(M_{XY}\)) для выборки Фазы II, сгенерированной на шаге 1, и эталонной выборки.

\noindent\textbf{Шаг 4.} Если, с одной стороны, статистика для нанесения на карту \(M_{XY}\) попадает на контрольные границы или за них, записывается наблюдение 1 как наблюдение из распределения длины серии. Если, с другой стороны, статистика для нанесения на карту \(M_{XY}\) попадает внутрь контрольных границ, шаги 1 и 2 повторяются с той же эталонной выборкой, но с заново сгенерированной выборкой Фазы II каждый раз, пока статистика для нанесения на карту \(M_{XY}\) впервые не попадёт на контрольные границы или за них, так что будет зафиксирован сигнал. В этот момент число выборок Фазы II, потребовавшихся для получения этого первого сигнала, записывается как наблюдение из распределения длины серии.

\noindent\textbf{Шаг 5.} Повторяем шаги 1--4 в общей сложности, например, 100\,000 раз. Это создаёт 100\,000 наблюдений из распределения длины серии контрольной карты MW (или любой другой карты Фазы II).

\noindent\textbf{Шаг 6.} Наблюдения, полученные на шаге 5, можно проанализировать для вычисления оценок характеристик распределения длины серии, таких как среднее, медиана, стандартное отклонение и т.~д.

Для программирования шагов симуляции можно использовать программный пакет, такой как SAS или R. Отметим, что при \(\delta = 0\) мы получаем характеристики длины серии в состоянии IC, а при \(\delta \ne 0\) --- характеристики длины серии в состоянии OOC.

\begin{figure}[!ht]
	\centering
	\includegraphics[width=0.75\linewidth]{"./images/fig-4.24.png"}
	\caption[]{Вывод SAS, где параметры проектирования (\(m = 500\), \(n = 5\), \(b_{mn} = 2172\)) карты Shewhart-MW были выбраны для номинальной \(ARL_{IC} = 500\).}
	\label{fig-4.24}
\end{figure}

\subsection{Контрольные карты CUSUM}
\subsubsection{Контрольные карты CUSUM на основе статистики превышения}
Хотя карта Shewhart-Prec является картой Шухарта и поэтому, как хорошо известно, эффективна при обнаружении больших и более резких сдвигов, в некоторых ситуациях могут оказаться желательными карты, более чувствительные к малым и более устойчивым сдвигам. Как обсуждалось в главе 3, в литературе по параметрическим контрольным картам рассматриваются два класса таких карт, а именно карты CUSUM и EWMA. Поэтому естественно рассмотреть непараметрические аналоги этих карт в случае U, и здесь мы рассматриваем две такие карты. Хотя можно рассмотреть непараметрическую карту CUSUM, основанную на статистике предшествования, для удобства мы, следуя литературе, рассматриваем использование некоторых связанных статистик, называемых статистиками превышения (exceedance statistics). Статистика превышения --- это число наблюдений одной из двух выборок (выборки Фазы I и тестовой выборки), которые превышают \(r\)-ю порядковую статистику другой выборки. Таким образом, статистики предшествования и превышения связаны линейно. Критерии, основанные на статистиках предшествования, называются критериями предшествования, а критерии, основанные на статистиках превышения, называются критериями превышения. Последние оказались полезными в ряде приложений, включая контроль качества и исследования надёжности по данным о наработке на отказ, где вероятность превышения может быть связана с так называемым «гарантийным сроком» изделия (см., например, Chakraborti and van der Laan, 2000).

Как и ранее, предположим, что имеется эталонная выборка Фазы I \(X_1, X_2, \ldots, X_m\) (\(m > 1\)) из процесса в состоянии IC с функцией распределения \(F(x)\). Пусть \(Y_{j1}, Y_{j2}, \ldots, Y_{jn}\) (\(n \ge 1\)), \(j = 1, 2, \ldots\), обозначают \(j\)-ю тестовую (Фазу II) выборку объёма \(n\) из функции распределения \(G(y)\). Обе функции \(F\) и \(G\) --- неизвестные непрерывные функции распределения, и процесс находится в состоянии IC, когда \(F = G\). Для обнаружения изменения положения предположим, что функции распределения равны \(F(t) = H(t - \theta_1)\) и \(G(t) = H(t - \theta_2)\) соответственно, где \(H\) --- некоторая неизвестная непрерывная функция распределения, а \(\theta_1\) и \(\theta_2\) представляют соответствующие параметры положения. Пусть \(U_{j,r}\) обозначает число наблюдений \(Y\) в \(j\)-й выборке Фазы II, которые превышают \(X_{(r)}\), \(r\)-е упорядоченное наблюдение в эталонной выборке Фазы I. Статистика \(U_{j,r}\) называется статистикой превышения, а вероятность \(p_r = P\left( Y > X_{(r)} \mid X_{(r)} \right)\) называется вероятностью превышения. Отметим, что эта вероятность превышения является условной вероятностью того, что наблюдение \(Y\) превышает \(X_{(r)}\), при заданном значении \(X_{(r)}\).

Далее рассматривается карта CUSUM, основанная на статистике превышения (обозначаемая CUSUM-EX). Это обсуждение основано на работе Mukherjee, Graham and Chakraborti (2013), к которой можно обратиться за подробностями.

Верхняя односторонняя карта CUSUM-EX определяется через статистику для нанесения на карту
\begin{equation}
	C_j^+ = \max\left[ 0, C_{j-1}^+ + \left( U_{j,r} - nd \right) - k \right], \quad j = 1, 2, 3, \ldots
	\label{eq:4.52}
\end{equation}
с начальным значением \(C_0^+ = 0\). Она подаёт сигнал при первом \(j\), для которого \(C_j^+ \ge h\).

Аналогично для нижней односторонней карты CUSUM-EX статистика для нанесения на карту задаётся как
\begin{equation}
	C_j^- = \min\left[ 0, C_{j-1}^- + \left( U_{j,r} - nd \right) + k \right], \quad j = 1, 2, 3, \ldots
	\label{eq:4.53}
\end{equation}
с начальным значением \(C_0^- = 0\). Она подаёт сигнал при первом \(j\), для которого \(C_j^- \le -h\).

Ключевая идея предлагаемых карт состоит в том, что при заданном \(X_{(r)} = x_{(r)}\) каждая \(U_{j,r}\) имеет биномиальное распределение с параметрами \((n, p_r)\). Следовательно, условно на \(X_{(r)}\) и следуя идее существующих карт CUSUM для биномиальной вероятности, можно построить карту CUSUM биномиального типа на основе величин \(U_{j,r}\) для мониторинга положения процесса (через вероятности превышения). Теперь из свойств биномиального распределения \(E\left( U_{j,r} \mid X_{(r)} \right) = np_r\) для любого \(j = 1, 2, 3, \ldots\). Поскольку условная вероятность \(p_r\) неизвестна и зависит от лежащих в основе распределений, в литературе было предложено заменить \(p_r\) её безусловным значением в состоянии IC, равным \(P\left( Y > X_{(r)} \right) = \frac{m - r + 1}{m + 1} = d\).

Соответствующая двусторонняя карта CUSUM-EX подаёт сигнал при первом \(j\), для которого подаёт сигнал любая из односторонних карт, то есть всякий раз, когда \(C_j^+ \ge h\) или \(C_j^- \le -h\). Отметим, что здесь интервал решения \(h > 0\) и референсное значение \(k > 0\) являются параметрами процедуры контрольной карты.

Если взять \(X_{(r)}\) в качестве медианы эталонной выборки, то получим \(d = 0.5\). Подставив \(d = 0.5\) в уравнения \eqref{eq:4.52} и \eqref{eq:4.53}, получаем, что верхняя односторонняя карта CUSUM-EX, основанная на медиане эталонной выборки, задаётся статистикой для нанесения на карту
\begin{equation}
	C_j^+ = \max\left[ 0, C_{j-1}^+ + \left( U_{j,r} - n/2 \right) - k \right], \quad j = 1, 2, 3, \ldots
	\label{eq:4.54}
\end{equation}
а нижняя односторонняя карта CUSUM-EX, основанная на медиане эталонной выборки, задаётся статистикой для нанесения на карту
\begin{equation}
	C_j^- = \min\left[ 0, C_{j-1}^- + \left( U_{j,r} - n/2 \right) + k \right], \quad j = 1, 2, 3, \ldots
	\label{eq:4.55}
\end{equation}
Опять же, соответствующая двусторонняя симметричная карта CUSUM-EX подаёт сигнал при первом \(j\), для которого подаёт сигнал любая из односторонних карт, то есть всякий раз, когда \(C_j^+ \ge h\) или \(C_j^- \le -h\). Карту CUSUM-EX, основанную на медиане эталонной выборки, в дальнейшем будем называть медианной картой CUSUM-EX.

\paragraph{Замечание 4.6. Реализация карты: выбор параметров проектирования.} Реализация карты CUSUM-EX требует задания следующих величин: (i) \(m\) --- объёма эталонной выборки Фазы I в состоянии IC; (ii) \(n\) --- объёма каждой тестовой выборки Фазы II (объёма подгруппы); (iii) нужной \(ARL_{IC}\); (iv) \(k\) --- референсного значения, по существу рационального числа; и (v) \(h\) --- интервала решения/UCL, зависящего от \(m\), \(n\), \(k\) и нужной \(ARL_{IC}\). Экспериментатор сам задаёт параметры \(m\), \(n\) и нужную \(ARL_{IC}\) в конкретной ситуации. Выбор объёма эталонной выборки Фазы I в состоянии IC может быть весьма значимым; он подробно обсуждается в приложении B работы Mukherjee, Graham and Chakraborti (2013). Параметры проектирования \((k, h)\) выбираются так, чтобы карта имела заданную номинальную \(ARL_{IC}\) и была способна обнаружить сдвиг, особенно малый, как можно скорее. Первый шаг --- выбрать \(k\). В разделе 3.2.2 мы рассматривали параметрическую карту CUSUM для отдельных данных (\(n = 1\)) без эталонной выборки. Мы изучили влияние \(k\) (читателю рекомендуется обратиться к рисунку 2.5 и его обсуждению) и пришли к выводу, что когда имеется мало априорной информации или она отсутствует относительно величины сдвига, меньший выбор \(k\) является наиболее безопасным (чтобы защититься от ненужных задержек обнаружения). Mukherjee, Graham and Chakraborti (2013) приходят к аналогичным выводам для карты CUSUM-EX. Следующий шаг --- выбрать \(h\) в сочетании с выбранным \(k\) так, чтобы была достигнута нужная номинальная \(ARL_{IC}\). Пример приведён ниже.

\paragraph{Методы вычисления распределения длины серии.} Аналитические выражения для распределения длины серии карты CUSUM-EX получены в работе Mukherjee, Graham and Chakraborti (2013). Однако удобнее использовать симуляции, как описано ниже.

\paragraph{Симуляционный подход.} Для получения характеристик длины серии карты CUSUM-EX можно использовать следующий алгоритм симуляции. Соответствующая программа SAS приведена в приложении B, программа SAS номер 9, а вывод показан ниже на рисунке \ref{fig-4.25}.

\noindent\textbf{Шаг 1.} Задаём объём эталонной выборки Фазы I \(m \ge 1\), объём тестовой выборки Фазы II \(n \ge 1\), нужную номинальную \(ARL_{IC}\) и величину сдвига (\(\delta\)), который нужно обнаружить. Затем задаём два распределения для генерации выборок Фазы I и Фазы II соответственно. В случае IC оба распределения совпадают (говорят, что сдвиг равен нулю (\(\delta = 0\))). В случае OOC распределение выборки Фазы II имеет тот же вид, что и у выборки Фазы I, но со сдвигом положения в единицах стандартного отклонения совокупности. Выборка Фазы I, например, может быть извлечена из нормального распределения со средним 0 и стандартным отклонением 1, тогда как в случае OOC выборка Фазы II может быть извлечена из нормального распределения с некоторым средним \(\delta \ne 0\) и стандартным отклонением 1 (говорят, что сдвиг среднего равен \(\delta \ne 0\)). Можно использовать любое непрерывное распределение, симметричное или асимметричное.

\noindent\textbf{Шаг 2.} Задаём два параметра проектирования \(k\) и \(h\) (читателю рекомендуется обратиться к разделу 3.2.2 за руководством по выбору параметров проектирования).

\noindent\textbf{Шаг 3.} Для каждой подгруппы вычисляем статистику \(U_{j,r}\) и соответствующие статистики для нанесения на карту \(C_j^-\) и \(C_j^+\).

\noindent\textbf{Шаг 4.} Сравниваем статистику для нанесения на карту с интервалом решения \(h\).

\noindent\textbf{Шаг 5.} Число подгрупп, необходимых до первого попадания \(C_j^-\) на \(-h\) или ниже либо \(C_j^+\) на \(h\) или выше, записывается как наблюдение из распределения длины серии.

\noindent\textbf{Шаг 6.} Повторяем шаги 1--5 в общей сложности 100\,000 раз.

\noindent\textbf{Шаг 7.} Наблюдения, полученные на шаге 5, можно проанализировать для вычисления оценок характеристик распределения длины серии, таких как среднее, медиана, стандартное отклонение и т.~д.

Для программирования шагов симуляции можно использовать программный пакет, такой как SAS или R. Отметим, что при \(\delta = 0\) мы получаем характеристики длины серии в состоянии IC, а при \(\delta \ne 0\) --- характеристики длины серии в состоянии OOC.

\begin{figure}[!ht]
	\centering
	\includegraphics[width=0.75\linewidth]{"./images/fig-4.25.png"}
	\caption[]{Вывод SAS, где параметры проектирования (\(m = 125\), \(n = 5\), \(j = 3\), \(h = 9.775\), \(k = n(d^* - d)\)) карты CUSUM-EX были выбраны для номинальной \(ARL_{IC} = 500\).}
	\label{fig-4.25}
\end{figure}

\subsubsection{Контрольные карты CUSUM на основе статистики суммы рангов Уилкоксона}
Как и в случае статистики предшествования/превышения, естественно рассмотреть карты CUSUM и EWMA, основанные на статистике MW или, эквивалентно, статистике WRS. Именно это сделали Li, Tang и Ng (2010), следуя работе Chakraborti и Van de Wiel (2008), и ввиду отмеченной ранее эквивалентности статистик WRS и MW мы в дальнейшем называем эти карты картами Уилкоксона--Манна--Уитни (WMW). Ожидается, что эти карты будут более чувствительными, чем карты Шухарта, к малым устойчивым сдвигам положения. Сначала рассмотрим карты CUSUM.

Свободные от распределения карты CUSUM, основанные на статистике WMW, в дальнейшем будем называть картами CUSUM-rank. Верхняя односторонняя карта CUSUM-rank определяется как
\begin{equation}
	C_j^+ = \max\left[ 0, C_{j-1}^+ + W_j - \frac{n(m + n + 1)}{2} - k \right], \quad j = 1, 2, 3, \ldots
	\label{eq:4.56}
\end{equation}
с начальным значением \(C_0^+ = 0\). Эта карта подаёт сигнал при первом \(j\), для которого \(C_j^+ \ge h\).

Напротив, для нижней односторонней карты CUSUM-rank определим
\begin{equation}
	C_j^- = \min\left[ 0, C_{j-1}^- + W_j - \frac{n(m + n + 1)}{2} + k \right], \quad j = 1, 2, 3, \ldots
	\label{eq:4.57}
\end{equation}
с начальным значением \(C_0^- = 0\). Она подаёт сигнал при первом \(j\), для которого \(C_j^- \le -h\).

Соответствующая двусторонняя карта CUSUM-rank подаёт сигнал при первом \(j\), для которого подаёт сигнал любая из односторонних карт, то есть всякий раз, когда \(C_j^+ \ge h\) или \(C_j^- \le -h\). Отметим, что здесь интервал решения \(h > 0\) и референсное значение \(k > 0\) являются (как и раньше) параметрами процедуры контрольной карты.

\paragraph{Методы вычисления распределения длины серии.} Выражения для распределения длины серии карты CUSUM-rank получены в работе Li, Tang и Ng (2010) с использованием метода CUC. Некоторые выводы приведены в приложении 4.2. Однако удобнее использовать симуляции, как описано ниже.

\paragraph{Симуляционный подход.} Для получения характеристик длины серии карты CUSUM-rank можно использовать следующий алгоритм симуляции. Соответствующая программа SAS приведена в приложении B, программа SAS номер 10, а вывод показан ниже на рисунке \ref{fig-4.26}.

\noindent\textbf{Шаг 1.} Задаём объём эталонной выборки Фазы I \(m \ge 1\), объём тестовой выборки Фазы II \(n \ge 1\) и величину сдвига (\(\delta\)), который нужно обнаружить. Затем задаём два распределения для генерации выборок Фазы I и Фазы II соответственно. В случае IC оба распределения совпадают (говорят, что сдвиг равен нулю (\(\delta = 0\))). В случае OOC распределение выборки Фазы II имеет тот же вид, что и у выборки Фазы I, но со сдвигом положения в единицах стандартного отклонения совокупности. Выборка Фазы I, например, может быть извлечена из нормального распределения со средним 0 и стандартным отклонением 1, тогда как в случае OOC выборка Фазы II может быть извлечена из нормального распределения с некоторым средним \(\delta \ne 0\) и стандартным отклонением 1 (говорят, что сдвиг среднего равен \(\delta \ne 0\)). Можно использовать любое непрерывное распределение, симметричное или асимметричное.

\noindent\textbf{Шаг 2.} Задаём два параметра проектирования \(k\) и \(h\) (читателю рекомендуется обратиться к разделу 3.2.2 за руководством по выбору параметров проектирования).

\noindent\textbf{Шаг 3.} Для каждой подгруппы вычисляем статистики \(W_j\) и соответствующие статистики для нанесения на карту \(C_j^-\) и \(C_j^+\).

\noindent\textbf{Шаг 4.} Сравниваем статистики для нанесения на карту с интервалом решения \(h\).

\noindent\textbf{Шаг 5.} Число подгрупп, необходимых до первого попадания \(C_j^-\) на \(-h\) или ниже либо \(C_j^+\) на \(h\) или выше, записывается как наблюдение из распределения длины серии.

\noindent\textbf{Шаг 6.} Повторяем шаги 1--5 в общей сложности 100\,000 раз. Это создаёт 100\,000 наблюдений из распределения длины серии.

\noindent\textbf{Шаг 7.} Наблюдения, полученные на шаге 5, можно проанализировать для вычисления оценок характеристик распределения длины серии, таких как среднее, медиана, стандартное отклонение и т.~д.

Опять же, для программирования шагов симуляции можно использовать программный пакет, такой как SAS или R. Отметим, что при \(\delta = 0\) мы получаем характеристики длины серии в состоянии IC, а при \(\delta \ne 0\) --- характеристики длины серии в состоянии OOC.

\begin{figure}[!ht]
	\centering
	\includegraphics[width=0.75\linewidth]{"./images/fig-4.26.png"}
	\caption[]{Вывод SAS, где параметры проектирования (\(m = 100\), \(n = 5\), \(h = 353\), \(k = 0\)) карты CUSUM-rank были выбраны для номинальной \(ARL_{IC} = 500\).}
	\label{fig-4.26}
\end{figure}

Далее рассмотрим некоторые свободные от распределения карты EWMA в случае U. Это обсуждение основано на работе Graham, Mukherjee and Chakraborti (2012).

\subsection{Контрольные карты EWMA}
\subsubsection{Контрольные карты EWMA на основе статистики превышения}
Процедура контрольной карты EWMA, основанная на статистике превышения \(U_{j,r}\), накапливает статистики \(U_{1,r}, U_{2,r}, U_{3,r}, \ldots\) с EWMA-статистиками для нанесения на карту, определяемыми как
\begin{equation}
	Z_j = \lambda U_{j,r} + (1 - \lambda)Z_{j-1}, \quad j = 1, 2, 3, \ldots
	\label{eq:4.58}
\end{equation}
с весовой константой \(0 < \lambda \le 1\).

Центральная линия CL и точные контрольные границы задаются как
\begin{equation}
\begin{aligned}
	UCL &= n(1 - a)\left( 1 - (1 - \lambda)^j \right) \\
	    &\quad + L\sqrt{\frac{na(1 - a)}{m + 2}n\left( 1 - (1 - \lambda)^j \right)^2 + \frac{\lambda(m + 1)}{2 - \lambda}\left( 1 - (1 - \lambda)^{2j} \right)} \\
	CL &= n(1 - a)\left( 1 - (1 - \lambda)^j \right) \\
	LCL &= n(1 - a)\left( 1 - (1 - \lambda)^j \right) \\
	    &\quad - L\sqrt{\frac{na(1 - a)}{m + 2}n\left( 1 - (1 - \lambda)^j \right)^2 + \frac{\lambda(m + 1)}{2 - \lambda}\left( 1 - (1 - \lambda)^{2j} \right)},
\end{aligned}
\label{eq:4.59}
\end{equation}
где \(a = r/(m + 1)\), а начальное значение \(Z_0 = E\left( U_{j-k,r} \mid X_{(r)} \right) = 0\) при использовании точных контрольных границ. Подробный вывод точных меняющихся во времени контрольных границ для карты EWMA-EX см. в приложении 4.3.3, результат 4.10. Асимптотические, или стационарные, контрольные границы, то есть при \(j \to \infty\), задаются как
\begin{equation}
\begin{aligned}
	UCL &= n(1 - a) + L\sqrt{\frac{na(1 - a)}{m + 2}\left( n + \frac{\lambda(m + 1)}{2 - \lambda} \right)} \\
	CL &= n(1 - a) \\
	LCL &= n(1 - a) - L\sqrt{\frac{na(1 - a)}{m + 2}\left( n + \frac{\lambda(m + 1)}{2 - \lambda} \right)},
\end{aligned}
\label{eq:4.60}
\end{equation}
где начальное значение \(Z_0 = E\left( U_{j-k,r} \mid X_{(r)} \right) = np_r\) при использовании стационарных контрольных границ. Подробный вывод стационарных контрольных границ для карты EWMA-EX см. в приложении 4.3.3, результат 4.11. Процесс считается находящимся в состоянии IC, пока все статистики для нанесения на карту \(Z_j\), \(j = 1, 2, 3, \ldots\), находятся между двумя контрольными границами, но как только какая-либо статистика для нанесения на карту попадает на верхнюю или нижнюю контрольную границу или за них, процесс объявляется вышедшим из-под контроля (OOC), и обычно начинается поиск назначаемых причин. Вывод точных контрольных границ приведён в приложении 4.3.3, результат 4.10.

\paragraph{Методы вычисления распределения длины серии.} Выражения для распределения длины серии карты EWMA-EX получены в работе Graham, Mukherjee and Chakraborti (2012), к которой может обратиться заинтересованный читатель. Однако, опять же, удобнее использовать симуляции, как описано ниже.

\paragraph{Симуляционный подход.} Для получения характеристик длины серии карты EWMA-EX можно использовать следующий алгоритм симуляции. Соответствующая программа SAS приведена в приложении B, программа SAS номер 11, а вывод показан ниже на рисунке \ref{fig-4.27}.

\noindent\textbf{Шаг 1.} Задаём объём эталонной выборки Фазы I \(m \ge 1\), объём тестовой выборки Фазы II \(n \ge 1\), нужную номинальную \(ARL_{IC}\) и величину сдвига (\(\delta\)), который нужно обнаружить. Затем задаём два распределения для генерации выборок Фазы I и Фазы II соответственно. В случае IC оба распределения совпадают (говорят, что сдвиг равен нулю (\(\delta = 0\))). В случае OOC распределение выборки Фазы II имеет тот же вид, что и у выборки Фазы I, но со сдвигом положения в единицах стандартного отклонения совокупности. Выборка Фазы I, например, может быть извлечена из нормального распределения со средним 0 и стандартным отклонением 1, тогда как в случае OOC выборка Фазы II может быть извлечена из нормального распределения с некоторым средним \(\delta \ne 0\) и стандартным отклонением 1 (говорят, что сдвиг среднего равен \(\delta \ne 0\)). Можно использовать любое непрерывное распределение, симметричное или асимметричное.

\noindent\textbf{Шаг 2.} Задаём два параметра проектирования \(\lambda\) и \(L\) (читателю рекомендуется обратиться к разделу 3.2.3 за руководством по выбору параметров проектирования).

\noindent\textbf{Шаг 3.} Для каждой подгруппы вычисляем статистику \(U_{j,r}\) и соответствующую статистику для нанесения на карту \(Z_j\).

\noindent\textbf{Шаг 4.} Сравниваем статистики для нанесения на карту с контрольными границами.

\noindent\textbf{Шаг 5.} Число подгрупп, необходимых до первого попадания \(Z_j\) на контрольные границы или за них, записывается как наблюдение из распределения длины серии.

\noindent\textbf{Шаг 6.} Повторяем шаги 1--5 в общей сложности 100\,000 раз. Это создаёт 100\,000 наблюдений из распределения длины серии.

\noindent\textbf{Шаг 7.} Наблюдения, полученные на шаге 5, можно проанализировать для вычисления оценок характеристик распределения длины серии, таких как среднее, медиана, стандартное отклонение и т.~д.

Как отмечалось ранее, для программирования исследования симуляций можно использовать программный пакет, такой как SAS или R. Отметим, что при \(\delta = 0\) мы получаем характеристики длины серии в состоянии IC, а при \(\delta \ne 0\) --- характеристики длины серии в состоянии OOC.

\begin{figure}[!ht]
	\centering
	\includegraphics[width=0.75\linewidth]{"./images/fig-4.27.png"}
	\caption[]{Вывод SAS, где параметры проектирования (\(m = 100\), \(n = 5\), \(\lambda = 0.10\), \(L = 2.220\)) карты EWMA-EX были выбраны для номинальной \(ARL_{IC} = 500\).}
	\label{fig-4.27}
\end{figure}

Подобно карте CUSUM-rank, можно рассмотреть карту EWMA-rank, основанную на статистике WRS. Она описана далее.

\subsubsection{Контрольные карты EWMA на основе статистики суммы рангов Уилкоксона}
Li et al. (2010) рассмотрели контрольную карту EWMA, основанную на статистике WRS \(W_j\); мы называем её картой EWMA-rank. Эта карта накапливает статистики \(W_1, W_2, W_3, \ldots\) в соответствии с формулой EWMA. Результирующая статистика для нанесения на карту EWMA-rank определяется как
\begin{equation}
	Z_j = \lambda W_j + (1 - \lambda)Z_{j-1}, \quad j = 1, 2, 3, \ldots
	\label{eq:4.61}
\end{equation}
где весовая константа \(0 < \lambda \le 1\).

Центральная линия и точные контрольные границы задаются как
\begin{equation}
\begin{aligned}
	UCL &= \mu_W + L\sqrt{\sigma_W^2\frac{\lambda}{2 - \lambda}\left( 1 - (1 - \lambda)^{2j} \right)} \\
	CL &= \mu_W \\
	LCL &= \mu_W - L\sqrt{\sigma_W^2\frac{\lambda}{2 - \lambda}\left( 1 - (1 - \lambda)^{2j} \right)},
\end{aligned}
\label{eq:4.62}
\end{equation}
где \(\mu_W = \frac{n(m + n + 1)}{2}\) и \(\sigma_W^2 = \frac{mn(m + n + 1)}{12}\) обозначают среднее и дисперсию \(W\) в состоянии IC соответственно, а начальное значение \(Z_0 = \mu_W\) при использовании точных контрольных границ. Асимптотические, или стационарные, контрольные границы, то есть при \(j \to \infty\), задаются как
\begin{equation}
\begin{aligned}
	UCL &= \mu_W + L\sqrt{\sigma_W^2\frac{\lambda}{2 - \lambda}} \\
	CL &= \mu_W \\
	LCL &= \mu_W - L\sqrt{\sigma_W^2\frac{\lambda}{2 - \lambda}},
\end{aligned}
\label{eq:4.63}
\end{equation}
где начальное значение \(Z_0 = \mu_W\) при использовании стационарных контрольных границ. Процесс считается находящимся в состоянии IC, пока все статистики для нанесения на карту \(Z_j\), \(j = 1, 2, 3, \ldots\), находятся между двумя контрольными границами, но как только какая-либо статистика для нанесения на карту попадает на UCL или LCL либо за них, процесс объявляется вышедшим из-под контроля (OOC), и обычно начинается поиск назначаемых причин.

\begin{figure}[!ht]
	\centering
	\includegraphics[width=0.75\linewidth]{"./images/fig-4.28.png"}
	\caption[]{Вывод SAS, где параметры проектирования (\(m = 100\), \(n = 5\), \(\lambda = 0.20\), \(LCL = 195.5\), \(UCL = 334.5\)) карты EWMA-rank были выбраны для номинальной \(ARL_{IC} = 370\).}
	\label{fig-4.28}
\end{figure}

\paragraph{Методы вычисления распределения длины серии.} Шаги аналогичны шагам получения условного и безусловного распределения длины серии для карты CUSUM-rank, описанным в работе Li, Tang и Ng (2010). Читателю рекомендуется обратиться к приложению 4.2 за подробностями. Однако использовать симуляции гораздо проще и удобнее. Алгоритм, аналогичный приведённым для карт Шухарта и CUSUM, показан ниже.

\paragraph{Симуляционный подход.} Для получения характеристик длины серии карты EWMA-rank можно использовать следующий алгоритм симуляции. Соответствующая программа SAS приведена в приложении B, программа SAS номер 12, а вывод показан ниже на рисунке \ref{fig-4.28}.

\noindent\textbf{Шаг 1.} Задаём объём эталонной выборки Фазы I \(m \ge 1\), объём тестовой выборки Фазы II \(n \ge 1\), нужную номинальную \(ARL_{IC}\) и величину сдвига (\(\delta\)), который нужно обнаружить. Затем задаём два распределения для генерации выборок Фазы I и Фазы II соответственно. В случае IC оба распределения совпадают (говорят, что сдвиг равен нулю (\(\delta = 0\))). В случае OOC распределение выборки Фазы II имеет тот же вид, что и у выборки Фазы I, но со сдвигом положения в единицах стандартного отклонения совокупности. Выборка Фазы I, например, может быть извлечена из нормального распределения со средним 0 и стандартным отклонением 1, тогда как в случае OOC выборка Фазы II может быть извлечена из нормального распределения с некоторым средним \(\delta \ne 0\) и стандартным отклонением 1 (говорят, что сдвиг среднего равен \(\delta \ne 0\)). Можно использовать любое непрерывное распределение, симметричное или асимметричное.

\noindent\textbf{Шаг 2.} Задаём параметр проектирования \(\lambda\) (читателю рекомендуется обратиться к разделу 3.2.3 за руководством по выбору этого параметра) и контрольные границы LCL и UCL так, чтобы получить номинальную \(ARL_{IC}\). Следует отметить в этом месте, что обычно для карты EWMA есть два параметра проектирования, \(\lambda\) и \(L\), которые нужно задать. Однако Li, Tang и Ng (2010), судя по всему, в своей работе напрямую выбрали контрольные границы (LCL и UCL), а не ширину контрольных границ (\(L\)).

\noindent\textbf{Шаг 3.} Для каждой подгруппы вычисляем статистику \(W_j\) и соответствующую статистику для нанесения на карту \(Z_j\).

\noindent\textbf{Шаг 4.} Сравниваем статистики для нанесения на карту с контрольными границами.

\noindent\textbf{Шаг 5.} Число подгрупп, необходимых до первого попадания \(Z_j\) на контрольные границы или за них, записывается как наблюдение из распределения длины серии.

\noindent\textbf{Шаг 6.} Повторяем шаги 1--5 в общей сложности 100\,000 раз. Это создаёт 100\,000 наблюдений из распределения длины серии.

\noindent\textbf{Шаг 7.} Наблюдения, полученные на шаге 5, можно проанализировать для вычисления оценок характеристик распределения длины серии, таких как среднее, медиана, стандартное отклонение и т.~д. 

Для программирования исследования симуляций можно использовать программный пакет, такой как SAS или R. Отметим, что при \(\delta = 0\) мы получаем характеристики длины серии в состоянии IC, а при \(\delta \ne 0\) --- характеристики длины серии в состоянии OOC.

Далее мы приведём несколько иллюстративных примеров свободных от распределения карт в случае U.

\section{Свободные от распределения контрольные карты в случае U: иллюстративные примеры}
\subsection{Контрольные карты Шухарта}

\noindent\textbf{Пример 4.18. Карта Шухарта Med.}

В таблице \ref{tab:4.21} приведены \(m = 100\) эталонных наблюдений (расположенных в порядке возрастания) из неизвестного непрерывного распределения процесса. Предположим, что мы хотим контролировать положение распределения процесса с помощью выборочных медиан будущих тестовых выборок объёма \(n = 5\). Кроме того, предположим, что номинальное значение \(ARL_{IC}\) равно 370. Тогда по таблице \ref{tab:4.18} находим, что \(a = 6\) и \(b = 100 - 6 + 1 = 95\), так что контрольные границы карты Med задаются 6-й и 95-й порядковыми статистиками эталонной выборки, \(\widehat{LCL} = X_{(6:100)}\) и \(\widehat{UCL} = X_{(95:100)}\), которые равны 4.39 и 9.83 соответственно и дают достигнутую \(ARL\), равную 359.6. Отметим, что достигнутая \(ARL_{IC}\) не равна номинальному значению 370, но достаточно близка к нему. Это поднимает важный и интересный вопрос для многих используемых на практике контрольных карт, статистики для нанесения на карту которых имеют дискретное распределение вероятностей. Именно так обстоит дело, в частности, с картами предшествования и в целом с непараметрическими контрольными картами. Фактически так же обстоит дело с контрольными картами по альтернативному признаку. В этих случаях часто нельзя достичь заданной номинальной \(ARL_{IC}\) точно и приходится довольствоваться достаточно близким значением либо использовать идею рандомизированного критерия (см. Gibbons and Chakraborti, 2010, с.~25), которая не рекомендуется на практике. Мы применяем консервативный подход, при котором константа \(a\) находится так, чтобы достигнутая \(ARL_{IC}\) была не больше номинальной \(ARL_{IC}\).

Далее были получены 25 тестовых выборок объёма \(n = 5\), и мониторинг процесса продолжался на основе выборочной медианы \(Y_{(3:5)}\). Эти значения приведены в столбцах (a) и (b) соответственно таблицы \ref{tab:4.22}, а двусторонняя контрольная карта Med показана на рисунке \ref{fig-4.29}.

Из рисунка \ref{fig-4.29} видно, что отсутствие статистической управляемости не наблюдается, и, соответственно, корректирующие меры не требуются; поэтому мониторинг процесса продолжается.

\begin{table}[!ht]
	\centering
	\caption{Данные Фазы I и порядковые статистики для контрольной карты Med.}
	\label{tab:4.21}
	\small
	\begin{tabular}{@{}rrrrrrrrrr@{}}
		\toprule
		\(i\) & \(X_{(i:100)}\) & \(i\) & \(X_{(i:100)}\) & \(i\) & \(X_{(i:100)}\) & \(i\) & \(X_{(i:100)}\) & \(i\) & \(X_{(i:100)}\) \\
		\midrule
		1 & 3.13 & 21 & 5.59 & 41 & 6.45 & 61 & 7.32 & 81 & 8.30 \\
		2 & 3.72 & 22 & 5.62 & 42 & 6.46 & 62 & 7.37 & 82 & 8.31 \\
		3 & 3.82 & 23 & 5.63 & 43 & 6.51 & 63 & 7.48 & 83 & 8.35 \\
		4 & 4.23 & 24 & 5.73 & 44 & 6.52 & 64 & 7.67 & 84 & 8.64 \\
		5 & 4.34 & 25 & 5.74 & 45 & 6.53 & 65 & 7.68 & 85 & 8.67 \\
		6 & 4.39 & 26 & 5.77 & 46 & 6.55 & 66 & 7.70 & 86 & 8.80 \\
		7 & 4.46 & 27 & 5.84 & 47 & 6.64 & 67 & 7.73 & 87 & 8.86 \\
		8 & 4.58 & 28 & 5.90 & 48 & 6.65 & 68 & 7.81 & 88 & 8.91 \\
		9 & 4.66 & 29 & 5.96 & 49 & 6.72 & 69 & 7.84 & 89 & 8.96 \\
		10 & 4.72 & 30 & 6.02 & 50 & 6.87 & 70 & 7.89 & 90 & 9.01 \\
		11 & 4.83 & 31 & 6.02 & 51 & 6.87 & 71 & 7.93 & 91 & 9.17 \\
		12 & 4.94 & 32 & 6.05 & 52 & 6.95 & 72 & 7.96 & 92 & 9.17 \\
		13 & 5.08 & 33 & 6.15 & 53 & 6.96 & 73 & 8.01 & 93 & 9.49 \\
		14 & 5.08 & 34 & 6.21 & 54 & 7.00 & 74 & 8.04 & 94 & 9.60 \\
		15 & 5.20 & 35 & 6.23 & 55 & 7.04 & 75 & 8.07 & 95 & 9.83 \\
		16 & 5.33 & 36 & 6.32 & 56 & 7.11 & 76 & 8.14 & 96 & 9.88 \\
		17 & 5.37 & 37 & 6.36 & 57 & 7.17 & 77 & 8.14 & 97 & 9.96 \\
		18 & 5.42 & 38 & 6.39 & 58 & 7.20 & 78 & 8.24 & 98 & 10.29 \\
		19 & 5.46 & 39 & 6.43 & 59 & 7.20 & 79 & 8.24 & 99 & 10.31 \\
		20 & 5.53 & 40 & 6.44 & 60 & 7.21 & 80 & 8.29 & 100 & 10.56 \\
		\bottomrule
	\end{tabular}
\end{table}

\begin{table}[!ht]
	\centering
	\caption{Данные Фазы II и соответствующие статистики для нанесения на карту для контрольной карты Med.}
	\label{tab:4.22}
	\small
	\begin{tabular}{@{}ccccccr@{}}
		\toprule
		& \multicolumn{5}{c}{(a)} & (b) \\
		\(i\) & \(Y_{i:1}\) & \(Y_{i:2}\) & \(Y_{i:3}\) & \(Y_{i:4}\) & \(Y_{i:5}\) & \(Y_{(3:5)}\) \\
		\midrule
		1 & 7.97 & 8.08 & 8.62 & 6.82 & 6.56 & 7.97 \\
		2 & 8.80 & 7.89 & 2.82 & 8.56 & 7.13 & 7.89 \\
		3 & 6.70 & 7.26 & 6.51 & 8.00 & 5.61 & 6.70 \\
		4 & 3.82 & 5.85 & 7.32 & 7.47 & 6.92 & 6.92 \\
		5 & 6.77 & 5.13 & 7.13 & 5.96 & 8.25 & 6.77 \\
		6 & 7.37 & 8.18 & 4.53 & 7.14 & 6.87 & 7.14 \\
		7 & 6.01 & 6.78 & 9.14 & 7.85 & 4.10 & 6.78 \\
		8 & 6.44 & 5.26 & 6.72 & 7.38 & 7.25 & 6.72 \\
		9 & 6.87 & 5.19 & 6.78 & 5.69 & 7.37 & 6.78 \\
		10 & 6.54 & 7.46 & 8.60 & 7.25 & 8.14 & 7.46 \\
		11 & 8.66 & 7.92 & 6.54 & 6.18 & 7.97 & 7.92 \\
		12 & 4.43 & 4.20 & 6.34 & 6.24 & 6.13 & 6.13 \\
		13 & 7.35 & 7.54 & 10.28 & 7.69 & 6.30 & 7.54 \\
		14 & 4.54 & 9.60 & 5.66 & 6.55 & 6.85 & 6.55 \\
		15 & 8.43 & 7.32 & 5.65 & 9.72 & 8.77 & 8.43 \\
		16 & 8.51 & 8.60 & 7.09 & 7.66 & 6.83 & 7.66 \\
		17 & 9.11 & 5.06 & 6.09 & 8.54 & 7.61 & 7.61 \\
		18 & 7.52 & 9.15 & 7.18 & 5.97 & 8.44 & 7.52 \\
		19 & 6.55 & 6.05 & 5.41 & 4.68 & 8.00 & 6.05 \\
		20 & 7.07 & 8.29 & 6.59 & 4.55 & 3.75 & 6.59 \\
		21 & 5.81 & 4.89 & 6.57 & 6.53 & 6.21 & 6.21 \\
		22 & 6.86 & 8.65 & 6.67 & 7.99 & 7.60 & 7.60 \\
		23 & 5.47 & 7.29 & 5.92 & 9.25 & 9.32 & 7.29 \\
		24 & 8.20 & 7.62 & 5.97 & 3.72 & 5.83 & 5.97 \\
		25 & 6.26 & 7.98 & 6.38 & 6.09 & 7.45 & 6.38 \\
		\bottomrule
	\end{tabular}
\end{table}

\begin{figure}[!ht]
	\centering
	\includegraphics[width=0.7\linewidth]{"./images/fig-4.29.png"}
	\caption[]{Контрольная карта Med.}
	\label{fig-4.29}
\end{figure}

\noindent\textbf{Пример 4.19. Карта Shewhart-MW.}

В таблице \ref{tab:4.23} показаны 25 ретроспективных, то есть выборок Фазы I, объёма пять каждая, собранных, когда процесс считался находящимся в состоянии IC, то есть \(m = 125\). Эти данные рассматриваются как эталонные данные Фазы I, для которых критерий согласия для нормальности не отвергается. Из работы Chakraborti и Van de Wiel (2008) используем \(b_{mn} = 540\) и \(a_{mn} = mn - b_{mn} = (125)(5) - 540 = 85\). Как найти контрольные границы, мы обсудим позже в этом разделе.

Столбец (a) таблицы \ref{tab:4.24} содержит 15 перспективных (Фазы II) выборок объёма \(n = 5\) значений внутренних диаметров поршневых колец, изготовленных в процессе ковки. Соответствующие статистики для нанесения на карту MW показаны в столбце (b).

Графическое представление контрольной карты Shewhart-MW показано на рисунке \ref{fig-4.30}; карта подаёт сигнал на выборке номер 12, что указывает на необходимость поиска назначаемых причин. Фактически выборки 12, 13 и 14, по-видимому, все находятся в состоянии OOC.

Далее рассмотрим свободные от распределения карты CUSUM.

\begin{table}[!ht]
	\centering
	\caption{Данные Фазы I для карты Shewhart-MW.}
	\label{tab:4.23}
	\small
	\begin{tabular}{@{}rccccc@{}}
		\toprule
		\(j\) & \(X_{j1}\) & \(X_{j2}\) & \(X_{j3}\) & \(X_{j4}\) & \(X_{j5}\) \\
		\midrule
		1 & 74.030 & 74.002 & 74.019 & 73.992 & 74.008 \\
		2 & 73.995 & 73.992 & 74.001 & 74.011 & 74.004 \\
		3 & 73.998 & 74.024 & 74.021 & 74.005 & 74.002 \\
		4 & 74.002 & 73.996 & 73.993 & 74.015 & 74.009 \\
		5 & 73.992 & 74.007 & 74.015 & 73.989 & 74.014 \\
		6 & 74.009 & 73.994 & 73.997 & 73.985 & 73.993 \\
		7 & 73.995 & 74.006 & 73.994 & 74.000 & 74.005 \\
		8 & 73.985 & 74.003 & 73.993 & 74.015 & 73.988 \\
		9 & 74.008 & 73.995 & 74.009 & 74.005 & 74.004 \\
		10 & 73.998 & 74.000 & 73.990 & 74.007 & 73.995 \\
		11 & 73.994 & 73.998 & 73.994 & 73.995 & 73.990 \\
		12 & 74.004 & 74.000 & 74.007 & 74.000 & 73.996 \\
		13 & 73.983 & 74.002 & 73.998 & 73.997 & 74.012 \\
		14 & 74.006 & 73.967 & 73.994 & 74.000 & 73.984 \\
		15 & 74.012 & 74.014 & 73.998 & 73.999 & 74.007 \\
		16 & 74.000 & 73.984 & 74.005 & 73.998 & 73.996 \\
		17 & 73.994 & 74.012 & 73.986 & 74.005 & 74.007 \\
		18 & 74.006 & 74.010 & 74.018 & 74.003 & 74.000 \\
		19 & 73.984 & 74.002 & 74.003 & 74.005 & 73.997 \\
		20 & 74.000 & 74.010 & 74.013 & 74.020 & 74.003 \\
		21 & 73.982 & 74.001 & 74.015 & 74.005 & 73.996 \\
		22 & 74.004 & 73.999 & 73.990 & 74.006 & 74.009 \\
		23 & 74.010 & 73.989 & 73.990 & 74.009 & 74.014 \\
		24 & 74.015 & 74.008 & 73.993 & 74.000 & 74.010 \\
		25 & 73.982 & 73.984 & 73.995 & 74.017 & 74.013 \\
		\bottomrule
	\end{tabular}
\end{table}

\begin{table}[!ht]
	\centering
	\caption{Данные Фазы II для карты Shewhart-MW и соответствующие статистики для нанесения на карту.}
	\label{tab:4.24}
	\small
	\begin{tabular}{@{}rccccc r@{}}
		\toprule
		& \multicolumn{5}{c}{(a)} & (b) \\
		\multicolumn{1}{c}{Выборка \(j\)} & \(X_{j1}\) & \(X_{j2}\) & \(X_{j3}\) & \(X_{j4}\) & \(X_{j5}\) & \(M_{YX}\) \\
		\midrule
		1 & 74.012 & 74.015 & 74.030 & 73.986 & 74.000 & 429 \\
		2 & 73.995 & 74.010 & 73.990 & 74.015 & 74.001 & 333 \\
		3 & 73.987 & 73.999 & 73.985 & 74.000 & 73.990 & 143 \\
		4 & 74.008 & 74.010 & 74.003 & 73.991 & 74.006 & 371 \\
		5 & 74.003 & 74.000 & 74.001 & 73.986 & 73.997 & 242 \\
		6 & 73.994 & 74.003 & 74.015 & 74.020 & 74.004 & 411 \\
		7 & 74.008 & 74.002 & 74.018 & 73.995 & 74.005 & 393 \\
		8 & 74.001 & 74.004 & 73.990 & 73.996 & 73.998 & 241 \\
		9 & 74.015 & 74.000 & 74.016 & 74.025 & 74.000 & 471 \\
		10 & 74.030 & 74.005 & 74.000 & 74.016 & 74.012 & 486 \\
		11 & 74.001 & 73.990 & 73.995 & 74.010 & 74.024 & 341 \\
		12 & 74.015 & 74.020 & 74.024 & 74.005 & 74.019 & 561 \\
		13 & 74.035 & 74.010 & 74.012 & 74.015 & 74.026 & 576 \\
		14 & 74.017 & 74.013 & 74.036 & 74.025 & 74.026 & 602 \\
		15 & 74.010 & 74.005 & 74.029 & 74.000 & 74.020 & 485 \\
		\bottomrule
	\end{tabular}
\end{table}

\begin{figure}[!ht]
	\centering
	\includegraphics[width=0.7\linewidth]{"./images/fig-4.30.png"}
	\caption[]{Карта Shewhart-MW.}
	\label{fig-4.30}
\end{figure}

\subsection{Контрольные карты CUSUM}

\noindent\textbf{Пример 4.20. Двусторонняя медианная карта CUSUM-EX.}

Столбец (a) таблицы \ref{tab:4.25} содержит 15 перспективных (Фазы II) выборок объёма \(n = 5\) значений внутренних диаметров поршневых колец, изготовленных в процессе ковки. Для данных таблицы \ref{tab:4.25} вычислены две односторонние статистики для нанесения на карту CUSUM \(C_j^+\) и \(C_j^-\) с \(C_0^+ = C_0^- = 0\), а также два счётчика \(N^+\) и \(N^-\); они показаны в столбцах (c) и (d) таблицы \ref{tab:4.25}. Используем \(k = \delta \cdot STDEV(U_{j,r})\), где \(\delta = \gamma\sigma/\sqrt{n}\). Предполагая, что нужно обнаружить малый сдвиг, берём \(\gamma\) малым, скажем \(\gamma = 0.25\). Напомним из замечания 4.6 в разделе 4.4.2, что самым безопасным выбором для \(k\) является меньшее значение \(k\) или даже \(k = 0\); здесь значение \(k\) равно 0. Следующий шаг --- выбрать \(h\) в сочетании с выбранным \(k\) так, чтобы была достигнута нужная номинальная \(ARL_{IC}\). Здесь интервал решения был взят равным \(h = 7.5\). Опять же, читателю рекомендуется обратиться к замечанию 4.6 в разделе 4.4.2 и к работе Mukherjee, Graham and Chakraborti (2013) для подробного обсуждения выбора параметров проектирования.

Для иллюстрации расчётов рассмотрим выборку номер 1. Выражения для статистик для нанесения на карту CUSUM \(C_1^+\) и \(C_1^-\) имеют вид
\begin{align*}
	C_1^+ &= \max\left[ 0, C_0^+ + \left( U_{1,r} - \frac{n}{2} \right) - k \right] \\
	      &= \max\left[ 0, 0 + (3 - 5/2) - 0 \right] \\
	      &= 0.5,
\end{align*}
и
\begin{align*}
	C_1^- &= \min\left[ 0, C_0^- + \left( U_{1,r} - \frac{n}{2} \right) + k \right] \\
	      &= \min\left[ 0, 0 + (3 - 5/2) + 0 \right] \\
	      &= 0.
\end{align*}

Графическое представление медианной контрольной карты CUSUM-EX показано на рисунке \ref{fig-4.31}. Из рисунка \ref{fig-4.31} видно, что процесс выходит из-под контроля начиная с выборки номер 13, поскольку соответствующая статистика для нанесения на карту попадает выше UCL. Таким образом, необходимы корректирующая мера и поиск назначаемых причин. Кроме того, поскольку счётчик \(N^+ = 8\) на выборке номер 13 и \(N^+ = 0\) на выборке номер 5, сдвиг, скорее всего, произошёл на выборке номер 5.

\begin{table}[!ht]
	\centering
	\caption{Данные для медианной карты CUSUM-EX и соответствующие статистики для нанесения на карту.}
	\label{tab:4.25}
	\small
	\begin{tabular}{@{}rcccccccccc@{}}
		\toprule
		& \multicolumn{5}{c}{(a)} & (b) & \multicolumn{2}{c}{(c)} & \multicolumn{2}{c}{(d)} \\ \cmidrule(lr){2-6}\cmidrule(lr){8-9}\cmidrule(lr){10-11}
		\multicolumn{1}{c}{Выборка \(j\)} & \(X_{j1}\) & \(X_{j2}\) & \(X_{j3}\) & \(X_{j4}\) & \(X_{j5}\) & \(U_{j,r}\) & \(C_j^-\) & \(N^-\) & \(C_j^+\) & \(N^+\) \\
		\midrule
		1 & 74.012 & 74.015 & 74.030 & 73.986 & 74.000 & 3 & 0.0 & 0 & 0.5 & 1 \\
		2 & 73.995 & 74.010 & 73.990 & 74.015 & 74.001 & 2 & \(-0.5\) & 1 & 0.0 & 0 \\
		3 & 73.987 & 73.999 & 73.985 & 74.000 & 73.990 & 0 & \(-3.0\) & 2 & 0.0 & 0 \\
		4 & 74.008 & 74.010 & 74.003 & 73.991 & 74.006 & 4 & \(-1.5\) & 3 & 1.5 & 1 \\
		5 & 74.003 & 74.000 & 74.001 & 73.986 & 73.997 & 1 & \(-3.0\) & 4 & 0.0 & 0 \\
		6 & 73.994 & 74.003 & 74.015 & 74.020 & 74.004 & 4 & \(-1.5\) & 5 & 1.5 & 1 \\
		7 & 74.008 & 74.002 & 74.018 & 73.995 & 74.005 & 4 & 0.0 & 0 & 3.0 & 2 \\
		8 & 74.001 & 74.004 & 73.990 & 73.996 & 73.998 & 1 & \(-1.5\) & 1 & 1.5 & 3 \\
		9 & 74.015 & 74.000 & 74.016 & 74.025 & 74.000 & 3 & \(-1.0\) & 2 & 2.0 & 4 \\
		10 & 74.030 & 74.005 & 74.000 & 74.016 & 74.012 & 4 & 0.0 & 0 & 3.5 & 5 \\
		11 & 74.001 & 73.990 & 73.995 & 74.010 & 74.024 & 2 & \(-0.5\) & 1 & 3.0 & 6 \\
		12 & 74.015 & 74.020 & 74.024 & 74.005 & 74.019 & 5 & 0.0 & 0 & 5.5 & 7 \\
		13 & 74.035 & 74.010 & 74.012 & 74.015 & 74.026 & 5 & 0.0 & 0 & 8.0 & 8 \\
		14 & 74.017 & 74.013 & 74.036 & 74.025 & 74.026 & 5 & 0.0 & 0 & 10.5 & 9 \\
		15 & 74.010 & 74.005 & 74.029 & 74.000 & 74.020 & 4 & 0.0 & 0 & 12.0 & 10 \\
		\bottomrule
	\end{tabular}
\end{table}

\begin{figure}[!ht]
	\centering
	\includegraphics[width=0.7\linewidth]{"./images/fig-4.31.png"}
	\caption[]{Двусторонняя медианная карта CUSUM-EX.}
	\label{fig-4.31}
\end{figure}

\noindent\textbf{Пример 4.21. Двусторонняя карта CUSUM-rank.}

Столбец (a) таблицы \ref{tab:4.26} содержит 15 перспективных (Фазы II) выборок объёма \(n = 5\) значений внутренних диаметров поршневых колец, изготовленных в процессе ковки. Для данных таблицы \ref{tab:4.26} вычислены две односторонние статистики для нанесения на карту CUSUM \(C_j^+\) и \(C_j^-\) с \(C_0^+ = C_0^- = 0\), а также два счётчика \(N^+\) и \(N^-\); они показаны в таблице \ref{tab:4.26}. Используем референсное значение \(k = \tau_l SD(W_j)\), где \(\tau_l = c_l \delta\), \(c_l\) --- положительная константа, а \(\delta\) --- величина обнаруживаемого сдвига в единицах стандартного отклонения процесса. Для иллюстрации желаемый обнаруживаемый сдвиг был взят малым, \(\delta = 0.25\), и \(c_l = 1.00\), так что
\[
	k = \tau_l SD(W_j) = (1)(0.25)\sqrt{\frac{(125)(5)(125 + 5 + 1)}{12}} = 20.650.
\]
Интервал решения был взят равным \(h = 946\); пояснение следует ниже.

Кратко напомним о выборе параметров проектирования \((k, h)\). Величины \((k, h)\) выбираются так, чтобы карта имела заданную номинальную \(ARL_{IC} = t\) (или заданную номинальную \(MRL_{IC} = s\)) и была способна обнаружить сдвиг, особенно малый, как можно скорее. Отметим, что Mukherjee, Graham и Chakraborti (2013), среди прочих, рекомендуют использовать MRL, а не ARL в качестве характеристики эффективности карты. Первый шаг --- выбрать \(k\). В разделе 3.2.2 мы рассматривали параметрическую карту CUSUM для отдельных данных (\(n = 1\)) без эталонной выборки. Мы изучили влияние \(k\) (читателю рекомендуется обратиться к рисунку 3.2 и его обсуждению) и пришли к выводу, что когда имеется мало априорной информации или она отсутствует относительно величины сдвига, меньший выбор \(k\) является наиболее безопасным (чтобы защититься от ненужных задержек обнаружения). Mukherjee, Graham и Chakraborti (2013) приходят к аналогичным выводам для непараметрических карт CUSUM. Следующий шаг --- выбрать \(h\) в сочетании с выбранным \(k\) так, чтобы была достигнута нужная номинальная \(ARL_{IC}\) или \(MRL_{IC}\). В работе Mukherjee, Graham и Chakraborti (2013) \(h\) и \(k\) выбирались для номинальной \(MRL_{IC} \approx 350\). Для иллюстрации расчётов для карты CUSUM-rank рассмотрим выборку номер 1. Выражения для статистик для нанесения на карту CUSUM \(C_1^+\) и \(C_1^-\) имеют вид
\begin{align*}
	C_1^+ &= \max\left[ 0, C_{j-1}^+ + W_j - \frac{n(m + n + 1)}{2} - k \right] \\
	      &= \max\left[ 0, 0 + 429 - \frac{5(125 + 5 + 1)}{2} - 20.650 \right] \\
	      &= 80.850
\end{align*}
и
\begin{align*}
	C_1^- &= \min\left[ 0, C_{j-1}^- + W_j - \frac{n(m + n + 1)}{2} + k \right] \\
	      &= \min\left[ 0, 0 + 429 - \frac{5(125 + 5 + 1)}{2} + 20.650 \right] \\
	      &= 0.
\end{align*}

Эти значения и соответствующие счётчики показаны в столбцах (c) и (d) таблицы \ref{tab:4.26}. Графическое представление карты CUSUM-rank показано на рисунке \ref{fig-4.32}; из него видно, что процесс выходит из-под контроля на предпоследней выборке, то есть карта подаёт сигнал на выборке номер 14 при счётчике \(N^+ = 6\), что означает, что сдвиг, скорее всего, произошёл на выборке номер 8.

\begin{table}[!ht]
	\centering
	\caption{Данные для карты CUSUM-rank и соответствующие статистики для нанесения на карту.}
	\label{tab:4.26}
	\small
	\begin{tabular}{@{}rcccccccccc@{}}
		\toprule
		& \multicolumn{5}{c}{(a)} & (b) & \multicolumn{2}{c}{(c)} & \multicolumn{2}{c}{(d)} \\ \cmidrule(lr){2-6}\cmidrule(lr){8-9}\cmidrule(lr){10-11}
		\multicolumn{1}{c}{Выборка \(j\)} & \(X_{j1}\) & \(X_{j2}\) & \(X_{j3}\) & \(X_{j4}\) & \(X_{j5}\) & \(W_j\) & \(C_j^-\) & \(N^-\) & \(C_j^+\) & \(N^+\) \\
		\midrule
		1 & 74.012 & 74.015 & 74.030 & 73.986 & 74.000 & 429 & 0 & 0 & 80.850 & 1 \\
		2 & 73.995 & 74.010 & 73.990 & 74.015 & 74.001 & 333 & 0 & 0 & 65.700 & 2 \\
		3 & 73.987 & 73.999 & 73.985 & 74.000 & 73.990 & 143 & \(-164.350\) & 1 & 0 & 0 \\
		4 & 74.008 & 74.010 & 74.003 & 73.991 & 74.006 & 371 & \(-100.700\) & 2 & 22.350 & 1 \\
		5 & 74.003 & 74.000 & 74.001 & 73.986 & 73.997 & 242 & \(-166.049\) & 3 & 0 & 0 \\
		6 & 73.994 & 74.003 & 74.015 & 74.020 & 74.004 & 411 & \(-62.399\) & 4 & 62.350 & 1 \\
		7 & 74.008 & 74.002 & 74.018 & 73.995 & 74.005 & 393 & 0 & 0 & 107.200 & 2 \\
		8 & 74.001 & 74.004 & 73.990 & 73.996 & 73.998 & 241 & \(-66.350\) & 1 & 0 & 0 \\
		9 & 74.015 & 74.000 & 74.016 & 74.025 & 74.000 & 471 & 0 & 0 & 122.850 & 1 \\
		10 & 74.030 & 74.005 & 74.000 & 74.016 & 74.012 & 486 & 0 & 0 & 260.700 & 2 \\
		11 & 74.001 & 73.990 & 73.995 & 74.010 & 74.024 & 341 & 0 & 0 & 253.049 & 3 \\
		12 & 74.015 & 74.020 & 74.024 & 74.005 & 74.019 & 561 & 0 & 0 & 465.899 & 4 \\
		13 & 74.035 & 74.010 & 74.012 & 74.015 & 74.026 & 576 & 0 & 0 & 693.249 & 5 \\
		14 & 74.017 & 74.013 & 74.036 & 74.025 & 74.026 & 602 & 0 & 0 & 946.599 & 6 \\
		15 & 74.010 & 74.005 & 74.029 & 74.000 & 74.020 & 485 & 0 & 0 & 1082.948 & 7 \\
		\bottomrule
	\end{tabular}
\end{table}

\begin{figure}[!ht]
	\centering
	\includegraphics[width=0.7\linewidth]{"./images/fig-4.32.png"}
	\caption[]{Двусторонняя карта CUSUM-rank.}
	\label{fig-4.32}
\end{figure}

\subsection{Контрольные карты EWMA}

\noindent\textbf{Пример 4.22. Двусторонняя медианная карта EWMA-EX.}

Двусторонняя медианная карта CUSUM-EX была рассмотрена в примере 4.20 на данных таблицы \ref{tab:4.25}. В качестве альтернативы теперь к тем же данным применяется контрольная карта EWMA. Для иллюстрации используем \(\lambda = 0.1\) для вычисления статистик для нанесения на карту и \(L = 2.31\) для нахождения контрольных границ. Для иллюстрации расчётов рассмотрим выборку номер 1 с \(U_{1,r} = 3\) и предположим, что начальное значение \(Z_0 = np_r = (5)(0.5) = 2.5\). Тогда первая статистика для нанесения на карту на карте EWMA-EX равна
\begin{align*}
	Z_1 &= \lambda U_{1,r} + (1 - \lambda)Z_0 \\
	    &= 0.1(3) + (1 - 0.1)2.5 \\
	    &= 2.550.
\end{align*}
Значение второй статистики для нанесения на карту равно
\begin{align*}
	Z_2 &= \lambda U_{2,r} + (1 - \lambda)Z_1 \\
	    &= 0.1(3) + (1 - 0.1)2.550 \\
	    &= 2.595,
\end{align*}
а остальные расчёты сведены в столбец (b) таблицы \ref{tab:4.27}.

Для простоты используем стационарные контрольные границы на карте EWMA-EX, которые находятся по уравнению \eqref{eq:4.60}, то есть
\begin{align*}
	UCL &= n(1 - a) + L\sqrt{\frac{na(1 - a)}{m + 2}\left( n + \frac{\lambda(m + 1)}{2 - \lambda} \right)} \\
	    &= (5)(1 - 0.5) + 2.31\sqrt{\frac{(5)(0.5)(1 - 0.5)}{125 + 2}\left( 5 + \frac{(0.1)(125 + 1)}{2 - 0.1} \right)} \\
	    &= 3.282
\end{align*}
и
\begin{align*}
	LCL &= n(1 - a) - L\sqrt{\frac{na(1 - a)}{m + 2}\left( n + \frac{\lambda(m + 1)}{2 - \lambda} \right)} \\
	    &= (5)(1 - 0.5) - 2.31\sqrt{\frac{(5)(0.5)(1 - 0.5)}{125 + 2}\left( 5 + \frac{(0.1)(125 + 1)}{2 - 0.1} \right)} \\
	    &= 1.718
\end{align*}
при \(a = \frac{r}{m + 1} = \frac{63}{125 + 1} = 0.5\), а CL задаётся как \(CL = n(1 - a) = (5)(1 - 0.5) = 2.5\).

\begin{table}[!ht]
	\centering
	\caption{Расчёты для медианной карты EWMA-EX.}
	\label{tab:4.27}
	\small
	\begin{tabular}{@{}ccc@{}}
		\toprule
		Выборка & (a) & (b) \\
		\(j\) & \(U_{j,r}\) & \(Z_j = \lambda U_{j,r} + (1 - \lambda)Z_{j-1}\) \\
		\midrule
		1 & 3 & 2.550 \\
		2 & 3 & 2.595 \\
		3 & 0 & 2.336 \\
		4 & 4 & 2.502 \\
		5 & 2 & 2.452 \\
		6 & 4 & 2.607 \\
		7 & 4 & 2.746 \\
		8 & 2 & 2.671 \\
		9 & 3 & 2.704 \\
		10 & 4 & 2.834 \\
		11 & 3 & 2.850 \\
		12 & 5 & 3.065 \\
		13 & 5 & 3.259 \\
		14 & 5 & 3.433 \\
		15 & 4 & 3.490 \\
		\midrule
		16 & 3 & 2.550 \\
		17 & 3 & 2.595 \\
		18 & 0 & 2.336 \\
		19 & 4 & 2.502 \\
		20 & 2 & 2.452 \\
		\bottomrule
	\end{tabular}
	\par\smallskip
	\footnotesize Строки 16--20 воспроизведены по книге как есть; они повторяют строки 1--5.
\end{table}

\begin{figure}[!ht]
	\centering
	\includegraphics[width=0.7\linewidth]{"./images/fig-4.33.png"}
	\caption[]{Двусторонняя медианная карта EWMA-EX.}
	\label{fig-4.33}
\end{figure}

Результирующая карта EWMA-EX показана на рисунке \ref{fig-4.33}, и мы видим, что процесс выходит из-под контроля начиная с выборки номер 14, поскольку статистика для нанесения на карту попадает на UCL. Таким образом, необходимы корректирующая мера и поиск назначаемых причин.

\noindent\textbf{Пример 4.23. Двусторонняя карта EWMA-rank.}

Двусторонняя карта CUSUM-rank была построена в примере 4.21 на данных таблицы \ref{tab:4.26}. Здесь в качестве альтернативы мы рассматриваем контрольную карту EWMA. Для этой цели используем \(\lambda = 0.1\) для вычисления статистик для нанесения на карту и \(L = 2.814\) для нахождения контрольных границ. Для иллюстрации расчётов рассмотрим выборку номер 1 с \(W_1 = 429\) и предположим, что начальное значение \(Z_0 = \mu_W = \frac{n(m + n + 1)}{2} = \frac{5(125 + 5 + 1)}{2} = 327.5\).

Тогда первая статистика для нанесения на карту на карте EWMA-rank равна
\begin{align*}
	Z_1 &= \lambda W_1 + (1 - \lambda)Z_0 \\
	    &= 0.1(429) + (1 - 0.1)327.5 \\
	    &= 337.65.
\end{align*}
Значение второй статистики для нанесения на карту равно
\begin{align*}
	Z_2 &= \lambda W_2 + (1 - \lambda)Z_1 \\
	    &= 0.1(333) + (1 - 0.1)337.65 \\
	    &= 337.19,
\end{align*}
а остальные расчёты сведены в столбец (b) таблицы \ref{tab:4.28}.

\begin{table}[!ht]
	\centering
	\caption{Расчёты для карты EWMA-rank.}
	\label{tab:4.28}
	\small
	\begin{tabular}{@{}ccc@{}}
		\toprule
		Выборка & (a) & (b) \\
		\(j\) & \(W_j\) & \(Z_j = \lambda W_j + (1 - \lambda)Z_{j-1}\) \\
		\midrule
		1 & 429 & 337.65 \\
		2 & 333 & 337.19 \\
		3 & 143 & 317.72 \\
		4 & 371 & 322.99 \\
		5 & 242 & 314.85 \\
		6 & 411 & 324.41 \\
		7 & 393 & 331.27 \\
		8 & 241 & 322.19 \\
		9 & 471 & 337.07 \\
		10 & 486 & 351.97 \\
		11 & 341 & 350.82 \\
		12 & 561 & 371.84 \\
		13 & 576 & 392.20 \\
		14 & 602 & 413.13 \\
		15 & 485 & 420.27 \\
		\bottomrule
	\end{tabular}
\end{table}

Для простоты используем стационарные контрольные границы на карте EWMA-rank, которые находятся по уравнению \eqref{eq:4.63}, то есть
\begin{align*}
	UCL &= \mu_W + L\sqrt{\sigma_W^2\frac{\lambda}{2 - \lambda}} \\
	    &= \frac{n(m + n + 1)}{2} + L\sqrt{\frac{mn(m + n + 1)}{12}\frac{\lambda}{2 - \lambda}} \\
	    &= \frac{5(125 + 5 + 1)}{2} + 2.814\sqrt{\frac{(125)(5)(125 + 5 + 1)}{12}\frac{0.1}{2 - 0.1}} \\
	    &= 380.825
\end{align*}
и
\begin{align*}
	LCL &= \mu_W - L\sqrt{\sigma_W^2\frac{\lambda}{2 - \lambda}} \\
	    &= \frac{n(m + n + 1)}{2} - L\sqrt{\frac{mn(m + n + 1)}{12}\frac{\lambda}{2 - \lambda}} \\
	    &= \frac{5(125 + 5 + 1)}{2} - 2.814\sqrt{\frac{(125)(5)(125 + 5 + 1)}{12}\frac{0.1}{2 - 0.1}} \\
	    &= 274.175,
\end{align*}
а CL задаётся как \(CL = \mu_W = \frac{n(m + n + 1)}{2} = \frac{5(125 + 5 + 1)}{2} = 327.5\).

\begin{figure}[!ht]
	\centering
	\includegraphics[width=0.7\linewidth]{"./images/fig-4.34.png"}
	\caption[]{Двусторонняя карта EWMA-rank.}
	\label{fig-4.34}
\end{figure}

Результирующая карта EWMA-rank показана на рисунке \ref{fig-4.34}, и мы видим, что процесс выходит из-под контроля при увеличении параметра положения начиная с выборки номер 13, поскольку статистика для нанесения на карту попадает выше UCL. Таким образом, может потребоваться поиск назначаемых причин и некоторые последующие корректирующие действия.

\section{Влияние оценивания параметров}
Влияние оценивания параметров на эффективность контрольной карты Фазы II может быть значительным. Читателю рекомендуется обратиться к Jensen et al. (2006) за обзором литературы и некоторыми общими рекомендациями в этой области. Здесь мы просто отметим, что влияние объёма и качества эталонной выборки важно не только для непараметрических контрольных карт, но и для всех карт Фазы II, включая карты нормальной теории, когда параметры оцениваются по эталонным данным Фазы I. Влияние и эффекты оценивания параметров остаются активной областью текущих исследований в SPC, и читателю рекомендуется обратиться к этим работам. В заключение заметим, что, хотя хорошо известно, что использование оцениваемых параметров существенно влияет на статистическую эффективность карт Фазы II, исследование того, какими способами и средствами действительно получить эталонную выборку в состоянии IC, то есть выборку Фазы I, получило меньше внимания и должно быть рассмотрено в будущих исследованиях.

\section{Объём данных Фазы I}
Читателю рекомендуется обратиться к Jensen et al. (2006) за обзором литературы и некоторыми рекомендациями в этой области. По этой теме выполнено гораздо больше недавних исследований. Общее мнение состоит в том, что в Фазе I требуется гораздо больше эталонных данных, чем обычно рекомендуемые 20--30 выборок объёма 5, чтобы достичь сопоставимой эффективности в случае K, то есть в случае известных параметров. Важно отметить, что объём эталонных данных Фазы I зависит от многих факторов: от того, является ли карта картой по альтернативному признаку или по количественному признаку, от распределения процесса, от типа контрольной карты (Шухарта, CUSUM или EWMA), а также от размерности данных. В некоторых случаях, когда требуемый объём эталонных данных может оказаться недоступным, можно рассмотреть так называемые самозапускающиеся (self-starting) контрольные карты, которые используют последовательные наблюдения для одновременного обновления оценок параметров и статистики для нанесения на карту. Некоторые недавние публикации о самозапускающихся картах включают Maboudou-Tchao and Hawkins (2011), Zhang, Xie and Jin (2012) и Keefe, Woodall and Jones-Farmer (2015). Другой вариант --- каким-либо образом скорректировать константы для нанесения на карту. Эффективность как в состоянии IC, так и в состоянии OOC должна быть удовлетворительной. Мы не углубляемся здесь в эти темы; читатель может обратиться к некоторым недавним работам о влиянии оценивания параметров на контрольные карты.

\section{Улучшения контрольных карт}
\subsection{Чувствительность и правила сигналов на основе серий}
Правила сигналов на основе серий обычно используются для улучшения эффективности контрольных карт Шухарта. В отношении добавления правил сигналов к картам CUSUM и EWMA сделано очень мало. Для последних читателю рекомендуется обратиться к Riaz, Abbas and Does (2011), которые предложили две схемы правил серий для параметрических карт CUSUM, и к Abbas, Riaz and Does (2011), которые предлагают две схемы правил серий для параметрических карт EWMA. Однако эти работы посвящены параметрическим контрольным картам и, насколько нам известно, до настоящего времени никто не рассматривал добавление правил сигналов к картам CUSUM-EX, CUSUM-rank, EWMA-EX и EWMA-rank. Chakraborti, Eryilmaz and Human (2009) рассмотрели карту Shewhart-Prec с правилами сигналов на основе серий и пришли к выводу, что они превосходят типичную карту Shewhart-Prec «1 из 1». Malela-Majika, Chakraborti and Graham (2016) предлагают карту Shewhart-MW с правилами сигналов на основе серий и также обнаружили, что она превосходит карту Shewhart-MW «1 из 1», что и следует ожидать. Читателю рекомендуется обратиться к этим работам за дальнейшими подробностями. Здесь мы приведём лишь один пример, иллюстрирующий карты Шухарта с правилами сигналов в случае U. Для этого сначала получим безусловную ARL непараметрической карты Shewhart-Prec с правилами сигналов на основе серий, как предложено в работе Chakraborti, Eryilmaz and Human (2009).

Опять же предполагается, что имеется эталонная выборка объёма \(m\) из процесса в состоянии IC с неизвестной непрерывной функцией распределения \(F\). Пусть \(X_{1:m} < X_{2:m} < \cdots < X_{m:m}\) обозначают порядковые статистики эталонной выборки с \(LCL = X_{a:m}\) и \(UCL = X_{b:m}\), где \(1 \le a < b \le m\). В Фазе II тестовые выборки объёма \(n\) извлекаются последовательно и независимо друг от друга, а также независимо от эталонной выборки. Напомним, что для карты Shewhart-Prec статистикой для нанесения на карту для \(h\)-й тестовой выборки \(Y_{jh:n}\) может быть любая порядковая статистика, но для простоты выбирается медиана, а результирующая карта называется картой Med. Для \(h\)-й тестовой выборки определяются следующие индикаторные переменные: \(Z_h = 1\), если \(Y_{jh:n} \notin (LCL, UCL)\), и \(Z_h = 0\), если \(Y_{jh:n} \in (LCL, UCL)\), где \(h = 1, 2, 3, \ldots\). Поскольку контрольные границы являются порядковыми статистиками эталонной выборки, индикаторы сигнала \(Z_1, Z_2, \ldots\) являются зависимыми бинарными переменными. Также предполагается, что тестовая выборка происходит из непрерывного распределения с функцией распределения \(G\).

В качестве примера для правила сигнала «1 из 1» (которое соответствует исходной карте Med) условная вероятность отсутствия сигнала при заданных \(X_{a:m} = x_1\) и \(X_{b:m} = x_2\) равна \(P\left( Z_h = 0 \mid X_{a:m} = x_1, X_{b:m} = x_2 \right) = P\left( x_1 < Y_{j:n} < x_2 \mid X_{a:m} = x_1, X_{b:m} = x_2 \right) = G_j(x_2) - G_j(x_1)\), где \(G_j\) обозначает функцию распределения \(j\)-й порядковой статистики в выборке объёма \(n\) из распределения с функцией распределения \(G\). Поскольку \(G_j(x) = I_{G(x)}(j, n - j + 1)\), где \(I_x(y, z)\) обозначает неполный бета-интеграл (то есть функцию распределения бета-распределения), вероятность сигнала можно переписать как \(I_{G(x_2)}(j, n - j + 1) - I_{G(x_1)}(j, n - j + 1)\).

Следовательно, используя метод CUC, безусловную вероятность отсутствия сигнала можно найти, усредняя условную вероятность по совместному распределению \(X_{a:m}\) и \(X_{b:m}\). Хотя здесь это подробно не обсуждается, заинтересованному читателю рекомендуется обратиться к Chakraborti (2000) за подробным обсуждением метода CUC.

Эффективность контрольных карт оценивается через их распределение длины серии, которое можно рассматривать как время ожидания до первого сигнала. Времена ожидания определяются следующим образом:
\begin{align*}
	T_1 &= \min\left( t : Z_t = 1 \right), \\
	T_2^{(1)} &= \min\left( t : Z_{t-1} = 1, Z_t = 1 \right), \\
	T_2^{(2)} &= \min\left( t : Z_{t-1} = 1, Z_t = 2 \right), \\
	T_2 &= \min\left( T_2^{(1)}, T_2^{(2)} \right),
\end{align*}
где индикаторные переменные \(Z_h\) равны \(Z_h = 0\), если \(Y_{j:n} \in (X_{a:m}, X_{b:m})\), \(Z_h = 1\), если \(Y_{j:n} \ge X_{b:m}\), и \(Z_h = 2\), если \(Y_{j:n} \le X_{a:m}\).
Связывая это с правилами сигналов на основе серий, \(T_1\) --- это длина серии карты с правилом «1 из 1», тогда как \(T_2\) представляет длину серии карты с правилом «2 из 2». Используя метод MC, условное распределение \(T_2^*\), то есть времени ожидания до первого сигнала (или длины серии) для двусторонней карты KL «2 из 2», можно получить, применяя теорему 5.2 из Fu and Lou (2003, с.~68):
\[
	P\left( T_2^* = x \mid X_{a:m}, X_{b:m} \right) = \xi Q^{x-1}(I - Q)\mathbf{1}, \quad x \ge 2,
\]
где \(\xi = (1\ \ 0\ \ 0\ \ 0)\), а существенная матрица переходных вероятностей (TPM) \(Q\) приведена в таблице \ref{tab:4.29} (заинтересованному читателю также рекомендуется обратиться к Chakraborti, Human and Graham, 2009, с.~1058).

\begin{table}[!ht]
	\centering
	\caption{Существенная TPM (\(Q\)) для двусторонней контрольной карты, дополненной правилом сигналов на основе серий «2 из 2».}
	\label{tab:4.29}
	\begin{tabular}{@{}lcccc@{}}
		\toprule
		Состояния & \(S_\varphi\) & \(S_0\) & \(S_U\) & \(S_L\) \\
		\midrule
		\(S_\varphi\) & 0 & \(1 - p_L - p_U\) & \(p_U\) & \(p_L\) \\
		\(S_0\) & 0 & \(1 - p_L - p_U\) & \(p_U\) & \(p_L\) \\
		\(S_U\) & 0 & \(1 - p_L - p_U\) & 0 & \(p_L\) \\
		\(S_L\) & 0 & \(1 - p_L - p_U\) & \(p_U\) & 0 \\
		\bottomrule
	\end{tabular}
\end{table}

Отметим, что состояния перехода в таблице \ref{tab:4.29} обозначены \(S_\varphi\), \(S_0\), \(S_U\) и \(S_L\), где по определению George and Bowman (1995) (уравнение на с.~83) \(S_n\) обозначает общее число успехов в последовательности из \(n\) бинарных испытаний. Состояния перехода, упомянутые в таблице \ref{tab:4.29}, определяются как пустое состояние, обозначаемое \(S_\varphi\), состояние, обозначаемое \(S_0\) и определяемое как состояние, в которое переходят, если статистика для нанесения на карту попадает между UCL и LCL, состояние, обозначаемое \(S_U\) и определяемое как состояние, в которое переходят, если статистика для нанесения на карту попадает на UCL или выше, и, наконец, состояние, обозначаемое \(S_L\) и определяемое как состояние, в которое переходят, если статистика для нанесения на карту попадает на LCL или ниже.

Применяя метод CUC, безусловное распределение длины серии (времени ожидания) \(T_2\) получается из распределения \(T_2^*\) (путём усреднения по совместному распределению \(X_{a:m}\) и \(X_{b:m}\)):
\[
	P(T_2 = x) = \int_{-\infty}^{\infty}\int_{-\infty}^{\infty} P\left( T_2^* = x \mid X_{a:m} = x_1, X_{b:m} = x_2 \right) h_{a,b}(x_1, x_2)\, dx_1 dx_2, \quad x > 2,
\]
где \(P\left( T_2^* = x \mid X_{a:m}, X_{b:m} \right) = \xi Q^{x-1}(I - Q)\mathbf{1}\), \(x \ge 2\), а \(h_{a,b}(x_1, x_2)\) --- совместная плотность \(X_{a:m}\) и \(X_{b:m}\).

Chakraborti, Human and Graham (2009) вывели выражения для точных ARL как для карты DR, так и для карты KL. Идея, опять же, основана на CUC. При заданных \(X_{a:m}\) и \(X_{b:m}\) случайная величина \(T_2\) имеет так называемое геометрическое распределение порядка 2; см. Balakrishnan and Koutras (2002). Апеллируя к свойствам этого распределения, можно показать, что условная средняя длина серии (математическое ожидание) \(E\left( T_2 \mid X_{a:m} = x_1, X_{b:m} = x_2 \right)\) равна \(\frac{1 + p}{p^2}\), где
\begin{align*}
	p &= 1 - \left( G_j(x_2) - G_j(x_1) \right) \\
	  &= 1 - \left( I_{G(x_2)}(j, n - j + 1) - I_{G(x_1)}(j, n - j + 1) \right).
\end{align*}
Используя метод CUC, безусловная ARL для карты DR получается взятием математического ожидания условного математического ожидания и равна
\begin{align*}
	ARL_{DR} &= E_{X_{a:m}, X_{b:m}} E\left( T_2 \mid X_{a:m}, X_{b:m} \right) \\
	         &= \int_{0}^{1}\int_{0}^{x_2} \frac{2 - \left( G_j\left( F^{-1}(x_2) \right) - G_j\left( F^{-1}(x_1) \right) \right)}{\left( 1 - \left( G_j\left( F^{-1}(x_2) \right) - G_j\left( F^{-1}(x_1) \right) \right) \right)^2} \\
	         &\qquad \times f_{a,b}(x_1, x_2)\, dx_1 dx_2.
\end{align*}
Чтобы получить точную ARL карты KL «2 из 2», можно использовать вывод ARL из Klein (2000) вместе с методом CUC. Таким образом, при заданных \(X_{a:m}\) и \(X_{b:m}\) условное математическое ожидание \(T_2\) равно
\[
	E\left( T_2 \mid X_{a:m}, X_{b:m} \right) = \frac{1}{p - \dfrac{p_U}{1 + p_U} - \dfrac{p_L}{1 + p_L}},
\]
где \(p_U\) и \(p_L\) определены выше. Используя этот результат, безусловная ARL равна
\[
	\begin{aligned}
		ARL_{KL} = \int_{0}^{1}\int_{0}^{x_2}
		\frac{\left( 1 + \left( 1 - G_j\left( F^{-1}(x_2) \right) \right) \right)
			\left( 1 + G_j\left( F^{-1}(x_1) \right) \right)}{%
			\left( G_j\left( F^{-1}(x_1) \right) \right)^2
			\left( 1 + \left( 1 - G_j\left( F^{-1}(x_2) \right) \right) \right)} \\
		\frac{1}{\left( 1 - G_j\left( F^{-1}(x_2) \right) \right)^2
			\left( 1 + G_j\left( F^{-1}(x_1) \right) \right)}
		f_{a,b}(x_1, x_2)\, dx_1 dx_2.
	\end{aligned}
\]


\section{Приложения к главе 4}
\subsection{Контрольные карты Шухарта}
\subsubsection{Контрольная карта Shewhart-Prec}

\paragraph{Результат 4.1.} Безусловная вероятность превышения \(P\left( Y > X_{(r)} \right)\) равна \(\frac{m - r + 1}{m + 1}\), когда процесс находится в состоянии IC.

\noindent\textbf{Доказательство.} Отметим, что \(P\left( Y > X_{(r)} \right) = E\left( P\left( Y > X_{(r)} \mid X_{(r)} \right) \right) = E(p_r)\). Когда процесс находится в состоянии IC, \(F = G\) и тогда \(p_r = 1 - F\left( X_{(r)} \right) = 1 - V_r\) (скажем). Поскольку \(X_{(r)}\) является \(r\)-й порядковой статистикой в случайной выборке объёма \(m\) из \(F\), с использованием PIT можно показать, что \(V_r\) имеет распределение \(r\)-й порядковой статистики в случайной выборке объёма \(m\) из \(U(0, 1)\)-распределения, если \(F\) непрерывна. Это последнее распределение, как известно, является бета-распределением с параметрами \(r\) и \(m - r + 1\) соответственно. Таким образом, когда процесс находится в состоянии IC, \(p_r = 1 - V_r\) имеет бета-распределение с параметрами \(m - r + 1\) и \(r\). Значит,
\[
	E(p_r) = P\left( Y > X_{(r)} \right) = 1 - \frac{r}{m + 1} = \frac{m - r + 1}{m + 1} = d \ (\text{скажем}),
\]
когда процесс находится в состоянии IC, при использовании формулы математического ожидания бета-распределения. \(\square\)

\paragraph{Результат 4.2.} Распределение длины серии карты Shewhart-Prec и его свойства при некоторых конкретных альтернативах.

\noindent\textbf{Доказательство.} Можно видеть, что распределение длины серии и связанные с ним свойства зависят в общем случае от функции распределения \(F\) эталонного распределения и функции распределения \(G\) распределения Фазы II через композиционную функцию \(\Psi = GF^{-1}\). Таким образом, чтобы вычислить и/или проанализировать распределение длины серии и любые меры эффективности, нужно задать функцию \(\Psi\). Например, когда процесс находится в состоянии IC, то есть когда \(G = F\), имеем \(\Psi(t) = t\), и распределение длины серии в состоянии IC можно получить; оно приведено в уравнении \eqref{eq:4.26}. Однако в случае \(G \ne F\), что означает, что процесс тем или иным образом находится в состоянии OOC, следующие результаты полезны для нахождения функции \(\Psi\):

\begin{enumerate}
	\item Для альтернатив сдвига положения, когда процесс следует одному и тому же параметрическому распределению в Фазе I и Фазе II с изменением лишь параметра положения, пусть \(F(x) = H(x - \theta_1)\) и \(G(x) = H(x - \theta_2)\), где \(H\) --- непрерывная функция распределения, \(x \in \mathbb{R}\) и \(\theta_1, \theta_2 \in \mathbb{R}\), так что \(\Psi(t) = H\left( \theta_1 - \theta_2 + H^{-1}(t) \right)\). Отметим, что в случае IC \(\Psi(t) = t\).
	\item Для альтернатив масштаба, то есть когда процесс следует тому же параметрическому распределению в Фазе II, что и в Фазе I, но в этом случае возможен лишь сдвиг разброса или вариации процесса, пусть \(F(x) = H\left( \frac{x}{\gamma_1} \right)\) и \(G(x) = H\left( \frac{x}{\gamma_2} \right)\), где \(H\) --- непрерывная функция распределения, \(x \in \mathbb{R}\) и \(\gamma_1, \gamma_2 \in \mathbb{R}^+\), так что \(\Psi(t) = H\left( \frac{\gamma_1}{\gamma_2}H^{-1}(t) \right)\). Отметим, что в случае IC \(\Psi(t) = t\).
	\item Для альтернатив сдвига положения и масштаба пусть \(F(x) = H\left( \frac{x - \theta_1}{\gamma_1} \right)\) и \(G(x) = H\left( \frac{x - \theta_2}{\gamma_2} \right)\), где \(H\) --- непрерывная функция распределения, \(x \in \mathbb{R}\), \(\theta_1, \theta_2 \in \mathbb{R}\) и \(\gamma_1, \gamma_2 \in \mathbb{R}^+\), так что \(\Psi(t) = H\left( \frac{\theta_1 - \theta_2}{\gamma_2} + \frac{\gamma_1}{\gamma_2}H^{-1}(t) \right)\). Отметим, что в случае IC, опять же, \(\Psi(t) = t\).
\end{enumerate}

Таким образом, ясно, что распределение длины серии можно вычислить и/или проанализировать только тогда, когда известно изменение распределения процесса при переходе от IC к OOC или, в общем случае, когда нам известно соотношение между \(G\) и \(F\) через композиционную функцию \(\Psi\), а также величина сдвига параметров положения в единицах параметра масштаба распределения \(G\) и величина сдвига параметров масштаба. \(\square\)

\subsection{Контрольные карты CUSUM}
\subsubsection{Контрольная карта CUSUM-EX}

\paragraph{Результат 4.3.} При заданной (условно на) \(X_{(r)}\), то есть \(r\)-й порядковой статистике эталонной выборки, величины \(U_{j,r}\) независимы и имеют биномиальное распределение с параметрами \((n, p_r)\) для любого \(j = 1, 2, 3, \ldots\)

\noindent\textbf{Доказательство.} Поскольку \(U_{j,r}\) --- число наблюдений \(Y\) в \(j\)-й выборке Фазы II, которые превышают \(X_{(r)}\), при заданной \(X_{(r)}\) случайная величина \(U_{j,r}\) имеет биномиальное распределение с параметрами \((n, p_r)\), где \(p_r = P\left[ Y > X_{(r)} \mid X_{(r)} \right]\). Независимость \(U_{j,r}\) следует из независимости тестовых выборок. \(\square\)

\paragraph{Результат 4.4.} (Безусловное) распределение \(U_{j,r}\) в состоянии IC свободно от распределения и задаётся функцией вероятности
\[
	P(U_{j,r} = u) = \frac{\binom{u + m - r}{u}\binom{n - u + r - 1}{n - u}}{\binom{m + n}{n}}, \quad u = 0, 1, 2, \ldots, n.
\]

\noindent\textbf{Доказательство.} Используя результат 4.3 и усредняя по распределению \(X_{(r)}\), находим безусловное распределение \(U_{j,r}\). Таким образом,
\begin{align*}
	P(U_{j,r} = u) &= E\left[ P\left( U_{j,r} = u \mid X_{(r)} \right) \right] = \binom{n}{u}\int p_r^u(1 - p_r)^{n-u}\, dF\left( X_{(r)} \right) \\
	               &= \binom{n}{u}\int \left[ 1 - G(t) \right]^u\left[ G(t) \right]^{n-u} dF_{X_{(r)}}(t).
\end{align*}
Когда процесс находится в состоянии IC, имеем \(F = G\), и можно показать, что приведённый выше интеграл упрощается с помощью бета-функции до приведённого результата. Следовательно, статистики \(U_{j,r}\) безусловно свободны от распределения, когда процесс находится в состоянии IC. Тот же результат можно получить комбинаторными рассуждениями, детали которых здесь опущены. Отметим, что приведённая выше функция вероятности \(U_{j,r}\) известна как отрицательное гипергеометрическое распределение.

Далее расширим результат 4.4 и покажем, что (безусловное) совместное распределение \(U_{j,r}\) для \(j = 1, 2, \ldots, v\) свободно от распределения, когда процесс находится в состоянии IC. Это устанавливает, что предлагаемая карта свободна от распределения. \(\square\)

\paragraph{Результат 4.5.} Безусловное распределение \(U_{j,r}\) в состоянии IC для \(j = 1, 2, \ldots, v\), где \(v > 1\) --- положительное целое число, свободно от распределения.

\noindent\textbf{Доказательство.} Отметив, что по результату 4.3 при заданной \(X_{(r)}\) величины \(U_{j,r}\), \(j = 1, 2, \ldots, v\), независимы и имеют биномиальное распределение \((n, p_r)\), находим совместное распределение \((U_{1,r}, U_{2,r}, \ldots, U_{v,r})\), когда процесс находится в состоянии IC:
\begin{align*}
	&P\left[ U_{1,r} = u_1, U_{2,r} = u_2, \ldots, U_{v,r} = u_v \right] \\
	&\quad = E\left[ P\left[ U_{1,r} = u_1, U_{2,r} = u_2, \ldots, U_{v,r} = u_v \right] \mid X_{(r)} \right] \\
	&\quad = \binom{n}{u_1}\binom{n}{u_2}\cdots\binom{n}{u_v}\int \left( \frac{p_r}{1 - p_r} \right)^{\sum_{j=1}^{v} u_j}(1 - p_r)^{nv}\, dF\left( X_{(r)} \right) \\
	&\quad = \binom{n}{u_1}\binom{n}{u_2}\cdots\binom{n}{u_v}\frac{B\left( \sum_{j=1}^{v} u_j + m - r + 1,\; nv + r - \sum_{j=1}^{v} u_j \right)}{B(r, m - r + 1)}.
\end{align*}
Последний результат снова следует из того, что \(F = G\), когда процесс находится в состоянии IC, и упрощения интеграла с помощью бета-функции. Следовательно, безусловное совместное распределение в состоянии IC статистик превышения для любого числа независимых выборок Фазы II свободно от распределения. Это доказывает, что предлагаемая карта свободна от распределения. \(\square\)

\subsubsection{Контрольная карта CUSUM-rank}

\paragraph{Результат 4.6.} Условное распределение длины серии для карты CUSUM-rank.

Рассмотрим верхнюю одностороннюю карту CUSUM-rank, предложенную Li, Tang и Ng (2010). Из уравнения \eqref{eq:4.56} имеем, что статистика для нанесения на карту для этой карты задаётся как
\[
	C_j^+ = \max\left[ 0, C_{j-1}^+ + W_j - \frac{n(m + n + 1)}{2} - k \right]
\]
для \(j = 1, 2, 3, \ldots\). Предположим, что карта начинается со значения \(u\), то есть \(C_0^+ = u\), где \(0 \le u \le h\). Для длины серии (\(N\)), равной единице, имеем
\begin{align*}
	P_C\left( N = 1 \mid u, X \right)
	&= P_C\left( C_1^+ \ge h \mid u, X \right) \\
	&= P_C\left( u + W_1 - \frac{n(m + n + 1)}{2} - k \ge h \mid u, X \right) \\
	&= P_C\left( W_1 \ge \frac{n(m + n + 1)}{2} + k + h - u \mid u, X \right) \\
	&= P_C\left( W \ge \frac{n(m + n + 1)}{2} + k + h - u \mid u, X \right).
\end{align*}
Для \(t > 1\) имеем
\begin{align*}
	P_C\left( N = t \mid u, X \right)
	&= P_C\left( N = t, C_1^+ = 0 \mid u, X \right) + P_C\left( N = t, 0 < C_1^+ < h \mid u, X \right) \\
	&= P_C\left( N = t \mid C_1^+ = 0, u, X \right)\cdot P_C\left( C_1^+ = 0 \mid u, X \right) \\
	&\quad + P_C\left( N = t \mid 0 < C_1^+ < h, u, X \right)\cdot P_C\left( 0 < C_1^+ < h \mid u, X \right)
\end{align*}
по свойствам маргинальных, совместных и условных вероятностей.

Теперь, когда \(C_1^+ = 0\), имеем \(P_C\left( N = t \mid C_1^+ = 0, u, X \right) = P_C\left( N = t - 1 \mid u = 0, X \right)\).

Зная, что
\[
	C_1^+ = \max\left[ 0, u + W_1 - \frac{n(m + n + 1)}{2} - k \right]
\]
и что \(C_1^+ = 0\), можно заключить, что
\begin{align*}
	P_C\left( C_1^+ = 0 \mid u, X \right)
	&= P_C\left( u + W_1 - \frac{n(m + n + 1)}{2} - k \le 0 \mid u, X \right) \\
	&= P_C\left( W_1 \le k + \frac{n(m + n + 1)}{2} - u \mid u, X \right).
\end{align*}
Следовательно,
\begin{align*}
	&P_C\left( N = t \mid C_1^+ = 0, u, X \right)\cdot P_C\left( C_1^+ = 0 \mid u, X \right) \\
	&\qquad + P_C\left( N = t \mid 0 < C_1^+ < h, u, X \right)\cdot P_C\left( 0 < C_1^+ < h \mid u, X \right) \\
	&\quad = P_C\left( N = t - 1 \mid u = 0, X \right)\cdot P_C\left( W_1 \le k + \frac{n(m + n + 1)}{2} - u \mid u, X \right) \\
	&\qquad + \sum_{v \in V}\left( P_C\left( N = t \mid C_1^+ = v, C_0^+ = u, X \right)\cdot P_C\left( C_1^+ = v \mid C_0^+ = u, X \right) \right) \\
	&\quad = P_C\left( N = t - 1 \mid u = 0, X \right)\cdot P_C\left( W \le k + \frac{n(m + n + 1)}{2} - u \mid u, X \right) \\
	&\qquad + \sum_{v \in V} P_C\left( N = t - 1 \mid v, X \right)\cdot P_C\left( W = v + \frac{n(m + n + 1)}{2} + k - u \mid u, X \right),
\end{align*}
где множество \(V\) содержит все значения, которые \(C_1^+\) может принимать строго между 0 и \(h\), исключая 0 и \(h\). После этого условное распределение длины серии можно вычислить рекурсивно. Условное распределение длины серии для нижней односторонней карты CUSUM-rank можно получить аналогично; оно приведено ниже. Опять же, предположим, что \(u\) означает, что карта начинается со значения \(u\), то есть \(C_0^- = u\), где \(0 \le u \le h\). Имеем \(C_1^- = \max\left[ 0, C_0^- - \left( W_1 - \frac{n(m + n + 1)}{2} \right) - k \right]\), так что
\begin{align*}
	P_C\left( N = 1 \mid u, X \right)
	&= P_C\left( C_1^- \ge h \mid u, X \right) \\
	&= P_C\left( u - W_1 + \frac{n(m + n + 1)}{2} - k \ge h \mid u, X \right) \\
	&= P_C\left( W_1 \ge k - u + \frac{n(m + n + 1)}{2} - h \mid u, X \right) \\
	&= P_C\left( W \ge k - u + \frac{n(m + n + 1)}{2} - h \mid u, X \right).
\end{align*}
Хотя не все шаги для \(t > 1\) показаны для нижней односторонней карты CUSUM-rank, выводы аналогичны выводам для верхней односторонней карты CUSUM-rank, и ниже приведено лишь итоговое выражение. Для \(t > 1\) имеем
\begin{align*}
	P_C\left( N = t \mid u, X \right)
	&= P_C\left( N = t - 1 \mid u = 0, X \right)\cdot P_C\left( W \ge u + \frac{n(m + n + 1)}{2} - k \mid u, X \right) \\
	&\quad + \sum_{v \in V} P_C\left( N = t - 1 \mid v, X \right)\cdot P_C\left( W = u + \frac{n(m + n + 1)}{2} - k - v \mid u, X \right),
\end{align*}
где множество \(V\) содержит все значения, которые \(C_1^-\) может принимать строго между 0 и \(h\), исключая 0 и \(h\).

\paragraph{Результат 4.7.} Безусловное распределение длины серии для карты CUSUM-rank.

В работе Li, Tang и Ng (2010) безусловное распределение длины серии можно найти из
\[
	P_U(N = t) = \int_{-\infty}^{\infty}\cdots\int_{-\infty}^{\infty} P_C\left( N = t \mid u, X \right) dF(x_1)\cdots dF(x_n).
\]
Поскольку это \(n\)-мерный интеграл, а \(P_C\left( N = t \mid u, X \right)\) в общем случае не имеет замкнутой формы, прямое численное решение приведённого выше \(n\)-мерного интеграла получить трудно. В качестве альтернативы можно приблизить \(n\)-мерный интеграл, используя
\[
	P_U(N = t) \approx \frac{1}{Sim}\sum_{i=1}^{Sim} P_C\left( N = t \mid u, X_j \right),
\]
где \(Sim\) --- число симуляций Монте-Карло, а \(X_j\) --- \(j\)-я случайно сгенерированная эталонная выборка.

\subsection{Контрольные карты EWMA}
\subsubsection{Контрольная карта EWMA-SN}

\paragraph{Результат 4.8.} Вывод CL и контрольных границ карты EWMA-SN.

По определению статистики для нанесения на карту (см. уравнение \eqref{eq:4.13}) и используя рекурсивную подстановку, находим, что
\begin{align*}
	E(Z_i \mid IC) &= E\left( \lambda SN_i + (1 - \lambda)Z_{i-1} \mid IC \right) \\
	               &= E\left( \lambda\sum_{j=0}^{i-1}(1 - \lambda)^j SN_{i-j} + (1 - \lambda)^i Z_0 \mid IC \right).
\end{align*}
Теперь, используя тот факт, что \(E(SN_{i-j} \mid IC) = 0\) и \(Z_0 = 0\), следует, что
\[
	E(Z_i \mid IC) = \lambda\sum_{j=0}^{i-1}(1 - \lambda)^j E(SN_{i-j} \mid IC) + (1 - \lambda)^i Z_0 = 0.
\]
Отсюда получаем результат. Для вычисления дисперсии EWMA-статистики для нанесения на карту выполняются аналогичные шаги, то есть мы снова используем определение статистики для нанесения на карту (см. уравнение \eqref{eq:4.13}) и рекурсивную подстановку. Однако мы также используем результат для суммы конечной геометрической прогрессии, свойства дисперсии независимых случайных величин и тот факт, что \(\operatorname{var}(SN_{i-j} \mid IC) = n\) при \(p = 0.5\). Следовательно,
\begin{align*}
	\operatorname{var}(Z_i \mid IC) &= \operatorname{var}\left( \lambda\sum_{j=0}^{i-1}(1 - \lambda)^j SN_{i-j} + (1 - \lambda)^i Z_0 \mid IC \right) \\
	&= \lambda^2\sum_{j=0}^{i-1}(1 - \lambda)^{2j}\operatorname{var}(SN_{i-j} \mid IC) \\
	&= n\lambda^2\sum_{j=0}^{i-1}(1 - \lambda)^{2j} \\
	&= n\lambda^2\frac{1 - (1 - \lambda)^{2i}}{1 - (1 - \lambda)^2} \\
	&= n\lambda\frac{1 - (1 - \lambda)^{2i}}{2 - \lambda}.
\end{align*}
Следовательно, \(SD(Z_i \mid IC) = \sigma_{Z_i \mid IC} = \sqrt{\frac{\lambda n}{2 - \lambda}\left( 1 - (1 - \lambda)^{2i} \right)}\).

Наконец, поскольку \(LCL/UCL = E(Z_i \mid IC) \pm L\sigma_{Z_i \mid IC}\) и \(CL = E(Z_i \mid IC)\), получаем контрольные границы и CL, приведённые в уравнении \eqref{eq:4.14}.

\subsubsection{Контрольная карта EWMA-SR}

\paragraph{Результат 4.9.} Вывод CL и контрольных границ карты EWMA-SR.

По определению статистики для нанесения на карту (см. уравнение \eqref{eq:4.20}) и используя рекурсивную подстановку, получаем следующий результат:
\begin{align*}
	E(Z_i \mid IC) &= E\left( \lambda SR_i + (1 - \lambda)Z_{i-1} \mid IC \right) \\
	               &= E\left( \lambda\sum_{j=0}^{i-1}(1 - \lambda)^j SR_{i-j} + (1 - \lambda)^i Z_0 \mid IC \right).
\end{align*}
Используя тот факт, что \(E(SR_{i-j} \mid IC) = 0\) и \(Z_0 = 0\), имеем
\[
	E(Z_i \mid IC) = \lambda\sum_{j=0}^{i-1}(1 - \lambda)^j E(SR_{i-j} \mid IC) + (1 - \lambda)^i Z_0 = 0.
\]
Для получения дисперсии выполняются аналогичные шаги, то есть мы снова используем определение статистики для нанесения на карту (см. уравнение \eqref{eq:4.20}) и рекурсивную подстановку. Однако мы также используем результат для суммы конечной геометрической прогрессии и тот факт, что \(\operatorname{var}(SR_i \mid IC) = \frac{n(n + 1)(2n + 1)}{6}\). Следовательно, имеем
\begin{align*}
	\operatorname{var}(Z_i \mid IC) &= \operatorname{var}\left( \lambda\sum_{j=0}^{i-1}(1 - \lambda)^j SR_{i-j} + (1 - \lambda)^i Z_0 \mid IC \right) \\
	&= \lambda^2\sum_{j=0}^{i-1}(1 - \lambda)^{2j}\operatorname{var}(SR_{i-j} \mid IC) \\
	&= \frac{n(n + 1)(2n + 1)}{6}\lambda^2\sum_{j=0}^{i-1}(1 - \lambda)^{2j} \\
	&= \frac{n(n + 1)(2n + 1)}{6}\lambda^2\frac{1 - (1 - \lambda)^{2i}}{1 - (1 - \lambda)^2}.
\end{align*}
Следовательно,
\[
	SD(Z_i \mid IC) = \sigma_{Z_i \mid IC} = \sqrt{\frac{n(n + 1)(2n + 1)}{6}\cdot\frac{\lambda}{2 - \lambda}\left( 1 - (1 - \lambda)^{2i} \right)}.
\]
Наконец, поскольку \(LCL/UCL = E(Z_i \mid IC) \pm L\sigma_{Z_i \mid IC}\) и \(CL = E(Z_i \mid IC)\), получаем контрольные границы и CL, приведённые в уравнении \eqref{eq:4.21}.

\subsubsection{Контрольная карта EWMA-EX}

\paragraph{Результат 4.10.} Точные меняющиеся во времени контрольные границы для карты EWMA-EX.

Для точных меняющихся во времени контрольных границ имеем
\[
	E(Z_j \mid IC) = n(1 - a)\left( 1 - (1 - \lambda)^j \right)
\]
и
\[
	SD(Z_j \mid IC) = \sqrt{\frac{na(1 - a)}{m + 2}n\left( 1 - (1 - \lambda)^j \right)^2 + \frac{\lambda(m + 1)}{2 - \lambda}\left( 1 - (1 - \lambda)^{2j} \right)}
\]
при \(LCL/UCL = E(Z_j \mid IC) \pm L \times SD(Z_j \mid IC)\) и \(CL = E(Z_j \mid IC)\).

\noindent\textbf{Доказательство.} Условное математическое ожидание статистики для нанесения на карту \(Z_j = \lambda U_{j,r} + (1 - \lambda)Z_{j-1}\) в состоянии IC. Используя свойства математических ожиданий, рекурсивную подстановку и формулу суммы конечной геометрической прогрессии вместе с тем фактом, что \(E\left( U_{j,r} \mid X_{(r)} \right) = np_r\) для всех \(j\) (\(j = 1, 2, \ldots\)) и \(Z_0 = 0\), получаем следующее:
\begin{align*}
	E\left( Z_j \mid X_{(r)} \right) &= E\left( \lambda U_{j,r} + (1 - \lambda)Z_{j-1} \mid X_{(r)} \right) = E\left( \lambda\sum_{k=0}^{j-1}(1 - \lambda)^k U_{j-k,r} + (1 - \lambda)^j Z_0 \mid X_{(r)} \right) \\
	&= np_r\left( 1 - (1 - \lambda)^j \right).
\end{align*}
Условная дисперсия и стандартное отклонение статистики для нанесения на карту \(Z_j = \lambda U_{j,r} + (1 - \lambda)Z_{j-1}\) в состоянии IC. Используя свойства дисперсии, рекурсивную подстановку и формулу суммы конечной геометрической прогрессии вместе с тем фактом, что \(VAR\left( U_{j,r} \mid X_{(r)} \right) = np_r(1 - p_r)\) для всех \(j\) (\(j = 1, 2, \ldots\)), получаем
\begin{align*}
	\operatorname{var}\left( Z_j \mid X_{(r)} \right) &= \operatorname{var}\left( \lambda\sum_{k=0}^{j-1}(1 - \lambda)^k U_{j-k,r} + (1 - \lambda)^j Z_0 \mid X_{(r)} \right) \\
	&= np_r(1 - p_r)\lambda^2\frac{1 - (1 - \lambda)^{2j}}{1 - (1 - \lambda)^2}.
\end{align*}
Последнее выражение упрощается к \(np_r(1 - p_r)\lambda\frac{1 - (1 - \lambda)^{2j}}{2 - \lambda}\), так что
\[
	SD\left( Z_j \mid X_{(r)} \right) = \sqrt{\frac{\lambda}{2 - \lambda}\left( 1 - (1 - \lambda)^{2j} \right)np_r(1 - p_r)}.
\]
Чтобы получить безусловные свойства статистики для нанесения на карту, сначала нужно привести некоторые свойства \(p_r\).

\emph{Свойства \(p_r\).} В случае IC \(F\left( X_{(r)} \right) = 1 - p_r\) имеет бета-распределение с параметрами \(r\) и \(m + 1 - r\). Следовательно,
\[
	E(p_r) = \frac{m + 1 - r}{m + 1} \quad \text{и} \quad E\left( p_r^2 \right) = \frac{m^2 + 3m + 2 - 2rm - 3r + r^2}{(m + 1)(m + 2)},
\]
так что \(\operatorname{var}(p_r) = \frac{r(m + 1 - r)}{(m + 1)^2(m + 2)}\).

Безусловное математическое ожидание статистики для нанесения на карту \(Z_j = \lambda U_{j,r} + (1 - \lambda)Z_{j-1}\) в состоянии IC:
\begin{align*}
	E(Z_j) &= E_{X_{(r)}} E\left( Z_j \mid X_{(r)} \right) \\
	       &= E_{X_{(r)}}\left( np_r\left( 1 - (1 - \lambda)^j \right) \right) = n\left( 1 - (1 - \lambda)^j \right)E_{X_{(r)}}(p_r) \\
	       &= n\left( 1 - (1 - \lambda)^j \right)\frac{m + 1 - r}{m + 1} = n\left( 1 - (1 - \lambda)^j \right)(1 - a),
\end{align*}
где \(a = r/(m + 1)\).

Безусловная дисперсия и стандартное отклонение статистики для нанесения на карту \(Z_j = \lambda U_{j,r} + (1 - \lambda)Z_{j-1}\) в состоянии IC:
\begin{align*}
	\operatorname{var}(Z_j) &= \operatorname{var}\left[ E\left( Z_j \mid X_{(r)} \right) \right] + E\left[ \operatorname{var}\left( Z_j \mid X_{(r)} \right) \right] \\
	&= \operatorname{var}\left[ np_r\left( 1 - (1 - \lambda)^j \right) \right] + E\left[ \frac{\lambda}{2 - \lambda}\left( 1 - (1 - \lambda)^{2j} \right)np_r(1 - p_r) \right] \\
	&= n^2\left( 1 - (1 - \lambda)^j \right)^2\operatorname{var}(p_r) + \frac{\lambda}{2 - \lambda}\left( 1 - (1 - \lambda)^{2j} \right)n\left( E(p_r) - E\left( p_r^2 \right) \right) \\
	&= n^2\frac{r(m - r + 1)}{(m + 1)^2(m + 2)}\left( 1 - (1 - \lambda)^j \right)^2 + \frac{\lambda}{2 - \lambda}\left( 1 - (1 - \lambda)^{2j} \right)n\frac{r(m - r + 1)}{(m + 1)(m + 2)} \\
	&= n^2\frac{r(m - r + 1)}{(m + 1)(m + 1)(m + 2)}\left( 1 - (1 - \lambda)^j \right)^2 + \frac{\lambda}{2 - \lambda}\left( 1 - (1 - \lambda)^{2j} \right)n\frac{r(m - r + 1)}{(m + 1)(m + 2)} \\
	&= n\left[ \frac{a(1 - a)}{m + 2}n\left( 1 - (1 - \lambda)^j \right)^2 + \frac{\lambda(m + 1)}{2 - \lambda}\left( 1 - (1 - \lambda)^{2j} \right) \right],
\end{align*}
где \(a = r/(m + 1)\). Следовательно,
\[
	SD(Z_j \mid IC) = \sqrt{\frac{na(1 - a)}{m + 2}n\left( 1 - (1 - \lambda)^j \right)^2 + \frac{\lambda(m + 1)}{2 - \lambda}\left( 1 - (1 - \lambda)^{2j} \right)}. \ \square
\]

\paragraph{Результат 4.11.} Стационарные контрольные границы для карты EWMA-EX.

Для стационарных контрольных границ имеем
\[
	E(Z_j \mid IC) = n(1 - a)
\]
и
\[
	SD(Z_j \mid IC) = \sqrt{\frac{na(1 - a)}{m + 2}\left( n + \frac{\lambda(m + 1)}{2 - \lambda} \right)}
\]
при \(LCL/UCL = E(Z_j \mid IC) \pm L \times SD(Z_j \mid IC)\) и \(CL = E(Z_j \mid IC)\).

\noindent\textbf{Доказательство.} Условное математическое ожидание статистики для нанесения на карту \(Z_j = \lambda U_{j,r} + (1 - \lambda)Z_{j-1}\) в состоянии IC. Используя свойства математических ожиданий, рекурсивную подстановку и формулу суммы конечной геометрической прогрессии вместе с тем фактом, что \(E\left( U_{j,r} \mid X_{(r)} \right) = np_r\) для всех \(j\) (\(j = 1, 2, \ldots\)) и \(Z_0 = np_r\), получаем следующее:
\begin{align*}
	E\left( Z_j \mid X_{(r)} \right) &= E\left( \lambda U_{j,r} + (1 - \lambda)Z_{j-1} \mid X_{(r)} \right) \\
	&= E\left( \lambda\sum_{k=0}^{j-1}(1 - \lambda)^k U_{j-k,r} + (1 - \lambda)^j Z_0 \mid X_{(r)} \right) \\
	&= np_r\left( 1 - (1 - \lambda)^j \right) + np_r(1 - \lambda)^j = np_r.
\end{align*}
Условная дисперсия и стандартное отклонение статистики для нанесения на карту \(Z_j = \lambda U_{j,r} + (1 - \lambda)Z_{j-1}\) в состоянии IC. Используя свойства дисперсии, рекурсивную подстановку и формулу суммы конечной геометрической прогрессии вместе с тем фактом, что \(\operatorname{var}\left( U_{j,r} \mid X_{(r)} \right) = np_r(1 - p_r)\) для всех \(j\) (\(j = 1, 2, \ldots\)), получаем следующее:
\begin{align*}
	\operatorname{var}\left( Z_j \mid X_{(r)} \right) &= \operatorname{var}\left( \lambda\sum_{k=0}^{j-1}(1 - \lambda)^k U_{j-k,r} + (1 - \lambda)^j Z_0 \mid X_{(r)} \right) \\
	&= np_r(1 - p_r)\lambda^2\frac{1 - (1 - \lambda)^{2j}}{1 - (1 - \lambda)^2}.
\end{align*}
Последнее выражение упрощается к \(np_r(1 - p_r)\lambda\frac{1 - (1 - \lambda)^{2j}}{2 - \lambda}\), так что
\[
	SD\left( Z_j \mid X_{(r)} \right) = \sqrt{\frac{\lambda}{2 - \lambda}\left( 1 - (1 - \lambda)^{2j} \right)np_r(1 - p_r)}.
\]
Однако член \(\left( 1 - (1 - \lambda)^{2j} \right)\) стремится к единице при \(j \to \infty\), так что этим членом можно пренебречь:
\[
	SD\left( Z_j \mid X_{(r)} \right) = \sqrt{\frac{\lambda}{2 - \lambda}np_r(1 - p_r)}.
\]
Чтобы получить безусловные свойства статистики для нанесения на карту, сначала нужно привести некоторые свойства \(p_r\).

\emph{Свойства \(p_r\).} В случае IC \(F\left( X_{(r)} \right) = 1 - p_r\) имеет бета-распределение с параметрами \(r\) и \(m + 1 - r\). Следовательно,
\begin{align*}
	E(p_r) &= \frac{m + 1 - r}{m + 1} \quad \text{и} \quad E\left( p_r^2 \right) = \frac{m^2 + 3m + 2 - 2rm - 3r + r^2}{(m + 1)(m + 2)}, \\
	\text{так что} \quad \operatorname{var}(p_r) &= \frac{r(m + 1 - r)}{(m + 1)^2(m + 2)}.
\end{align*}
Безусловное математическое ожидание статистики для нанесения на карту \(Z_j = \lambda U_{j,r} + (1 - \lambda)Z_{j-1}\) в состоянии IC:
\begin{align*}
	E(Z_j) &= E_{X_{(r)}}\left( Z_j \mid X_{(r)} \right) = E_{X_{(r)}}(np_r) = nE_{X_{(r)}}(p_r) \\
	       &= n\frac{m + 1 - r}{m + 1} = n(1 - a),
\end{align*}
где \(a = r/(m + 1)\).

Безусловная дисперсия и стандартное отклонение статистики для нанесения на карту \(Z_j = \lambda U_{j,r} + (1 - \lambda)Z_{j-1}\) в состоянии IC:
\begin{align*}
	\operatorname{var}(Z_j) &= \operatorname{var}\left[ E\left( Z_j \mid X_{(r)} \right) \right] + E\left[ \operatorname{var}\left( Z_j \mid X_{(r)} \right) \right] \\
	&= \operatorname{var}(np_r) + E\left[ \frac{\lambda}{2 - \lambda}np_r(1 - p_r) \right] \\
	&= n^2\operatorname{var}(p_r) + \frac{\lambda}{2 - \lambda}n\left( E(p_r) - E\left( p_r^2 \right) \right) \\
	&= n^2\frac{r(m - r + 1)}{(m + 1)^2(m + 2)} + \frac{\lambda}{2 - \lambda}n\frac{r(m - r + 1)}{(m + 1)(m + 2)} \\
	&= n^2\frac{r(m - r + 1)}{(m + 1)(m + 1)(m + 2)} + \frac{\lambda}{2 - \lambda}n\frac{r(m - r + 1)}{(m + 1)(m + 2)} \\
	&= n\left[ \frac{a(1 - a)}{m + 2}\left( n + \frac{\lambda(m + 1)}{2 - \lambda} \right) \right],
\end{align*}
где \(a = r/(m + 1)\). Следовательно,
\[
	SD(Z_j \mid IC) = \sqrt{\frac{na(1 - a)}{m + 2}\left( n + \frac{\lambda(m + 1)}{2 - \lambda} \right)}. \ \square
\]

\subsubsection{Контрольная карта EWMA-rank}

\paragraph{Результат 4.12.} Условное распределение длины серии для карты EWMA-rank.

Статистика для нанесения на карту для карты EWMA-rank, предложенной Li, Tang и Ng (2010), задаётся как \(Z_j = \lambda W_j + (1 - \lambda)Z_{j-1}\) для \(j = 1, 2, 3, \ldots\) (см. уравнение \eqref{eq:4.61}). Предположим, что \(u\) обозначает начальное значение статистики для нанесения на карту, то есть \(Z_0 = u\).

Чтобы длина серии (\(N\)) была равна единице, должно выполняться
\begin{align*}
	P_C\left( N = 1 \mid u, X \right)
	&= P_C\left( Z_1 - \frac{n(m + n + 1)}{2} \le LCL \cup Z_1 - \frac{n(m + n + 1)}{2} \ge UCL \mid u, X \right) \\
	&= P_C\left( \lambda W_1 + (1 - \lambda)u - \frac{n(m + n + 1)}{2} \le LCL \right. \\
	&\qquad \left. \cup\ \lambda W_1 + (1 - \lambda)u - \frac{n(m + n + 1)}{2} \ge UCL \mid u, X \right) \\
	&= P_C\left( W_1 \le \frac{n(m + n + 1) + 2LCL - 2(1 - \lambda)u}{2\lambda} \mid u, X \right) \\
	&\qquad + P_C\left( W_1 \ge \frac{n(m + n + 1) + 2UCL - 2(1 - \lambda)u}{2\lambda} \mid u, X \right) \\
	&= P_C\left( W \le \frac{n(m + n + 1) + 2LCL - 2(1 - \lambda)u}{2\lambda} \mid u, X \right) \\
	&\qquad + P_C\left( W \ge \frac{n(m + n + 1) + 2UCL - 2(1 - \lambda)u}{2\lambda} \mid u, X \right).
\end{align*}
Используя рекурсивный подход, имеем для \(t = 1, 2, 3, \ldots\):
\begin{align*}
	P_C\left( N = t \mid u, X \right)
	&= P_C\left( N = t,\ \frac{n(m + n + 1)}{2} + LCL < Z_1 < \frac{n(m + n + 1)}{2} + UCL \mid u, X \right) \\
	&= P_C\left( N = t \mid \frac{n(m + n + 1)}{2} + LCL < Z_1 < \frac{n(m + n + 1)}{2} + UCL,\ u, X \right) \\
	&\quad \cdot P_C\left( \frac{n(m + n + 1)}{2} + LCL < Z_1 < \frac{n(m + n + 1)}{2} + UCL \mid u, X \right) \\
	&= \sum_{v \in V}\left( P_C\left( N = t \mid Z_1 = v, Z_0 = u, X \right) P_C\left( Z_1 = v \mid Z_0 = u, X \right) \right),
\end{align*}
где множество \(V\) содержит все возможные значения, которые \(Z_1\) может принимать между LCL и UCL. Поскольку известно, что \(Z_1\) принимает значение между LCL и UCL, мы можем воспользоваться тем фактом, что процесс находится в состоянии IC в момент 1, и получить
\begin{align*}
	&= \sum_{v \in V}\left( P_C\left( N = t - 1 \mid Z_0 = u, X \right) P_C\left( Z_1 = v \mid Z_0 = u, X \right) \right) \\
	&= \sum_{v \in V} P_C\left( N = t - 1 \mid v, X \right) P_C\left( W = \frac{v - (1 - \lambda)u}{\lambda} \mid u, X \right).
\end{align*}
